% Options for packages loaded elsewhere
\PassOptionsToPackage{unicode}{hyperref}
\PassOptionsToPackage{hyphens}{url}
\documentclass[
]{article}
\usepackage{xcolor}
\usepackage[margin=1in]{geometry}
\usepackage{amsmath,amssymb}
\setcounter{secnumdepth}{5}
\usepackage{iftex}
\ifPDFTeX
  \usepackage[T1]{fontenc}
  \usepackage[utf8]{inputenc}
  \usepackage{textcomp} % provide euro and other symbols
\else % if luatex or xetex
  \usepackage{unicode-math} % this also loads fontspec
  \defaultfontfeatures{Scale=MatchLowercase}
  \defaultfontfeatures[\rmfamily]{Ligatures=TeX,Scale=1}
\fi
\usepackage{lmodern}
\ifPDFTeX\else
  % xetex/luatex font selection
\fi
% Use upquote if available, for straight quotes in verbatim environments
\IfFileExists{upquote.sty}{\usepackage{upquote}}{}
\IfFileExists{microtype.sty}{% use microtype if available
  \usepackage[]{microtype}
  \UseMicrotypeSet[protrusion]{basicmath} % disable protrusion for tt fonts
}{}
\makeatletter
\@ifundefined{KOMAClassName}{% if non-KOMA class
  \IfFileExists{parskip.sty}{%
    \usepackage{parskip}
  }{% else
    \setlength{\parindent}{0pt}
    \setlength{\parskip}{6pt plus 2pt minus 1pt}}
}{% if KOMA class
  \KOMAoptions{parskip=half}}
\makeatother
\usepackage{longtable,booktabs,array}
\usepackage{calc} % for calculating minipage widths
% Correct order of tables after \paragraph or \subparagraph
\usepackage{etoolbox}
\makeatletter
\patchcmd\longtable{\par}{\if@noskipsec\mbox{}\fi\par}{}{}
\makeatother
% Allow footnotes in longtable head/foot
\IfFileExists{footnotehyper.sty}{\usepackage{footnotehyper}}{\usepackage{footnote}}
\makesavenoteenv{longtable}
\setlength{\emergencystretch}{3em} % prevent overfull lines
\providecommand{\tightlist}{%
  \setlength{\itemsep}{0pt}\setlength{\parskip}{0pt}}
\usepackage{amsmath}
\usepackage{amssymb}
\usepackage{graphicx}
\usepackage{bm}
\usepackage{bookmark}
\IfFileExists{xurl.sty}{\usepackage{xurl}}{} % add URL line breaks if available
\urlstyle{same}
\hypersetup{
  pdftitle={Math 2117: Final Exam Boss Notes},
  pdfauthor={Compiled for Exam Preparation},
  hidelinks,
  pdfcreator={LaTeX via pandoc}}

\title{Math 2117: Final Exam Boss Notes}
\author{Compiled for Exam Preparation}
\date{June 2026}


\usepackage{xcolor}
\definecolor{calloutbg}{HTML}{F7F7F7}
\definecolor{calloutframe}{HTML}{D0D0D0}
\hypersetup{
  colorlinks=true,
  linkcolor=blue,
  filecolor=blue,
  citecolor=blue,
  urlcolor=blue
}

\usepackage[most]{tcolorbox}
\renewenvironment{quote}
  {\begin{tcolorbox}[colback=calloutbg,colframe=calloutframe,arc=2pt,boxrule=0.5pt,left=3pt,right=3pt,top=3pt,bottom=3pt]}
  {\end{tcolorbox}}

\usepackage{enumitem}
\setlist{nosep}
\setlist[itemize]{leftmargin=15pt}
\setlist[enumerate]{leftmargin=15pt}

\setlength{\parskip}{3pt plus 1pt minus 1pt}

\begin{document}
\maketitle

{
\setcounter{tocdepth}{3}
\tableofcontents
}

\section{A01: Vector Basics and Direction
Cosines}\label{a01-vector-basics-and-direction-cosines}

\begin{quote}
\textbf{Exam Weight}: Tier 3 \textbar{} \textbf{Appeared in}: 2019,
2020, 2021, 2023 \textbar{} \textbf{Typical Marks}: 3--4
\end{quote}

This section covers the absolute basics of vectors: definitions, unit
vectors, and direction cosines. These concepts form the foundation for
all subsequent vector calculus operations.

\subsection{Definitions}\label{definitions}

\begin{itemize}
\tightlist
\item
  \textbf{Scalar Quantity}: A physical quantity that has only magnitude
  but no direction (e.g., Mass, Temperature, Time).
\item
  \textbf{Vector Quantity}: A physical quantity that has both magnitude
  and direction, and obeys vector addition rules (e.g., Velocity, Force,
  Acceleration).
\end{itemize}

A vector field
\[
\vec{V}(x, y) = P(x, y)\hat{i} + Q(x, y)\hat{j}
\]
assigns a vector to every point in space. - A \textbf{source} field
(e.g.,
\[
\vec{V} = x\hat{i} + y\hat{j}
\]
) has all vectors pointing radially outward. - A \textbf{sink} field
(e.g.,
\[
\vec{V} = -x\hat{i} - y\hat{j}
\]
) has all vectors pointing radially inward.

\subsection{Unit Vectors and
Resultants}\label{unit-vectors-and-resultants}

The \textbf{resultant} of two vectors \(\vec{r}_1\) and \(\vec{r}_2\) is
simply their vector sum:
\[
\vec{R} = \vec{r}_1 + \vec{r}_2.
\]
The \textbf{unit vector} parallel to any vector \(\vec{R}\) is the
vector divided by its magnitude:
\[
\hat{u} = \frac{\vec{R}}{\lvert \vec{R} \rvert}
\]
\subsubsection{Worked Example (PYQ 2020)}\label{worked-example-pyq-2020}

\textbf{Problem}: Find a unit vector parallel to the resultant of
\[
\vec{r}_1 = 2\hat{i} + 4\hat{j} - 5\hat{k}
\]
and
\[
\vec{r}_2 = \hat{i} + 2\hat{j} + \hat{k}.
\]
\textbf{Solution}:
\[
\vec{R} = \vec{r}_1 + \vec{r}_2 = 3\hat{i} + 6\hat{j} - 4\hat{k}
\]
\[
\lvert \vec{R} \rvert = \sqrt{3^2 + 6^2 + (-4)^2} = \sqrt{9 + 36 + 16} = \sqrt{61}
\]
Unit vector
\[
\hat{u} = \frac{3\hat{i} + 6\hat{j} - 4\hat{k}}{\sqrt{61}}
\]

\subsection{Direction Cosines}\label{direction-cosines}

Let a position vector be
\[
\vec{r} = x\hat{i} + y\hat{j} + z\hat{k}
\]
with magnitude
\[
r = \sqrt{x^2 + y^2 + z^2}.
\]
The angles \(\alpha, \beta, \gamma\) that the vector makes with the
positive \(x, y, z\) axes respectively are defined by their
\textbf{direction cosines}:
\[
\cos\alpha = \frac{x}{r}, \quad \cos\beta = \frac{y}{r}, \quad \cos\gamma = \frac{z}{r}
\]
A fundamental identity is that the sum of their squares is always 1:
\[
\cos^2\alpha + \cos^2\beta + \cos^2\gamma = \frac{x^2}{r^2} + \frac{y^2}{r^2} + \frac{z^2}{r^2} = \frac{x^2 + y^2 + z^2}{r^2} = 1
\]
\subsubsection{Worked Example (PYQ 2019)}\label{worked-example-pyq-2019}

\textbf{Problem}: Find the acute angles which the line joining
\(P(1, -3, 2)\) and \(Q(3, -5, 1)\) makes with the coordinate axes.
\textbf{Solution}:
\[
\vec{PQ} = (3 - 1)\hat{i} + (-5 - (-3))\hat{j} + (1 - 2)\hat{k} = 2\hat{i} - 2\hat{j} - \hat{k}
\]
Magnitude
\[
\lvert \vec{PQ} \rvert = \sqrt{4 + 4 + 1} = 3
\]
Direction cosines:
\[
\cos\alpha = \frac{2}{3} \implies \alpha = \cos^{-1}\left(\frac{2}{3}\right)
\]
\[
\cos\beta = -\frac{2}{3}
\]
. For the acute angle, take the absolute value:
\[
\beta = \cos^{-1}\left(\frac{2}{3}\right)
\]
\[
\cos\gamma = -\frac{1}{3}
\]
. For the acute angle, take the absolute value:
\[
\gamma = \cos^{-1}\left(\frac{1}{3}\right)
\]

\subsection{Exam Patterns}\label{exam-patterns}

\begin{itemize}
\tightlist
\item
  Direction cosine questions often ask for the acute angle. If your
  cosine is negative, drop the minus sign before taking the inverse
  cosine to find the acute angle relative to that axis line.
\item
  Basic definitions (scalar vs vector) sometimes appear as 1-2 mark
  theoretical add-ons to larger questions.
\end{itemize}

\section{A02: Dot Product and Cross
Product}\label{a02-dot-product-and-cross-product}

\begin{quote}
\textbf{Exam Weight}: Tier 2 \textbar{} \textbf{Appeared in}: 2020,
2021, 2023, CT \textbar{} \textbf{Typical Marks}: 3--4
\end{quote}

The dot product and cross product are the two fundamental ways to
multiply vectors. They have distinct physical interpretations and
mathematical properties.

\subsection{Dot Product (Scalar
Product)}\label{dot-product-scalar-product}

\textbf{Definition}:
\[
\vec{A} \cdot \vec{B} = \lvert \vec{A} \rvert\lvert \vec{B} \rvert\cos\theta
\]
\begin{itemize}
\tightlist
\item
  \textbf{Output}: A scalar (a number).
\item
  \textbf{Physical Significance}: Represents the projection of one
  vector onto another. It is used to calculate work done by a force (
\end{itemize}
\[
W = \vec{F} \cdot \vec{d}
\]
). - \textbf{Orthogonality}: If two non-zero vectors are perpendicular
(orthogonal), their dot product is 0.

\subsubsection{Worked Example: Mutually Orthogonal Vectors
(CT)}\label{worked-example-mutually-orthogonal-vectors-ct}

\textbf{Problem}: Show that
\[
\vec{A} = \frac{1}{3}(2\hat{i} - 2\hat{j} + \hat{k}),
\]
\[
\vec{B} = \frac{1}{3}(\hat{i} + 2\hat{j} + 2\hat{k}),
\]
\[
\vec{C} = \frac{1}{3}(2\hat{i} + \hat{j} - 2\hat{k})
\]
are mutually orthogonal unit vectors. \textbf{Solution}: Magnitude check
for unit vectors:
\[
\lvert \vec{A} \rvert = \sqrt{(2/3)^2 + (-2/3)^2 + (1/3)^2} = \sqrt{9/9} = 1
\]
. (Same applies for \(\vec{B}\) and \(\vec{C}\)). Dot product check for
orthogonality:
\[
\vec{A} \cdot \vec{B} = \frac{1}{9}(2(1) + (-2)(2) + 1(2)) = \frac{1}{9}(2 - 4 + 2) = 0.
\]
\[
\vec{B} \cdot \vec{C} = \frac{1}{9}(1(2) + 2(1) + 2(-2)) = \frac{1}{9}(2 + 2 - 4) = 0.
\]
\[
\vec{C} \cdot \vec{A} = \frac{1}{9}(2(2) + 1(-2) + (-2)(1)) = \frac{1}{9}(4 - 2 - 2) = 0.
\]
Since all magnitudes are 1 and all dot products are 0, they are mutually
orthogonal unit vectors.

\subsection{Cross Product (Vector
Product)}\label{cross-product-vector-product}

\textbf{Definition}:
\[
\vec{A} \times \vec{B} = \lvert \vec{A} \rvert\lvert \vec{B} \rvert\sin\theta \, \hat{n}
\]
\begin{itemize}
\tightlist
\item
  \textbf{Output}: A new vector, perpendicular to the plane containing
  \(\vec{A}\) and \(\vec{B}\).
\item
  \textbf{Physical Significance}: Represents a rotational effect. It is
  used to calculate torque (
\end{itemize}
\[
\vec{\tau} = \vec{r} \times \vec{F}
\]
). The magnitude equals the area of the parallelogram formed by the two
vectors. - \textbf{Parallelism}: If two non-zero vectors are parallel,
their cross product is \(\vec{0}\).

\subsubsection{Worked Example: Finding Perpendicular Vectors (PYQ
2020)}\label{worked-example-finding-perpendicular-vectors-pyq-2020}

\textbf{Problem}: Find a vector of magnitude 5 perpendicular to both
\[
\vec{A} = 2\hat{i} + \hat{j} - 3\hat{k}
\]
and
\[
\vec{B} = \hat{i} - 2\hat{j} + \hat{k}.
\]
\textbf{Solution}: The cross product yields a perpendicular vector:
\[
\vec{A} \times \vec{B} =
\begin{vmatrix}
\hat{i} & \hat{j} & \hat{k} \\
2 & 1 & -3 \\
1 & -2 & 1
\end{vmatrix} = \hat{i}(1 - 6) - \hat{j}(2 - (-3)) + \hat{k}(-4 - 1) = -5\hat{i} - 5\hat{j} - 5\hat{k}
\]
Simplify to a base direction vector
\[
\vec{v} = \hat{i} + \hat{j} + \hat{k}.
\]
Find the unit normal vector \(\hat{n}\):
\[
\hat{n} = \pm \frac{\hat{i} + \hat{j} + \hat{k}}{\sqrt{1^2 + 1^2 + 1^2}} = \pm \frac{\hat{i} + \hat{j} + \hat{k}}{\sqrt{3}}
\]
Multiply by the desired magnitude 5:
\[
\vec{P} = \pm \frac{5}{\sqrt{3}}(\hat{i} + \hat{j} + \hat{k})
\]

\subsection{Exam Patterns}\label{exam-patterns-1}

\begin{itemize}
\tightlist
\item
  When asked to find a vector perpendicular to two other vectors, always
  use the cross product, normalize it to a unit vector, and then
  multiply by the requested magnitude. Include the \(\pm\) sign to show
  both possible directions.
\item
  ``Show vectors are orthogonal'' simply means proving their dot product
  is zero.
\end{itemize}

\section{A03: Area, Volume, and Triple
Products}\label{a03-area-volume-and-triple-products}

\begin{quote}
\textbf{Exam Weight}: Tier 2 \textbar{} \textbf{Appeared in}: 2019,
2020, 2021, 2023 \textbar{} \textbf{Typical Marks}: 3--4
\end{quote}

Triple products allow you to find volumes of 3D shapes or determine if
three vectors lie in the same flat plane (coplanar). Cross products
natively provide the area of 2D shapes.

\subsection{Area of a Triangle}\label{area-of-a-triangle}

The magnitude of the cross product of two vectors is the area of the
parallelogram they form. Therefore, the area of a triangle formed by two
edge vectors is half of that:
\[
\text{Area} = \frac{1}{2}\lvert \vec{A} \times \vec{B} \rvert
\]
\subsubsection{Worked Example (PYQ 2021)}\label{worked-example-pyq-2021}

\textbf{Problem}: Find the area of a triangle with vertices at
\(P(3, -1, 2)\), \(Q(1, -1, -3)\) and \(R(4, -3, 1)\).
\textbf{Solution}: Construct two edge vectors originating from the same
point \(P\):
\[
\vec{PQ} = (1 - 3)\hat{i} + (-1 - (-1))\hat{j} + (-3 - 2)\hat{k} = -2\hat{i} - 5\hat{k}
\]
\[
\vec{PR} = (4 - 3)\hat{i} + (-3 - (-1))\hat{j} + (1 - 2)\hat{k} = \hat{i} - 2\hat{j} - \hat{k}
\]
Calculate the cross product:
\[
\vec{PQ} \times \vec{PR} =
\begin{vmatrix}
\hat{i} & \hat{j} & \hat{k} \\
-2 & 0 & -5 \\
1 & -2 & -1
\end{vmatrix} = \hat{i}(0 - 10) - \hat{j}(2 - (-5)) + \hat{k}(4 - 0) = -10\hat{i} - 7\hat{j} + 4\hat{k}
\]
Find magnitude and area:
\[
\lvert \vec{PQ} \times \vec{PR} \rvert = \sqrt{(-10)^2 + (-7)^2 + 4^2} = \sqrt{100 + 49 + 16} = \sqrt{165}
\]
Area
\[
= \frac{\sqrt{165}}{2}
\]

\subsection{Scalar Triple Product and
Volume}\label{scalar-triple-product-and-volume}

The scalar triple product
\[
\vec{A} \cdot (\vec{B} \times \vec{C})
\]
calculates the volume of a parallelepiped (a 3D box with sheared sides)
whose adjacent edges are \(\vec{A}, \vec{B}, \vec{C}\).

It is evaluated as a \(3 \times 3\) determinant:
\[
\vec{A} \cdot (\vec{B} \times \vec{C}) =
\begin{vmatrix}
A_x & A_y & A_z \\
B_x & B_y & B_z \\
C_x & C_y & C_z
\end{vmatrix}
\]
\textbf{Geometric Proof}: The cross product \(\vec{B} \times \vec{C}\)
gives a normal vector whose magnitude is the area of the base. The dot
product
\[
\vec{A} \cdot (\vec{B} \times \vec{C})
\]
projects the third edge \(\vec{A}\) onto this normal vector, effectively
multiplying the base area by the perpendicular height. Thus:
\[
\text{Volume} = \lvert \vec{A} \cdot (\vec{B} \times \vec{C}) \rvert
\]
\subsubsection{Worked Example: Volume (PYQ
2019)}\label{worked-example-volume-pyq-2019}

\textbf{Problem}: Find the volume of the parallelepiped whose edges are
\[
\vec{A} = 2\hat{i} - 3\hat{j} + 4\hat{k},
\]
\[
\vec{B} = \hat{i} + 2\hat{j} - \hat{k},
\]
\[
\vec{C} = 3\hat{i} - \hat{j} + 2\hat{k}.
\]
\textbf{Solution}:
\[
\vec{A} \cdot (\vec{B} \times \vec{C}) =
\begin{vmatrix}
2 & -3 & 4 \\
1 & 2 & -1 \\
3 & -1 & 2
\end{vmatrix}
\]
\[
= 2(4 - 1) - (-3)(2 - (-3)) + 4(-1 - 6) = 2(3) + 3(5) + 4(-7) = 6 + 15 - 28 = -7
\]
Volume is the absolute value:
\[
\lvert -7 \rvert = 7
\]
cubic units.

\subsection{Coplanar Vectors}\label{coplanar-vectors}

If three vectors lie in the exact same plane, the 3D parallelepiped they
form has a height of zero. Therefore, its volume is zero.

\textbf{Condition for Coplanarity}: Three vectors are coplanar if and
only if their scalar triple product is zero.

\subsubsection{Worked Example: Unknown Constant (PYQ
2023)}\label{worked-example-unknown-constant-pyq-2023}

\textbf{Problem}: Find the constant \(a\) such that
\[
2\hat{i} - \hat{j} + \hat{k},
\]
\[
\hat{i} + 2\hat{j} - 3\hat{k},
\]
\[
3\hat{i} + a\hat{j} + 5\hat{k}
\]
are coplanar. \textbf{Solution}: Set the scalar triple product
determinant to 0:
\[
\begin{vmatrix}
2 & -1 & 1 \\
1 & 2 & -3 \\
3 & a & 5
\end{vmatrix} = 0
\]
\[
2(10 - (-3a)) - (-1)(5 - (-9)) + 1(a - 6) = 0
\]
\[
2(10 + 3a) + 14 + a - 6 = 0 \implies 20 + 6a + 14 + a - 6 = 0 \implies 7a + 28 = 0 \implies a = -4
\]

\subsection{Exam Patterns}\label{exam-patterns-2}

\begin{itemize}
\tightlist
\item
  When given three points and asked for an area, always form two vectors
  originating from the \textbf{same point} (e.g., \(\vec{PQ}\) and
  \(\vec{PR}\)). Do not cross arbitrary position vectors.
\item
  Scalar triple product calculations are very prone to arithmetic sign
  errors. Always write out the \(2 \times 2\) determinant expansions
  carefully.
\end{itemize}

\section{A04: Vector Differentiation}\label{a04-vector-differentiation}

\begin{quote}
\textbf{Exam Weight}: Tier 3 \textbar{} \textbf{Appeared in}: 2019, 2024
\textbar{} \textbf{Typical Marks}: 3--4
\end{quote}

Differentiating a vector function works exactly like scalar
differentiation, component by component. The product rule still applies
to both dot and cross products.

\subsection{The Basics}\label{the-basics}

If a position vector is given as a function of time \(t\):
\[
\vec{r}(t) = x(t)\hat{i} + y(t)\hat{j} + z(t)\hat{k}
\]
The derivative with respect to \(t\) is:
\[
\frac{d\vec{r}}{dt} = \frac{dx}{dt}\hat{i} + \frac{dy}{dt}\hat{j} + \frac{dz}{dt}\hat{k}
\]
\subsubsection{Product Rules}\label{product-rules}

When differentiating the product of two vector functions \(\vec{A}(t)\)
and \(\vec{B}(t)\):

\textbf{Dot Product Rule}:
\[
\frac{d}{dt} (\vec{A} \cdot \vec{B}) = \vec{A} \cdot \frac{d\vec{B}}{dt} + \frac{d\vec{A}}{dt} \cdot \vec{B}
\]
\textbf{Cross Product Rule} (Order matters!):
\[
\frac{d}{dt} (\vec{A} \times \vec{B}) = \vec{A} \times \frac{d\vec{B}}{dt} + \frac{d\vec{A}}{dt} \times \vec{B}
\]

\subsection{Must-Know Proof: Constant Magnitude (PYQ
2019)}\label{must-know-proof-constant-magnitude-pyq-2019}

\textbf{Problem}: If \(\vec{A}\) has a constant magnitude, show that
\(\vec{A}\) and \(\frac{d\vec{A}}{dt}\) are perpendicular, provided
\[
\left\lvert \frac{d\vec{A}}{dt}\right\rvert  \neq 0.
\]
\textbf{Proof}: Let the constant magnitude of \(\vec{A}\) be
\[
\lvert \vec{A} \rvert = c
\]
(where \(c\) is a constant). We know that the dot product of a vector
with itself is its magnitude squared:
\[
\vec{A} \cdot \vec{A} = \lvert \vec{A} \rvert^2 = c^2
\]
Differentiate both sides with respect to \(t\):
\[
\frac{d}{dt} (\vec{A} \cdot \vec{A}) = \frac{d}{dt} (c^2)
\]
Apply the dot product rule to the left side, and note that the
derivative of a constant is 0:
\[
\vec{A} \cdot \frac{d\vec{A}}{dt} + \frac{d\vec{A}}{dt} \cdot \vec{A} = 0
\]
Since the dot product is commutative (
\[
\vec{u} \cdot \vec{v} = \vec{v} \cdot \vec{u}
\]
):
\[
2 \left( \vec{A} \cdot \frac{d\vec{A}}{dt} \right) = 0 \implies \vec{A} \cdot \frac{d\vec{A}}{dt} = 0
\]
Because their dot product is zero and neither vector is the zero vector,
\(\vec{A}\) and \(\frac{d\vec{A}}{dt}\) must be mutually perpendicular.
\(\blacksquare\)

\subsection{Exam Patterns}\label{exam-patterns-3}

\begin{itemize}
\tightlist
\item
  The ``constant magnitude implies perpendicular derivative'' proof is a
  classic theoretical question. Memorize the 4-step proof above.
\item
  The fundamental concept of taking a derivative component-by-component
  is heavily used in the next topic: finding velocity and acceleration.
\end{itemize}

\section{A05: Velocity, Acceleration, and Particle
Motion}\label{a05-velocity-acceleration-and-particle-motion}

\begin{quote}
\textbf{Exam Weight}: Tier 2 \textbar{} \textbf{Appeared in}: 2019,
2020, 2023, 2024 \textbar{} \textbf{Typical Marks}: 4
\end{quote}

A particle's motion in space is described by its position vector
\(\vec{r}(t)\). Taking derivatives with respect to time \(t\) gives its
velocity and acceleration.

\subsection{The Core Formulas}\label{the-core-formulas}

Given a position vector \(\vec{r}(t)\): 1. \textbf{Velocity}:
\[
\vec{v}(t) = \frac{d\vec{r}}{dt}
\]
\begin{enumerate}
\def\labelenumi{\arabic{enumi}.}
\setcounter{enumi}{1}
\tightlist
\item
  \textbf{Acceleration}:
\end{enumerate}
\[
\vec{a}(t) = \frac{d\vec{v}}{dt} = \frac{d^2\vec{r}}{dt^2}
\]
If asked for a component of velocity or acceleration in a specific
direction, take the dot product of your result with the \textbf{unit
vector} \(\hat{u}\) of that direction.

\subsection{Worked Example 1: Directional Components (PYQ 2019,
2023)}\label{worked-example-1-directional-components-pyq-2019-2023}

\textbf{Problem}: A particle moves along
\[
x = 2t^2, y = t^2 - 4t, z = 3t - 5
\]
. Find the components of velocity and acceleration at \(t=1\) in the
direction of
\[
\vec{d} = \hat{i} - 3\hat{j} + 2\hat{k}.
\]
\textbf{Solution}: \textbf{Step 1: Setup} Position:
\[
\vec{r}(t) = 2t^2\hat{i} + (t^2 - 4t)\hat{j} + (3t - 5)\hat{k}
\]
Direction Unit Vector:
\[
\hat{u} = \frac{\hat{i} - 3\hat{j} + 2\hat{k}}{\sqrt{1^2 + (-3)^2 + 2^2}} = \frac{\hat{i} - 3\hat{j} + 2\hat{k}}{\sqrt{14}}
\]
\textbf{Step 2: Velocity}
\[
\vec{v}(t) = \frac{d\vec{r}}{dt} = 4t\hat{i} + (2t - 4)\hat{j} + 3\hat{k}
\]
At \(t=1\):
\[
\vec{v}(1) = 4\hat{i} - 2\hat{j} + 3\hat{k}
\]
Component in direction \(\hat{u}\):
\[
v_d = \vec{v}(1) \cdot \hat{u} = (4\hat{i} - 2\hat{j} + 3\hat{k}) \cdot \left( \frac{\hat{i} - 3\hat{j} + 2\hat{k}}{\sqrt{14}} \right) = \frac{4(1) - 2(-3) + 3(2)}{\sqrt{14}} = \frac{4 + 6 + 6}{\sqrt{14}} = \frac{16}{\sqrt{14}}
\]
\textbf{Step 3: Acceleration}
\[
\vec{a}(t) = \frac{d\vec{v}}{dt} = 4\hat{i} + 2\hat{j}
\]
At \(t=1\):
\[
\vec{a}(1) = 4\hat{i} + 2\hat{j}
\]
(Acceleration is constant) Component in direction \(\hat{u}\):
\[
a_d = \vec{a}(1) \cdot \hat{u} = (4\hat{i} + 2\hat{j}) \cdot \left( \frac{\hat{i} - 3\hat{j} + 2\hat{k}}{\sqrt{14}} \right) = \frac{4(1) + 2(-3) + 0}{\sqrt{14}} = \frac{4 - 6}{\sqrt{14}} = -\frac{2}{\sqrt{14}}
\]

\subsection{Worked Example 2: Trigonometric Motion (PYQ
2024)}\label{worked-example-2-trigonometric-motion-pyq-2024}

\textbf{Problem}: A particle moves so that its position vector is
\[
\vec{r} = \cos \omega t \hat{i} + \sin \omega t \hat{j}
\]
where \(\omega\) is constant. Show that velocity \(\vec{v}\) is
perpendicular to \(\vec{r}\), and that \(\vec{r} \times \vec{v}\) is a
constant vector.

\textbf{Solution}: \textbf{Step 1: Velocity}
\[
\vec{v} = \frac{d\vec{r}}{dt} = -\omega \sin \omega t \hat{i} + \omega \cos \omega t \hat{j}
\]
\textbf{Step 2: Show Perpendicular}
\[
\vec{r} \cdot \vec{v} = (\cos \omega t)(-\omega \sin \omega t) + (\sin \omega t)(\omega \cos \omega t) = -\omega \sin \omega t \cos \omega t + \omega \sin \omega t \cos \omega t = 0
\]
Since the dot product is zero, \(\vec{v}\) is perpendicular to
\(\vec{r}\).

\textbf{Step 3: Show Constant Cross Product}
\[
\vec{r} \times \vec{v} =
\begin{vmatrix}
\hat{i} & \hat{j} & \hat{k} \\
\cos \omega t & \sin \omega t & 0 \\
-\omega \sin \omega t & \omega \cos \omega t & 0
\end{vmatrix}
\]
\[
= \hat{k} \left[ (\cos \omega t)(\omega \cos \omega t) - (\sin \omega t)(-\omega \sin \omega t) \right] = \hat{k}(\omega \cos^2 \omega t + \omega \sin^2 \omega t)
\]
\[
= \omega(\cos^2 \omega t + \sin^2 \omega t)\hat{k} = \omega\hat{k}
\]
Since \(\omega\) is a constant, \(\omega\hat{k}\) is a constant vector.

\subsection{Exam Patterns}\label{exam-patterns-4}

\begin{itemize}
\tightlist
\item
  When asked for a ``component in the direction of'', do not forget to
  divide the direction vector by its magnitude to make it a unit vector
  before dotting.
\item
  Simple derivatives involving \(e^t\), \(\sin\), and \(\cos\) are
  common. Remember the chain rule for arguments like \(\omega t\) or
  \(3t\).
\end{itemize}

\section{A06: Curvature, Torsion, and
Frenet-Serret}\label{a06-curvature-torsion-and-frenet-serret}

\begin{quote}
\textbf{Exam Weight}: Tier 3 \textbar{} \textbf{Appeared in}: 2020,
2023, CT \textbar{} \textbf{Typical Marks}: 4--5
\end{quote}

This section analyzes the geometric properties of a 3D curve in space.
It relies heavily on taking multiple derivatives and cross products of
the position vector \(\vec{r}(t)\).

\subsection{Formulas for Curvature and
Torsion}\label{formulas-for-curvature-and-torsion}

Let \(\vec{r}(t)\) be the position vector of a curve parameterized by
\(t\). You need the first three derivatives: 1. \(\vec{r}'\) (Velocity)
2. \(\vec{r}''\) (Acceleration) 3. \(\vec{r}'''\) (Jerk)

\textbf{Curvature (\(K\))}: Measures how sharply a curve bends.
\[
K = \frac{\lvert \vec{r}' \times \vec{r}'' \rvert}{\lvert \vec{r}' \rvert^3}
\]
\textbf{Torsion (\(T\))}: Measures how sharply a curve twists out of its
plane.
\[
T = \frac{(\vec{r}' \times \vec{r}'') \cdot \vec{r}'''}{\lvert \vec{r}' \times \vec{r}'' \rvert^2}
\]
\emph{(Note: The numerator is a scalar triple product).}

\subsection{Frenet-Serret Formulas}\label{frenet-serret-formulas}

The Frenet-Serret formulas describe the kinematic properties of a
particle moving along a continuous, differentiable curve in 3D space.

Let \(s\) be the arc length along a space curve. - \(\vec{T}\) = unit
tangent vector - \(\vec{N}\) = unit principal normal vector -
\(\vec{B}_b\) = unit binormal vector - \(\kappa\) = curvature - \(\tau\)
= torsion

The three formulas are: 1.
\[
\frac{d\vec{T}}{ds} = \kappa \vec{N}
\]
\begin{enumerate}
\def\labelenumi{\arabic{enumi}.}
\setcounter{enumi}{1}
\tightlist
\item
\end{enumerate}
\[
\frac{d\vec{N}}{ds} = -\kappa \vec{T} + \tau \vec{B}_b
\]
\begin{enumerate}
\def\labelenumi{\arabic{enumi}.}
\setcounter{enumi}{2}
\tightlist
\item
\end{enumerate}
\[
\frac{d\vec{B}_b}{ds} = -\tau \vec{N}
\]

\subsection{Worked Example: Calculating and (PYQ 2020,
2023)}\label{worked-example-calculating-and-pyq-2020-2023}

\textbf{Problem}: For the space curve
\[
x = t, y = t^2, z = \frac{2}{3}t^3
\]
, find the curvature \(K\) and the torsion \(T\).

\textbf{Solution}: \textbf{Step 1: Find Derivatives}
\[
\vec{r}(t) = t\hat{i} + t^2\hat{j} + \frac{2}{3}t^3\hat{k}
\]
\[
\vec{r}' = \hat{i} + 2t\hat{j} + 2t^2\hat{k}
\]
\[
\vec{r}'' = 2\hat{j} + 4t\hat{k}
\]
\[
\vec{r}''' = 4\hat{k}
\]
\textbf{Step 2: Cross Product and Magnitudes}
\[
\vec{r}' \times \vec{r}'' =
\begin{vmatrix}
\hat{i} & \hat{j} & \hat{k} \\
1 & 2t & 2t^2 \\
0 & 2 & 4t
\end{vmatrix} = \hat{i}(8t^2 - 4t^2) - \hat{j}(4t - 0) + \hat{k}(2 - 0) = 4t^2\hat{i} - 4t\hat{j} + 2\hat{k}
\]
Magnitude \(\lvert \vec{r}' \rvert\):
\[
\lvert \vec{r}' \rvert = \sqrt{1^2 + (2t)^2 + (2t^2)^2} = \sqrt{1 + 4t^2 + 4t^4} = \sqrt{(2t^2 + 1)^2} = 2t^2 + 1
\]
Magnitude
\[
\lvert \vec{r}' \times \vec{r}'' \rvert:
\]
\[
\lvert \vec{r}' \times \vec{r}'' \rvert = \sqrt{(4t^2)^2 + (-4t)^2 + 2^2} = \sqrt{16t^4 + 16t^2 + 4} = \sqrt{4(4t^4 + 4t^2 + 1)} = 2(2t^2 + 1)
\]
\textbf{Step 3: Calculate Curvature (\(K\))}
\[
K = \frac{\lvert \vec{r}' \times \vec{r}'' \rvert}{\lvert \vec{r}' \rvert^3} = \frac{2(2t^2 + 1)}{(2t^2 + 1)^3} = \frac{2}{(2t^2 + 1)^2}
\]
\textbf{Step 4: Calculate Torsion (\(T\))} Calculate the scalar triple
product
\[
(\vec{r}' \times \vec{r}'') \cdot \vec{r}''':
\]
\[
(4t^2\hat{i} - 4t\hat{j} + 2\hat{k}) \cdot (4\hat{k}) = 8
\]
\[
T = \frac{(\vec{r}' \times \vec{r}'') \cdot \vec{r}'''}{\lvert \vec{r}' \times \vec{r}'' \rvert^2} = \frac{8}{[2(2t^2 + 1)]^2} = \frac{8}{4(2t^2 + 1)^2} = \frac{2}{(2t^2 + 1)^2}
\]
In this specific curve, curvature equals torsion!

\subsection{Exam Patterns}\label{exam-patterns-5}

\begin{itemize}
\tightlist
\item
  This specific problem (
\end{itemize}
\[
x=t, y=t^2, z=\frac{2}{3}t^3
\]
) is practically the only curvature/torsion problem asked. The trick is
recognizing that \(1 + 4t^2 + 4t^4\) is a perfect square
\((2t^2 + 1)^2\). - Be able to state the Frenet-Serret formulas
verbatim.

\section{A07: Gradient and Normal to a
Surface}\label{a07-gradient-and-normal-to-a-surface}

\begin{quote}
\textbf{Exam Weight}: Tier 1 \textbar{} \textbf{Appeared in}: 2017,
2018, 2019, 2020, 2021, 2024 (6 papers) \textbar{} \textbf{Typical
Marks}: 2--5
\end{quote}

The gradient is the foundation of vector calculus. It appears directly
(find the gradient) and indirectly (in directional derivatives, surface
normals, angle between surfaces).

\subsection{The Del Operator}\label{the-del-operator}

The del operator (nabla) is the vector differential operator:
\[
\nabla = \hat{i}\frac{\partial}{\partial x} + \hat{j}\frac{\partial}{\partial y} + \hat{k}\frac{\partial}{\partial z}
\]
It acts on scalar functions (gradient), dot-products with vectors
(divergence), and cross-products with vectors (curl).

\subsection{The Gradient}\label{the-gradient}

\subsubsection{Definition}\label{definition}

The gradient of a scalar function \(\phi(x, y, z)\) is:
\[
\nabla\phi = \frac{\partial\phi}{\partial x}\hat{i} + \frac{\partial\phi}{\partial y}\hat{j} + \frac{\partial\phi}{\partial z}\hat{k}
\]
Input: a scalar function. Output: a vector field.

\subsubsection{Geometric Meaning}\label{geometric-meaning}

The gradient vector \(\nabla\phi\) at any point is \textbf{perpendicular
to the surface} \(\phi(x, y, z) = c\) at that point.

The magnitude \(\lvert \nabla\phi \rvert\) gives the maximum rate of
change of \(\phi\) at that point.

The direction of \(\nabla\phi\) points toward the direction of greatest
increase.

\subsubsection{Properties}\label{properties}

\begin{enumerate}
\def\labelenumi{\arabic{enumi}.}
\tightlist
\item
\end{enumerate}
\[
\nabla(f \pm g) = \nabla f \pm \nabla g
\]
\begin{enumerate}
\def\labelenumi{\arabic{enumi}.}
\setcounter{enumi}{1}
\tightlist
\item
\end{enumerate}
\[
\nabla(fg) = f\nabla g + g\nabla f
\]
\begin{enumerate}
\def\labelenumi{\arabic{enumi}.}
\setcounter{enumi}{2}
\tightlist
\item
\end{enumerate}
\[
\nabla\left(\frac{f}{g}\right) = \frac{g\nabla f - f\nabla g}{g^2}
\]
\begin{enumerate}
\def\labelenumi{\arabic{enumi}.}
\setcounter{enumi}{3}
\tightlist
\item
  \(\nabla r^n\):
\end{enumerate}
\[
\nabla r^n = nr^{n-2}\vec{r} \quad (\text{where } r = \lvert \vec{r} \rvert)
\]

\subsection{Proof: Gradient is Normal to a Surface (PYQ
2020)}\label{proof-gradient-is-normal-to-a-surface-pyq-2020}

\subsubsection{Statement}\label{statement}

\(\nabla\phi\) is perpendicular to the surface \(\phi(x, y, z) = c\).

\subsubsection{Proof}\label{proof}

Let
\[
\vec{r}(t) = x(t)\hat{i} + y(t)\hat{j} + z(t)\hat{k}
\]
be any curve lying on the surface \(\phi = c\).

Since all points of this curve satisfy \(\phi(x(t), y(t), z(t)) = c\),
differentiate with respect to \(t\):
\[
\frac{\partial\phi}{\partial x}\frac{dx}{dt} + \frac{\partial\phi}{\partial y}\frac{dy}{dt} + \frac{\partial\phi}{\partial z}\frac{dz}{dt} = 0
\]
This is a dot product:
\[
\nabla\phi \cdot \frac{d\vec{r}}{dt} = 0
\]
Since \(\frac{d\vec{r}}{dt}\) is tangent to the curve, and the dot
product is zero, \(\nabla\phi\) is perpendicular to every tangent vector
on the surface. So \(\nabla\phi\) is normal to the surface.
\[
\blacksquare
\]

\subsection{Unit Normal Vector}\label{unit-normal-vector}

The unit normal to a surface \(\phi = c\) at a point is:
\[
\hat{n} = \frac{\nabla\phi}{|\nabla\phi|}
\]
\subsubsection{Worked Example: Unit Normal
(CT)}\label{worked-example-unit-normal-ct}

\textbf{Problem}: Find the unit outward normal to
\[
(x-1)^2 + y^2 + (z+2)^2 = 9
\]
at \((3, 1, -4)\).

\textbf{Solution}:

Let
\[
f = (x-1)^2 + y^2 + (z+2)^2 - 9.
\]
\[
\nabla f = 2(x-1)\hat{i} + 2y\hat{j} + 2(z+2)\hat{k}
\]
At \((3, 1, -4)\):
\[
\nabla f = 4\hat{i} + 2\hat{j} - 4\hat{k}.
\]
Magnitude:
\[
\lvert \nabla f \rvert = \sqrt{16 + 4 + 16} = 6.
\]
\[
\hat{n} = \frac{4\hat{i} + 2\hat{j} - 4\hat{k}}{6} = \frac{2}{3}\hat{i} + \frac{1}{3}\hat{j} - \frac{2}{3}\hat{k}
\]

\subsection{Angle Between Two Surfaces (PYQ 2017, 2019, 2021 --- 3
Times)}\label{angle-between-two-surfaces-pyq-2017-2019-2021-3-times}

\subsubsection{The Method}\label{the-method}

The angle between two surfaces
\[
\phi_1 = c_1
\]
and
\[
\phi_2 = c_2
\]
at a point is the angle between their normal vectors:
\[
\cos\theta = \frac{\nabla\phi_1 \cdot \nabla\phi_2}{|\nabla\phi_1||\nabla\phi_2|}
\]
\subsubsection{Must-Know Problem (Appeared 3
Times)}\label{must-know-problem-appeared-3-times}

\textbf{Problem}: Find the angle between
\[
x^2 + y^2 + z^2 = 9
\]
and
\[
z = x^2 + y^2 - 3
\]
at \((2, -1, 2)\).

\textbf{Solution}:

Surface 1:
\[
f_1 = x^2 + y^2 + z^2 - 9
\]
. Gradient:
\[
\nabla f_1 = 2x\hat{i} + 2y\hat{j} + 2z\hat{k}.
\]
At \((2, -1, 2)\):
\[
\vec{n}_1 = 4\hat{i} - 2\hat{j} + 4\hat{k}.
\]
Surface 2:
\[
f_2 = x^2 + y^2 - z - 3
\]
. Gradient:
\[
\nabla f_2 = 2x\hat{i} + 2y\hat{j} - \hat{k}.
\]
At \((2, -1, 2)\):
\[
\vec{n}_2 = 4\hat{i} - 2\hat{j} - \hat{k}.
\]
Dot product:
\[
\vec{n}_1 \cdot \vec{n}_2 = 16 + 4 - 4 = 16.
\]
Magnitudes:
\[
\lvert \vec{n}_1 \rvert = 6,
\]
\[
\lvert \vec{n}_2 \rvert = \sqrt{21}.
\]
\[
\cos\theta = \frac{16}{6\sqrt{21}} = \frac{8}{3\sqrt{21}}
\]
\[
\theta = \cos^{-1}\left(\frac{8}{3\sqrt{21}}\right)
\]

\subsection{Finding the Gradient at a Point (PYQ
2024)}\label{finding-the-gradient-at-a-point-pyq-2024}

\textbf{Problem}: If
\[
Q = 3x^2y - y^3z^2
\]
, find \(\nabla Q\) at \((1, -2, -1)\).

\textbf{Solution}:
\[
\frac{\partial Q}{\partial x} = 6xy, \quad \frac{\partial Q}{\partial y} = 3x^2 - 3y^2z^2, \quad \frac{\partial Q}{\partial z} = -2y^3z
\]
At \((1, -2, -1)\):
\[
\frac{\partial Q}{\partial x} = 6(1)(-2) = -12
\]
\[
\frac{\partial Q}{\partial y} = 3(1) - 3(4)(1) = -9
\]
\[
\frac{\partial Q}{\partial z} = -2(-8)(-1) = -16
\]
\[
\nabla Q = -12\hat{i} - 9\hat{j} - 16\hat{k}
\]

\subsection{Exam Patterns}\label{exam-patterns-6}

\begin{itemize}
\tightlist
\item
  ``Find the angle between two surfaces'' has appeared 3 times (2017,
  2019, 2021). The exact same surfaces. Memorize it.
\item
  ``Find the gradient at a point'' is a 2-3 mark easy question.
\item
  ``Show that gradient is normal to the surface'' appeared in 2020 for 3
  marks.
\item
  The gradient formula is used in directional derivative (A08) and
  tangent plane (A09) problems. It is foundational.
\end{itemize}

\section{A08: Directional Derivative}\label{a08-directional-derivative}

\begin{quote}
\textbf{Exam Weight}: Tier 1 \textbar{} \textbf{Appeared in}: 2017,
2018, 2019, 2021 \textbar{} \textbf{Typical Marks}: 3--4
\end{quote}

The directional derivative measures how fast a scalar function changes
in a specific direction. It builds directly on the gradient.

\subsection{Definition}\label{definition-1}

The directional derivative of \(\phi\) at point \(P\) in the direction
of unit vector \(\hat{u}\) is:
\[
D_{\hat{u}}\phi = \nabla\phi \cdot \hat{u}
\]
It gives the rate of change of \(\phi\) at \(P\) in the direction
\(\hat{u}\).

\subsubsection{Key Facts}\label{key-facts}

\begin{itemize}
\tightlist
\item
  Maximum value of \(D_{\hat{u}}\phi\) is \(\lvert \nabla\phi \rvert\),
  achieved in the direction of \(\nabla\phi\) itself.
\item
  Minimum value is \(-\lvert \nabla\phi \rvert\), in the opposite
  direction.
\item
\end{itemize}
\[
D_{\hat{u}}\phi = 0
\]
when \(\hat{u}\) is perpendicular to \(\nabla\phi\) (moving along the
level surface).

\subsection{The Three-Step Method}\label{the-three-step-method}

\begin{enumerate}
\def\labelenumi{\arabic{enumi}.}
\tightlist
\item
  \textbf{Find the gradient} \(\nabla\phi\) and evaluate at the given
  point.
\item
  \textbf{Find the unit direction vector}
\end{enumerate}
\[
\hat{u} = \frac{\vec{d}}{\lvert \vec{d} \rvert}.
\]
\begin{enumerate}
\def\labelenumi{\arabic{enumi}.}
\setcounter{enumi}{2}
\tightlist
\item
  \textbf{Dot product}:
\end{enumerate}
\[
D_{\hat{u}}\phi = \nabla\phi \cdot \hat{u}.
\]

\subsection{Worked Example 1 (PYQ
2018)}\label{worked-example-1-pyq-2018}

\textbf{Problem}: Find the directional derivative of
\[
\phi = x^2yz + 4xz^2
\]
at \((1, -2, -1)\) in the direction
\[
2\hat{i} - \hat{j} - 2\hat{k}.
\]
\textbf{Solution}:

\textbf{Step 1}: Gradient.
\[
\nabla\phi = (2xyz + 4z^2)\hat{i} + x^2z\hat{j} + (x^2y + 8xz)\hat{k}
\]
At \((1, -2, -1)\):
\[
\nabla\phi = (4 + 4)\hat{i} + (1)(-1)\hat{j} + (-2 - 8)\hat{k} = 8\hat{i} - \hat{j} - 10\hat{k}
\]
\textbf{Step 2}: Unit vector.
\[
\hat{u} = \frac{2\hat{i} - \hat{j} - 2\hat{k}}{\sqrt{4 + 1 + 4}} = \frac{2\hat{i} - \hat{j} - 2\hat{k}}{3}
\]
\textbf{Step 3}: Dot product.
\[
D_{\hat{u}}\phi = \frac{8(2) + (-1)(-1) + (-10)(-2)}{3} = \frac{16 + 1 + 20}{3} = \frac{37}{3}
\]

\subsection{Worked Example 2 (PYQ
2017)}\label{worked-example-2-pyq-2017}

\textbf{Problem}: Find the directional derivative of
\[
\phi = 4e^{2x-y+z}
\]
at \((1, 1, -1)\) toward \((-3, 5, 6)\).

\textbf{Solution}:

\textbf{Step 1}: Direction from \(P(1,1,-1)\) to \(Q(-3,5,6)\):
\[
\vec{d} = -4\hat{i} + 4\hat{j} + 7\hat{k}, \quad |\vec{d}| = \sqrt{16+16+49} = 9
\]
\[
\hat{u} = \frac{-4\hat{i} + 4\hat{j} + 7\hat{k}}{9}
\]
\textbf{Step 2}: Gradient of
\[
\phi = 4e^{2x-y+z}:
\]
\[
\nabla\phi = 4e^{2x-y+z}(2\hat{i} - \hat{j} + \hat{k})
\]
At \((1,1,-1)\): exponent = \(2-1-1 = 0\), so
\[
e^0 = 1:
\]
\[
\nabla\phi = 8\hat{i} - 4\hat{j} + 4\hat{k}
\]
\textbf{Step 3}:
\[
D_{\hat{u}}\phi = \frac{8(-4) + (-4)(4) + 4(7)}{9} = \frac{-32 - 16 + 28}{9} = -\frac{20}{9}
\]

\subsection{Worked Example 3 (PYQ
2021)}\label{worked-example-3-pyq-2021}

\textbf{Problem}: Find the directional derivative of
\[
f = x^2 + xy + z^2
\]
at \(A(1,-1,-1)\) in the direction of \(AB\) where \(B(3, 2, 1)\).

\textbf{Solution}:

Direction:
\[
\vec{AB} = 2\hat{i} + 3\hat{j} + 2\hat{k}.
\]
Unit vector:
\[
\hat{u} = \frac{2\hat{i}+3\hat{j}+2\hat{k}}{\sqrt{17}}.
\]
Gradient:
\[
\nabla f = (2x+y)\hat{i} + x\hat{j} + 2z\hat{k}.
\]
At \((1,-1,-1)\):
\[
\nabla f = \hat{i} + \hat{j} - 2\hat{k}.
\]
\[
D_{\hat{u}}f = \frac{2 + 3 - 4}{\sqrt{17}} = \frac{1}{\sqrt{17}}
\]

\subsection{Worked Example 4 (PYQ
2019)}\label{worked-example-4-pyq-2019}

\textbf{Problem}: Find directional derivative of
\[
\phi = x^2yz + 4xz^2
\]
at \((1, -2, 1)\) in direction
\[
2\hat{i} - \hat{j} - 2\hat{k}.
\]
\textbf{Solution}:

At \((1,-2,1)\):
\[
\nabla\phi = 0\hat{i} + \hat{j} + 6\hat{k}.
\]
\[
D_{\hat{u}}\phi = \frac{0(2) + 1(-1) + 6(-2)}{3} = \frac{-13}{3}
\]

\subsection{Exam Patterns}\label{exam-patterns-7}

\begin{itemize}
\tightlist
\item
  Directional derivative appears in about 4 out of 7 papers.
\item
  The same function
\end{itemize}
\[
\phi = x^2yz + 4xz^2
\]
has appeared in 2018 and 2019 with different points. - Always show the
three steps clearly: gradient, unit vector, dot product. - The direction
is sometimes given as a vector (just normalize it) and sometimes as
``toward point Q'' (compute \(\vec{PQ}\) first, then normalize).

\section{A09: Tangent Plane and Normal
Line}\label{a09-tangent-plane-and-normal-line}

\begin{quote}
\textbf{Exam Weight}: Tier 2 \textbar{} \textbf{Appeared in}: 2018
\textbar{} \textbf{Typical Marks}: 3
\end{quote}

The gradient vector \(\nabla F\) is perpendicular to the surface
\(F(x, y, z) = C\). We use this fact to find the equations of the
tangent plane and the normal line at a given point.

\subsection{Formulas}\label{formulas}

Let a surface be given by \(F(x, y, z) = C\), and let
\(P(x_0, y_0, z_0)\) be a point on this surface.

The normal vector to the surface at \(P\) is the gradient evaluated at
\(P\):
\[
\vec{n} = \nabla F|_P = \frac{\partial F}{\partial x}\hat{i} + \frac{\partial F}{\partial y}\hat{j} + \frac{\partial F}{\partial z}\hat{k}
\]
Let the components of this normal vector be
\[
\vec{n} = A\hat{i} + B\hat{j} + C\hat{k}.
\]
\subsubsection{Equation of the Tangent
Plane}\label{equation-of-the-tangent-plane}

The tangent plane passes through \(P\) and is perpendicular to
\(\vec{n}\). Its equation is:
\[
A(x - x_0) + B(y - y_0) + C(z - z_0) = 0
\]
Or, in terms of partial derivatives:
\[
\frac{\partial F}{\partial x}(x - x_0) + \frac{\partial F}{\partial y}(y - y_0) + \frac{\partial F}{\partial z}(z - z_0) = 0
\]
\subsubsection{Equation of the Normal
Line}\label{equation-of-the-normal-line}

The normal line passes through \(P\) and is parallel to \(\vec{n}\). Its
symmetric equations are:
\[
\frac{x - x_0}{A} = \frac{y - y_0}{B} = \frac{z - z_0}{C}
\]

\subsection{Worked Example 1 (PYQ
2018)}\label{worked-example-1-pyq-2018-1}

\textbf{Problem}: Find an equation for the tangent plane to the surface
\[
2xz^2 - 3xy - 4x = 7
\]
at the point \((1, -1, 2)\).

\textbf{Solution}:

\textbf{Step 1}: Define the surface function \(F(x, y, z)\).
\[
F(x, y, z) = 2xz^2 - 3xy - 4x - 7 = 0
\]
\textbf{Step 2}: Find the partial derivatives (the components of the
gradient).
\[
\frac{\partial F}{\partial x} = 2z^2 - 3y - 4
\]
\[
\frac{\partial F}{\partial y} = -3x
\]
\[
\frac{\partial F}{\partial z} = 4xz
\]
\textbf{Step 3}: Evaluate the partial derivatives at the point
\((1, -1, 2)\).
\[
A = \frac{\partial F}{\partial x}(1, -1, 2) = 2(2)^2 - 3(-1) - 4 = 8 + 3 - 4 = 7
\]
\[
B = \frac{\partial F}{\partial y}(1, -1, 2) = -3(1) = -3
\]
\[
C = \frac{\partial F}{\partial z}(1, -1, 2) = 4(1)(2) = 8
\]
The normal vector is
\[
\vec{n} = 7\hat{i} - 3\hat{j} + 8\hat{k}.
\]
\textbf{Step 4}: Write the equation of the tangent plane.

Using the formula
\[
A(x - x_0) + B(y - y_0) + C(z - z_0) = 0:
\]
\[
7(x - 1) - 3(y - (-1)) + 8(z - 2) = 0
\]
\[
7(x - 1) - 3(y + 1) + 8(z - 2) = 0
\]
\[
7x - 7 - 3y - 3 + 8z - 16 = 0
\]
\[
7x - 3y + 8z = 26
\]
The equation of the tangent plane is \(7x - 3y + 8z = 26\).

*(Note: If the question also asked for the normal line, it would be:
\[
\frac{x-1}{7} = \frac{y+1}{-3} = \frac{z-2}{8}
\]
)*

\subsection{Exam Patterns}\label{exam-patterns-8}

\begin{itemize}
\tightlist
\item
  ``Find the equation of the tangent plane'' is a straightforward 3-mark
  application of the gradient.
\item
  Always bring all terms to one side to define \(F(x, y, z) = 0\) before
  taking derivatives.
\item
  The normal vector components \((A, B, C)\) become the coefficients of
  \(x, y, z\) in the final plane equation.
\end{itemize}

\section{A10: Divergence and Curl}\label{a10-divergence-and-curl}

\begin{quote}
\textbf{Exam Weight}: Tier 1 \textbar{} \textbf{Appeared in}: Core
concept used in almost every vector calculus question.
\end{quote}

While gradient operates on scalars, divergence and curl operate on
vector fields. They describe the physical behavior of the field
(expansion/compression and rotation).

\subsection{Divergence (Dot Product)}\label{divergence-dot-product}

\subsubsection{Definition}\label{definition-2}

The divergence of a vector field
\[
\vec{F} = F_1\hat{i} + F_2\hat{j} + F_3\hat{k}
\]
is the dot product of the Del operator and \(\vec{F}\):
\[
\text{div }\vec{F} = \nabla \cdot \vec{F} = \left(\hat{i}\frac{\partial}{\partial x} + \hat{j}\frac{\partial}{\partial y} + \hat{k}\frac{\partial}{\partial z}\right) \cdot (F_1\hat{i} + F_2\hat{j} + F_3\hat{k})
\]
\[
\text{div }\vec{F} = \frac{\partial F_1}{\partial x} + \frac{\partial F_2}{\partial y} + \frac{\partial F_3}{\partial z}
\]
Input: Vector field. Output: \textbf{Scalar quantity}.

\subsubsection{Physical Significance}\label{physical-significance}

Divergence represents the net outward flux (flow) per unit volume from a
point. - \textbf{Positive divergence}: The point is a source (fluid
expanding outward). - \textbf{Negative divergence}: The point is a sink
(fluid compressing inward).

\subsubsection{Solenoidal Vector}\label{solenoidal-vector}

If
\[
\nabla \cdot \vec{F} = 0
\]
everywhere, the vector field is called \textbf{solenoidal}. It behaves
like an incompressible fluid with no sources or sinks.

\subsection{Curl (Cross Product)}\label{curl-cross-product}

\subsubsection{Definition}\label{definition-3}

The curl of a vector field \(\vec{F}\) is the cross product of the Del
operator and \(\vec{F}\):
\[
\text{curl }\vec{F} = \nabla \times \vec{F} =
\begin{vmatrix}
\hat{i} & \hat{j} & \hat{k} \\
\frac{\partial}{\partial x} & \frac{\partial}{\partial y} & \frac{\partial}{\partial z} \\
F_1 & F_2 & F_3
\end{vmatrix}
\]
Expanding this gives:
\[
\nabla \times \vec{F} = \hat{i}\left( \frac{\partial F_3}{\partial y} - \frac{\partial F_2}{\partial z} \right) - \hat{j}\left( \frac{\partial F_3}{\partial x} - \frac{\partial F_1}{\partial z} \right) + \hat{k}\left( \frac{\partial F_2}{\partial x} - \frac{\partial F_1}{\partial y} \right)
\]
Input: Vector field. Output: \textbf{Vector quantity}.

\subsubsection{Physical Significance}\label{physical-significance-1}

Curl measures the rotational tendency or ``spin'' of the vector field at
a point (like a vortex or whirlpool). The direction of the curl vector
is the axis of rotation, given by the right-hand rule.

\subsubsection{Irrotational Vector}\label{irrotational-vector}

If
\[
\nabla \times \vec{F} = \vec{0}
\]
everywhere, the vector field is called \textbf{irrotational}. This means
the field has absolutely no rotational flow or twisting motion.

\subsection{Properties of Position
Vector}\label{properties-of-position-vector}

Let the position vector be
\[
\vec{r} = x\hat{i} + y\hat{j} + z\hat{k}.
\]
Then
\[
r = \lvert \vec{r} \rvert = \sqrt{x^2 + y^2 + z^2}.
\]
Two absolute must-know facts:

\begin{enumerate}
\def\labelenumi{\arabic{enumi}.}
\tightlist
\item
  \textbf{Divergence of \(\vec{r}\) is 3}:
\end{enumerate}
\[
\nabla \cdot \vec{r} = \frac{\partial x}{\partial x} + \frac{\partial y}{\partial y} + \frac{\partial z}{\partial z} = 1 + 1 + 1 = 3
\]
\begin{enumerate}
\def\labelenumi{\arabic{enumi}.}
\setcounter{enumi}{1}
\tightlist
\item
  \textbf{Curl of \(\vec{r}\) is zero}:
\end{enumerate}
\[
\nabla \times \vec{r} =
\begin{vmatrix}
\hat{i} & \hat{j} & \hat{k} \\
\frac{\partial}{\partial x} & \frac{\partial}{\partial y} & \frac{\partial}{\partial z} \\
x & y & z
\end{vmatrix} = \hat{i}(0) - \hat{j}(0) + \hat{k}(0) = \vec{0}
\]

\subsection{Worked Example 1: Solenoidal
Condition}\label{worked-example-1-solenoidal-condition}

\textbf{Problem}: Determine \(a\) such that the vector
\[
\vec{F} = (x + 3y)\hat{i} + (y - 2z)\hat{j} + (x - az)\hat{k}
\]
is solenoidal.

\textbf{Solution}:

For \(\vec{F}\) to be solenoidal,
\[
\nabla \cdot \vec{F} = 0.
\]
\[
\frac{\partial}{\partial x}(x + 3y) + \frac{\partial}{\partial y}(y - 2z) + \frac{\partial}{\partial z}(x - az) = 0
\]
\[
1 + 1 - a = 0 \implies a = 2
\]

\subsection{Worked Example 2: Evaluating
Divergence}\label{worked-example-2-evaluating-divergence}

\textbf{Problem}: Prove that
\[
\text{div}(r^n \vec{r}) = (n + 3)r^n
\]
. For what value of \(n\) is it solenoidal?

\textbf{Solution}:

We know
\[
\frac{\partial r}{\partial x} = \frac{x}{r}
\]
, and similarly for \(y\) and \(z\).
\[
\nabla \cdot (r^n x\hat{i} + r^n y\hat{j} + r^n z\hat{k}) = \frac{\partial}{\partial x}(r^n x) + \frac{\partial}{\partial y}(r^n y) + \frac{\partial}{\partial z}(r^n z)
\]
Apply the product rule to the \(x\) term:
\[
\frac{\partial}{\partial x}(r^n x) = r^n(1) + x\left(nr^{n-1}\frac{\partial r}{\partial x}\right) = r^n + n x r^{n-1}\left(\frac{x}{r}\right) = r^n + n x^2 r^{n-2}
\]
Summing the terms for \(x, y, z\):
\[
\text{div}(r^n \vec{r}) = 3r^n + n r^{n-2}(x^2 + y^2 + z^2)
\]
Since
\[
x^2 + y^2 + z^2 = r^2:
\]
\[
\text{div}(r^n \vec{r}) = 3r^n + n r^{n-2}(r^2) = 3r^n + n r^n = (n + 3)r^n
\]
For the vector to be solenoidal, the divergence must be zero:
\[
(n + 3)r^n = 0 \implies n = -3
\]

\section{A11: Irrotational and Conservative
Fields}\label{a11-irrotational-and-conservative-fields}

\begin{quote}
\textbf{Exam Weight}: Tier 1 \textbar{} \textbf{Appeared in}: 2018
(twice), 2019, 2020, 2021, CT \textbar{} \textbf{Typical Marks}: 3--9
\end{quote}

This is one of the most frequently tested topics in the entire course.
You must know how to prove a field is irrotational (curl = 0) and how to
find its scalar potential.

\subsection{Conservative Force Fields}\label{conservative-force-fields}

\subsubsection{Definition}\label{definition-4}

A force field \(\vec{F}\) is \textbf{conservative} if the work done in
moving a particle between two points is independent of the path taken.

Mathematical equivalent: A vector field \(\vec{F}\) is conservative (or
irrotational) if and only if its curl is zero:
\[
\nabla \times \vec{F} = \vec{0}
\]
\subsubsection{Scalar Potential}\label{scalar-potential}

If \(\vec{F}\) is a conservative field, it can be written as the
gradient of a scalar function \(\phi\):
\[
\vec{F} = \nabla\phi
\]
This function \(\phi(x, y, z)\) is called the \textbf{scalar potential}
of \(\vec{F}\).

*(Note:
\[
\nabla \times (\nabla\phi) = \vec{0}
\]
is a fundamental vector identity. The curl of a gradient is always
zero).*

\subsection{The Method: Finding Scalar
Potential}\label{the-method-finding-scalar-potential}

To find \(\phi\) such that
\[
\vec{F} = \nabla\phi
\]
, where
\[
\vec{F} = F_1\hat{i} + F_2\hat{j} + F_3\hat{k}::
\]
\begin{enumerate}
\def\labelenumi{\arabic{enumi}.}
\tightlist
\item
  Set up three equations:
\end{enumerate}
\[
\frac{\partial\phi}{\partial x} = F_1 \quad \text{(1)}
\]
\[
\frac{\partial\phi}{\partial y} = F_2 \quad \text{(2)}
\]
\[
\frac{\partial\phi}{\partial z} = F_3 \quad \text{(3)}
\]
\begin{enumerate}
\def\labelenumi{\arabic{enumi}.}
\setcounter{enumi}{1}
\tightlist
\item
  Integrate (1) with respect to \(x\). Add a function \(g(y, z)\)
  instead of a constant \(C\).
\item
  Differentiate your result with respect to \(y\), and equate it to (2)
  to find
\end{enumerate}
\[
\frac{\partial g}{\partial y}.
\]
\begin{enumerate}
\def\labelenumi{\arabic{enumi}.}
\setcounter{enumi}{3}
\tightlist
\item
  Integrate to find \(g(y, z)\), adding a function \(h(z)\).
\item
  Substitute \(g(y, z)\) back into your \(\phi\) equation.
\item
  Differentiate with respect to \(z\), equate to (3) to find \(h'(z)\).
\item
  Integrate to find \(h(z)\), adding a final constant \(C\).
\end{enumerate}

\subsection{Must-Know Problem 1 (PYQ
2020)}\label{must-know-problem-1-pyq-2020}

\textbf{Problem}: Show that
\[
\vec{A} = (6xy + z^3)\hat{i} + (3x^2 - z)\hat{j} + (3xz^2 - y)\hat{k}
\]
is irrotational. Find \(\phi\) such that
\[
\vec{A} = \nabla\phi.
\]
\textbf{Solution}:

\textbf{Step 1}: Show Irrotational.
\[
\nabla \times \vec{A} =
\begin{vmatrix}
\hat{i} & \hat{j} & \hat{k} \\
\frac{\partial}{\partial x} & \frac{\partial}{\partial y} & \frac{\partial}{\partial z} \\
6xy + z^3 & 3x^2 - z & 3xz^2 - y
\end{vmatrix}
\]
\begin{itemize}
\tightlist
\item
  \(\hat{i}\)-component:
\end{itemize}
\[
\frac{\partial}{\partial y}(3xz^2 - y) - \frac{\partial}{\partial z}(3x^2 - z) = -1 - (-1) = 0
\]
\begin{itemize}
\tightlist
\item
  \(\hat{j}\)-component:
\end{itemize}
\[
\frac{\partial}{\partial z}(6xy + z^3) - \frac{\partial}{\partial x}(3xz^2 - y) = 3z^2 - 3z^2 = 0
\]
\begin{itemize}
\tightlist
\item
  \(\hat{k}\)-component:
\end{itemize}
\[
\frac{\partial}{\partial x}(3x^2 - z) - \frac{\partial}{\partial y}(6xy + z^3) = 6x - 6x = 0
\]
Curl is zero. It is irrotational.

\textbf{Step 2}: Find Potential.
\[
\frac{\partial\phi}{\partial x} = 6xy + z^3 \implies \phi = 3x^2y + xz^3 + g(y, z) \quad \text{(Eq A)}
\]
Differentiate (Eq A) w.r.t \(y\):
\[
\frac{\partial\phi}{\partial y} = 3x^2 + \frac{\partial g}{\partial y} = 3x^2 - z \implies \frac{\partial g}{\partial y} = -z
\]
Integrate w.r.t \(y\): \(g(y, z) = -yz + h(z)\). Substitute back into
(Eq A):
\[
\phi = 3x^2y + xz^3 - yz + h(z) \quad \text{(Eq B)}
\]
Differentiate (Eq B) w.r.t \(z\):
\[
\frac{\partial\phi}{\partial z} = 3xz^2 - y + h'(z) = 3xz^2 - y \implies h'(z) = 0 \implies h(z) = C
\]
Final Potential:
\[
\phi = 3x^2y + xz^3 - yz + C
\]

\subsection{Must-Know Problem 2 (PYQ
2021)}\label{must-know-problem-2-pyq-2021}

\textbf{Problem}: Show that
\[
\vec{F} = (y^2 + 2xz^2)\hat{i} + (2xy - z)\hat{j} + (2x^2z - y + 2z)\hat{k}
\]
is irrotational and find scalar potential.

\textbf{Solution}:

\textbf{Step 1}: Show Irrotational.
\[
\nabla \times \vec{F} =
\begin{vmatrix}
\hat{i} & \hat{j} & \hat{k} \\
\frac{\partial}{\partial x} & \frac{\partial}{\partial y} & \frac{\partial}{\partial z} \\
y^2 + 2xz^2 & 2xy - z & 2x^2z - y + 2z
\end{vmatrix}
\]
\begin{itemize}
\tightlist
\item
  \(\hat{i}\): \((-1) - (-1) = 0\)
\item
  \(\hat{j}\): \((4xz) - (4xz) = 0\)
\item
  \(\hat{k}\): \((2y) - (2y) = 0\)
\end{itemize}

\textbf{Step 2}: Find Potential.
\[
\frac{\partial\phi}{\partial x} = y^2 + 2xz^2 \implies \phi = xy^2 + x^2z^2 + g(y, z)
\]
\[
\frac{\partial\phi}{\partial y} = 2xy + \frac{\partial g}{\partial y} = 2xy - z \implies \frac{\partial g}{\partial y} = -z \implies g(y, z) = -yz + h(z)
\]
\[
\phi = xy^2 + x^2z^2 - yz + h(z)
\]
\[
\frac{\partial\phi}{\partial z} = 2x^2z - y + h'(z) = 2x^2z - y + 2z \implies h'(z) = 2z \implies h(z) = z^2 + C
\]
Final Potential:
\[
\phi = xy^2 + x^2z^2 - yz + z^2 + C
\]

\subsection{Must-Know Problem 3: Finding Unknown Constants (PYQ
2018)}\label{must-know-problem-3-finding-unknown-constants-pyq-2018}

\textbf{Problem}: Find constants \(a, b, c\) so that
\[
\vec{V} = (x + 2y + az)\hat{i} + (bx - 3y - z)\hat{j} + (4x + cy + 2z)\hat{k}
\]
is irrotational.

\textbf{Solution}:

Set
\[
\nabla \times \vec{V} = 0:
\]
\[
\begin{vmatrix}
\hat{i} & \hat{j} & \hat{k} \\
\frac{\partial}{\partial x} & \frac{\partial}{\partial y} & \frac{\partial}{\partial z} \\
x + 2y + az & bx - 3y - z & 4x + cy + 2z
\end{vmatrix} = 0
\]
\begin{itemize}
\tightlist
\item
  \(\hat{i}\)-component:
\end{itemize}
\[
\frac{\partial}{\partial y}(4x + cy + 2z) - \frac{\partial}{\partial z}(bx - 3y - z) = c - (-1) = c + 1 = 0 \implies c = -1
\]
\begin{itemize}
\tightlist
\item
  \(\hat{j}\)-component:
\end{itemize}
\[
\frac{\partial}{\partial z}(x + 2y + az) - \frac{\partial}{\partial x}(4x + cy + 2z) = a - 4 = 0 \implies a = 4
\]
\begin{itemize}
\tightlist
\item
  \(\hat{k}\)-component:
\end{itemize}
\[
\frac{\partial}{\partial x}(bx - 3y - z) - \frac{\partial}{\partial y}(x + 2y + az) = b - 2 = 0 \implies b = 2
\]
Constants are \(a = 4, b = 2, c = -1\).

\subsection{Exam Patterns}\label{exam-patterns-9}

\begin{itemize}
\tightlist
\item
  ``Show \(\vec{F}\) is conservative and find scalar potential'' is a
  6-mark staple.
\item
  ``Find \(a, b, c\) such that \(\vec{F}\) is irrotational'' is a 3-mark
  easy question.
\item
  Always write the final potential function with the integration
  constant \(+ C\).
\item
  The determinant calculation for curl must be shown explicitly to get
  full marks.
\end{itemize}

\section{A12: Vector Identities and
Proofs}\label{a12-vector-identities-and-proofs}

\begin{quote}
\textbf{Exam Weight}: Tier 2 \textbar{} \textbf{Appeared in}: 2018,
2019, 2021, 2024 \textbar{} \textbf{Typical Marks}: 3--5
\end{quote}

Vector identity proofs test your ability to apply the Del operator
systematically using the chain, product, and quotient rules.

\subsection{The Four Fundamental
Identities}\label{the-four-fundamental-identities}

You should memorize these results, as they often simplify complex
problems:

\begin{enumerate}
\def\labelenumi{\arabic{enumi}.}
\tightlist
\item
  \textbf{Curl of Gradient is Zero}:
\end{enumerate}
\[
\nabla \times (\nabla\phi) = \vec{0}
\]
\begin{enumerate}
\def\labelenumi{\arabic{enumi}.}
\setcounter{enumi}{1}
\tightlist
\item
  \textbf{Divergence of Curl is Zero}:
\end{enumerate}
\[
\nabla \cdot (\nabla \times \vec{A}) = 0
\]
\begin{enumerate}
\def\labelenumi{\arabic{enumi}.}
\setcounter{enumi}{2}
\tightlist
\item
  \textbf{Curl of Curl}:
\end{enumerate}
\[
\nabla \times (\nabla \times \vec{A}) = \nabla(\nabla \cdot \vec{A}) - \nabla^2\vec{A}
\]
\begin{enumerate}
\def\labelenumi{\arabic{enumi}.}
\setcounter{enumi}{3}
\tightlist
\item
  \textbf{Divergence of Gradient is Laplacian}:
\end{enumerate}
\[
\nabla \cdot (\nabla\phi) = \nabla^2\phi
\]

\subsection{Proof 1: (PYQ 2018)}\label{proof-1-pyq-2018}

\textbf{Statement}: If \(\vec{A}\) is a constant vector, then
\[
\nabla(\vec{r} \cdot \vec{A}) = \vec{A}.
\]
\textbf{Proof}:

Let the constant vector be
\[
\vec{A} = A_1\hat{i} + A_2\hat{j} + A_3\hat{k}.
\]
Let the position vector be
\[
\vec{r} = x\hat{i} + y\hat{j} + z\hat{k}.
\]
Calculate the dot product:
\[
\vec{r} \cdot \vec{A} = A_1x + A_2y + A_3z
\]
Now calculate the gradient of this scalar:
\[
\nabla(\vec{r} \cdot \vec{A}) = \left(\hat{i}\frac{\partial}{\partial x} + \hat{j}\frac{\partial}{\partial y} + \hat{k}\frac{\partial}{\partial z}\right)(A_1x + A_2y + A_3z)
\]
\[
\nabla(\vec{r} \cdot \vec{A}) = \hat{i}\frac{\partial}{\partial x}(A_1x + A_2y + A_3z) + \hat{j}\frac{\partial}{\partial y}(A_1x + A_2y + A_3z) + \hat{k}\frac{\partial}{\partial z}(A_1x + A_2y + A_3z)
\]
\[
\nabla(\vec{r} \cdot \vec{A}) = A_1\hat{i} + A_2\hat{j} + A_3\hat{k} = \vec{A} \qquad \blacksquare
\]

\subsection{Proof 2: Curl of Curl Identity (PYQ 2018,
2021)}\label{proof-2-curl-of-curl-identity-pyq-2018-2021}

\textbf{Statement}:
\[
\nabla \times (\nabla \times \vec{A}) = \nabla(\nabla \cdot \vec{A}) - \nabla^2\vec{A}.
\]
\textbf{Proof}:

Let
\[
\vec{A} = A_x\hat{i} + A_y\hat{j} + A_z\hat{k}.
\]
Let
\[
\vec{V} = \nabla \times \vec{A}
\]
. The components of \(\vec{V}\) are:
\[
V_x = \frac{\partial A_z}{\partial y} - \frac{\partial A_y}{\partial z}, \quad V_y = \frac{\partial A_x}{\partial z} - \frac{\partial A_z}{\partial x}, \quad V_z = \frac{\partial A_y}{\partial x} - \frac{\partial A_x}{\partial y}
\]
We want to find:
\[
\nabla \times \vec{V}
\]
Evaluate its x-component:
\[
(\nabla \times \vec{V})_x = \frac{\partial V_z}{\partial y} - \frac{\partial V_y}{\partial z}
\]
Substitute \(V_z\) and \(V_y\):
\[
(\nabla \times \vec{V})_x = \frac{\partial}{\partial y}\left( \frac{\partial A_y}{\partial x} - \frac{\partial A_x}{\partial y} \right) - \frac{\partial}{\partial z}\left( \frac{\partial A_x}{\partial z} - \frac{\partial A_z}{\partial x} \right)
\]
\[
(\nabla \times \vec{V})_x = \frac{\partial^2 A_y}{\partial y \partial x} - \frac{\partial^2 A_x}{\partial y^2} - \frac{\partial^2 A_x}{\partial z^2} + \frac{\partial^2 A_z}{\partial z \partial x}
\]
Add and subtract
\[
\frac{\partial^2 A_x}{\partial x^2}:
\]
\[
(\nabla \times \vec{V})_x = \left( \frac{\partial^2 A_x}{\partial x^2} + \frac{\partial^2 A_y}{\partial y \partial x} + \frac{\partial^2 A_z}{\partial z \partial x} \right) - \left( \frac{\partial^2 A_x}{\partial x^2} + \frac{\partial^2 A_x}{\partial y^2} + \frac{\partial^2 A_x}{\partial z^2} \right)
\]
Factor out the operators:
\[
(\nabla \times \vec{V})_x = \frac{\partial}{\partial x}\left( \frac{\partial A_x}{\partial x} + \frac{\partial A_y}{\partial y} + \frac{\partial A_z}{\partial z} \right) - \left( \frac{\partial^2}{\partial x^2} + \frac{\partial^2}{\partial y^2} + \frac{\partial^2}{\partial z^2} \right)A_x
\]
\[
(\nabla \times \vec{V})_x = \frac{\partial}{\partial x}(\nabla \cdot \vec{A}) - \nabla^2 A_x
\]
By symmetry for y and z components, assembling the vector gives:
\[
\nabla \times (\nabla \times \vec{A}) = \nabla(\nabla \cdot \vec{A}) - \nabla^2\vec{A} \qquad \blacksquare
\]

\subsection{Proof 3: (PYQ 2019)}\label{proof-3-pyq-2019}

\textbf{Statement}: Prove that the Laplacian of \(1/r\) is zero (except
at origin).

\textbf{Proof}:

Let
\[
r = \sqrt{x^2 + y^2 + z^2}
\]
. We know
\[
\frac{\partial r}{\partial x} = \frac{x}{r}
\]
, etc.

First, find the gradient of \(r^{-1}\):
\[
\frac{\partial}{\partial x}(r^{-1}) = -r^{-2}\frac{\partial r}{\partial x} = -r^{-2}\left(\frac{x}{r}\right) = -\frac{x}{r^3}
\]
So,
\[
\nabla(r^{-1}) = -\frac{x\hat{i} + y\hat{j} + z\hat{k}}{r^3} = -\frac{\vec{r}}{r^3}.
\]
Now find the divergence of this gradient (which is \(\nabla^2\)):
\[
\nabla^2(r^{-1}) = \nabla \cdot \left(-\frac{\vec{r}}{r^3}\right) = -\left[ \frac{\partial}{\partial x}\left(\frac{x}{r^3}\right) + \frac{\partial}{\partial y}\left(\frac{y}{r^3}\right) + \frac{\partial}{\partial z}\left(\frac{z}{r^3}\right) \right]
\]
Using the quotient rule on the x-component:
\[
\frac{\partial}{\partial x}\left(\frac{x}{r^3}\right) = \frac{r^3(1) - x(3r^2\frac{\partial r}{\partial x})}{r^6} = \frac{r^3 - 3x^2r}{r^6} = \frac{1}{r^3} - \frac{3x^2}{r^5}
\]
Summing the x, y, and z terms:
\[
\nabla^2(r^{-1}) = -\left[ \left(\frac{1}{r^3} - \frac{3x^2}{r^5}\right) + \left(\frac{1}{r^3} - \frac{3y^2}{r^5}\right) + \left(\frac{1}{r^3} - \frac{3z^2}{r^5}\right) \right]
\]
\[
\nabla^2(r^{-1}) = -\left[ \frac{3}{r^3} - \frac{3(x^2 + y^2 + z^2)}{r^5} \right]
\]
Since
\[
x^2 + y^2 + z^2 = r^2:
\]
\[
\nabla^2(r^{-1}) = -\left[ \frac{3}{r^3} - \frac{3r^2}{r^5} \right] = -\left[ \frac{3}{r^3} - \frac{3}{r^3} \right] = 0 \qquad \blacksquare
\]

\subsection{Exam Patterns}\label{exam-patterns-10}

\begin{itemize}
\tightlist
\item
  The identity proof for
\end{itemize}
\[
\nabla \times (\nabla \times \vec{A})
\]
\subsection{is tedious but important (appeared in 2018 and 2021). Write
out the x-component clearly, then use
symmetry.}\label{is-tedious-but-important-appeared-in-2018-and-2021.-write-out-the-x-component-clearly-then-use-symmetry.}
\[
\nabla^2(1/r) = 0
\]
is a classic vector calculus result. The trick is using
\[
\frac{\partial r}{\partial x} = \frac{x}{r}
\]
correctly with the quotient rule.

\section{A13: Line Integrals and Work
Done}\label{a13-line-integrals-and-work-done}

\begin{quote}
\textbf{Exam Weight}: Tier 1-2 \textbar{} \textbf{Appeared in}: 2017,
2018, 2020, 2021, 2024, CT \textbar{} \textbf{Typical Marks}: 4--6
\end{quote}

A line integral evaluates a function along a curve. The most common
application is calculating the work done by a force field along a path.

\subsection{Formula for Work Done}\label{formula-for-work-done}

The work done by a force field \(\vec{F}\) in moving a particle along a
curve \(C\) is:
\[
W = \int_C \vec{F} \cdot d\vec{r}
\]
If
\[
\vec{F} = F_1\hat{i} + F_2\hat{j} + F_3\hat{k}
\]
and
\[
d\vec{r} = dx\hat{i} + dy\hat{j} + dz\hat{k}::
\]
\[
W = \int_C (F_1 dx + F_2 dy + F_3 dz)
\]
\subsubsection{How to Evaluate}\label{how-to-evaluate}

Convert everything to a single variable (a parameter like \(t\), or
\(x\) if the path is \(y = f(x)\)). Then integrate normally.

\subsection{Worked Example 1: Parametric Path (PYQ
2024)}\label{worked-example-1-parametric-path-pyq-2024}

\textbf{Problem}: Find the work done in a force field
\[
\vec{F} = 3xy\hat{i} - 5z\hat{j} + 10x\hat{k}
\]
along the curve
\[
x = t^2 + 1, y = 2t^2, z = t^3
\]
from \(t = 1\) to \(t = 2\).

\textbf{Solution}:

\textbf{Step 1}: Find differentials.
\[
x = t^2 + 1 \implies dx = 2t dt
\]
\[
y = 2t^2 \implies dy = 4t dt
\]
\[
z = t^3 \implies dz = 3t^2 dt
\]
\textbf{Step 2}: Substitute into the dot product
\[
\vec{F} \cdot d\vec{r} = 3xy dx - 5z dy + 10x dz.
\]
\begin{itemize}
\tightlist
\item
\end{itemize}
\[
3xy dx = 3(t^2 + 1)(2t^2)(2t dt) = (12t^5 + 12t^3) dt
\]
\begin{itemize}
\tightlist
\item
\end{itemize}
\[
-5z dy = -5(t^3)(4t dt) = -20t^4 dt
\]
\begin{itemize}
\tightlist
\item
\end{itemize}
\[
10x dz = 10(t^2 + 1)(3t^2 dt) = (30t^4 + 30t^2) dt
\]
\textbf{Step 3}: Sum and integrate.
\[
W = \int_1^2 (12t^5 + 10t^4 + 12t^3 + 30t^2) dt
\]
\[
W = \left[ 2t^6 + 2t^5 + 3t^4 + 10t^3 \right]_1^2
\]
Upper limit (\(t=2\)):
\(2(64) + 2(32) + 3(16) + 10(8) = 128 + 64 + 48 + 80 = 320\) Lower limit
(\(t=1\)): \(2 + 2 + 3 + 10 = 17\)
\[
W = 320 - 17 = 303
\]

\subsection{Worked Example 2: Planar Path (PYQ
2017)}\label{worked-example-2-planar-path-pyq-2017}

\textbf{Problem}: Evaluate
\[
\int_C (3xy\hat{i} - y^2\hat{j}) \cdot d\vec{r}
\]
when \(C\) is
\[
y = 2x^2
\]
from \((0,0)\) to \((1,2)\).

\textbf{Solution}:

We have
\[
\vec{F} \cdot d\vec{r} = 3xy dx - y^2 dy.
\]
Path:
\[
y = 2x^2 \implies dy = 4x dx.
\]
Limits for \(x\): 0 to 1.

Substitute \(y\) and \(dy\):
\[
\int_0^1 \left[ 3x(2x^2)dx - (2x^2)^2(4x)dx \right] = \int_0^1 (6x^3 - 16x^5) dx
\]
\[
= \left[ \frac{3}{2}x^4 - \frac{8}{3}x^6 \right]_0^1 = \frac{3}{2} - \frac{8}{3} = \frac{9 - 16}{6} = -\frac{7}{6}
\]

\subsection{Worked Example 3: Segmented Path (PYQ
2021)}\label{worked-example-3-segmented-path-pyq-2021}

\textbf{Problem}: Evaluate
\[
\int_C \vec{A} \cdot d\vec{r}
\]
where
\[
\vec{A} = (3x^2 + 6y)\hat{i} - 14yz\hat{j} + 20xz^2\hat{k}
\]
along straight lines from \((0,0,0)\) to \((1,0,0)\) to \((1,1,0)\) to
\((1,1,1)\).

\textbf{Solution}:

Break into three paths: \(C_1, C_2, C_3\).
\[
\vec{A} \cdot d\vec{r} = (3x^2 + 6y)dx - 14yz dy + 20xz^2 dz.
\]
\textbf{Path \(C_1\)}: \((0,0,0)\) to \((1,0,0)\).
\(y=0, z=0 \implies dy=0, dz=0\). \(x\) goes 0 to 1.
\[
\int_{C_1} = \int_0^1 (3x^2 + 0) dx = [x^3]_0^1 = 1
\]
\textbf{Path \(C_2\)}: \((1,0,0)\) to \((1,1,0)\).
\(x=1, z=0 \implies dx=0, dz=0\). \(y\) goes 0 to 1.
\[
\int_{C_2} = \int_0^1 -14y(0) dy = 0
\]
\textbf{Path \(C_3\)}: \((1,1,0)\) to \((1,1,1)\).
\(x=1, y=1 \implies dx=0, dy=0\). \(z\) goes 0 to 1.
\[
\int_{C_3} = \int_0^1 20(1)z^2 dz = \left[\frac{20}{3}z^3\right]_0^1 = \frac{20}{3}
\]
Total Integral =
\[
1 + 0 + \frac{20}{3} = \frac{23}{3}.
\]

\subsection{Exam Patterns}\label{exam-patterns-11}

\begin{itemize}
\tightlist
\item
  Parametric path questions (like Example 1) are very common. Substitute
  \(x, y, z\) and \(dx, dy, dz\) perfectly. Watch your algebra.
\item
  Segmented straight line paths (Example 3) are simple but test your
  logic. On each segment, two variables are constant (so their
  differentials are zero).
\item
  Always explicitly write the dot product
\end{itemize}
\[
\vec{F} \cdot d\vec{r} = F_1 dx + F_2 dy + F_3 dz
\]
first.

\section{A14: Surface Integrals}\label{a14-surface-integrals}

\begin{quote}
\textbf{Exam Weight}: Tier 3 \textbar{} \textbf{Appeared in}: 2017, 2020
\textbar{} \textbf{Typical Marks}: 6
\end{quote}

A surface integral calculates the total flux of a vector field passing
through a surface. It is the 2D equivalent of a line integral. Most
surface integrals in exams are evaluated using the Divergence Theorem or
Stokes' Theorem, but occasionally you must calculate them directly.

\subsection{The Formula}\label{the-formula}

The flux of a vector field \(\vec{A}\) through a surface \(S\) is:
\[
\text{Flux} = \iint_S \vec{A} \cdot \hat{n} dS
\]
where \(\hat{n}\) is the unit outward normal vector to the surface.

\subsubsection{Evaluation by Projection}\label{evaluation-by-projection}

To evaluate this directly, we project the surface \(S\) onto one of the
coordinate planes (e.g., the \(xy\)-plane).

Let \(R\) be the projected region in the \(xy\)-plane. The differential
area \(dS\) on the surface relates to the area \(dx dy\) on the plane
by:
\[
dS = \frac{dx dy}{|\hat{n} \cdot \hat{k}|}
\]
If projecting onto the \(yz\)-plane:
\[
dS = \frac{dy dz}{\lvert \hat{n} \cdot \hat{i} \rvert}.
\]
If projecting onto the \(zx\)-plane:
\[
dS = \frac{dz dx}{\lvert \hat{n} \cdot \hat{j} \rvert}.
\]

\subsection{Must-Know Problem: Cylinder Surface (PYQ 2017,
2020)}\label{must-know-problem-cylinder-surface-pyq-2017-2020}

\textbf{Problem}: Evaluate
\[
\iint_S \vec{A} \cdot \hat{n} dS
\]
where
\[
\vec{A} = z\hat{i} + x\hat{j} - 3y^2z\hat{k}
\]
and \(S\) is the surface of the cylinder
\[
x^2 + y^2 = 16
\]
included in the first octant between \(z=0\) and \(z=5\).

\textbf{Solution}:

\textbf{Step 1}: Find the normal vector \(\hat{n}\). The surface is
\[
g(x,y,z) = x^2 + y^2 - 16 = 0.
\]
Gradient:
\[
\nabla g = 2x\hat{i} + 2y\hat{j}.
\]
Magnitude:
\[
\lvert \nabla g \rvert = \sqrt{4x^2 + 4y^2} = 2\sqrt{16} = 8.
\]
Unit normal:
\[
\hat{n} = \frac{2x\hat{i} + 2y\hat{j}}{8} = \frac{x\hat{i} + y\hat{j}}{4}
\]
\textbf{Step 2}: Calculate \(\vec{A} \cdot \hat{n}\).
\[
\vec{A} \cdot \hat{n} = (z\hat{i} + x\hat{j} - 3y^2z\hat{k}) \cdot \left( \frac{x\hat{i} + y\hat{j}}{4} \right) = \frac{xz + xy}{4}
\]
\textbf{Step 3}: Choose a projection plane. We project the cylinder onto
the \(yz\)-plane. The region \(R\) is bounded by \(y \in [0, 4]\) and
\(z \in [0, 5]\). The area element is:
\[
dS = \frac{dy dz}{|\hat{n} \cdot \hat{i}|} = \frac{dy dz}{|x/4|} = \frac{4}{x} dy dz
\]
\textbf{Step 4}: Set up and evaluate the integral.
\[
\iint_S (\vec{A} \cdot \hat{n}) dS = \iint_R \left( \frac{xz + xy}{4} \right) \left( \frac{4}{x} dy dz \right)
\]
The \(x\) and \(4\) cancel perfectly:
\[
\iint_R (z + y) dy dz
\]
Integrate over the rectangular region \(R\):
\[
\int_0^4 \int_0^5 (y + z) dz dy
\]
Integrate with respect to \(z\):
\[
\int_0^4 \left[ yz + \frac{z^2}{2} \right]_0^5 dy = \int_0^4 \left( 5y + \frac{25}{2} \right) dy
\]
Integrate with respect to \(y\):
\[
\left[ \frac{5y^2}{2} + \frac{25}{2}y \right]_0^4 = \frac{5(16)}{2} + \frac{25(4)}{2} = 40 + 50 = 90
\]
The surface integral value is \(90\).

\subsection{Exam Patterns}\label{exam-patterns-12}

\begin{itemize}
\tightlist
\item
  Direct evaluation of surface integrals is rare. This single cylinder
  problem appeared identically in 2017 and 2020.
\item
  The projection method is crucial. If the surface equation contains
  \(x\) and \(y\) but no \(z\), projecting onto the \(yz\) or \(zx\)
  plane is usually easier than projecting onto the \(xy\) plane.
\item
  Most surface integrals in exams will ask you to use Gauss's Divergence
  Theorem or Stokes' Theorem instead of direct evaluation.
\end{itemize}

\section{A15: Volume Integrals}\label{a15-volume-integrals}

\begin{quote}
\textbf{Exam Weight}: Tier 2-3 \textbar{} \textbf{Appeared in}: 2018,
2019, 2021 \textbar{} \textbf{Typical Marks}: 4--6
\end{quote}

Volume integrals (triple integrals) are used to sum a quantity over a 3D
region. You must carefully determine the limits of integration for
\(x\), \(y\), and \(z\).

\subsection{The Formula}\label{the-formula-1}

The volume integral of a scalar function \(\phi(x,y,z)\) over a volume
\(V\) is:
\[
\iiint_V \phi \, dV = \iiint_V \phi \, dx \, dy \, dz
\]
For a vector field
\[
\vec{F} = F_1\hat{i} + F_2\hat{j} + F_3\hat{k},
\]
You integrate each component separately:
\[
\iiint_V \vec{F} \, dV = \left(\iiint_V F_1 \, dV\right)\hat{i} + \left(\iiint_V F_2 \, dV\right)\hat{j} + \left(\iiint_V F_3 \, dV\right)\hat{k}
\]

\subsection{Setting Up Integration
Limits}\label{setting-up-integration-limits}

The hardest part is finding the limits. For a region bounded by planes:
1. \textbf{Inner integral (\(z\))}: Express \(z\) in terms of \(x\) and
\(y\) using the top and bottom bounding surfaces. 2. \textbf{Middle
integral (\(y\))}: Set \(z=0\) (or analyze the projection on the
\(xy\)-plane) to find \(y\) in terms of \(x\). 3. \textbf{Outer integral
(\(x\))}: Find the absolute minimum and maximum constant values for
\(x\).

\subsection{Worked Example 1: Scalar Volume Integral (PYQ
2018)}\label{worked-example-1-scalar-volume-integral-pyq-2018}

\textbf{Problem}: Let
\[
\phi = 45x^2y
\]
. Evaluate
\[
\iiint_V \phi \, dV
\]
where \(V\) is the closed region bounded by planes
\(4x + 2y + z = 8, x = 0, y = 0, z = 0\).

\textbf{Solution}:

\textbf{Step 1}: Find limits. - \(z\) goes from \(0\) to
\(8 - 4x - 2y\). - In the \(xy\)-plane (\(z=0\)), the boundary is
\(4x + 2y = 8 \implies y = 4 - 2x\). So \(y\) goes from \(0\) to
\(4 - 2x\). - Set \(y=0\) to find \(x\): \(4x = 8 \implies x = 2\). So
\(x\) goes from \(0\) to \(2\).

\textbf{Step 2}: Set up integral.
\[
\int_0^2 \int_0^{4-2x} \int_0^{8-4x-2y} 45x^2y \, dz \, dy \, dx
\]
\textbf{Step 3}: Integrate \(z\).
\[
45 \int_0^2 x^2 \int_0^{4-2x} y [z]_0^{8-4x-2y} dy dx = 45 \int_0^2 x^2 \int_0^{4-2x} (8y - 4xy - 2y^2) dy dx
\]
\textbf{Step 4}: Integrate \(y\).
\[
45 \int_0^2 x^2 \left[ 4y^2 - 2xy^2 - \frac{2}{3}y^3 \right]_0^{4-2x} dx
\]
Notice
\[
4y^2 - 2xy^2 = y^2(4 - 2x).
\]
Substitute upper limit \(y = 4-2x\):
\[
45 \int_0^2 x^2 \left[ (4-2x)^2(4-2x) - \frac{2}{3}(4-2x)^3 \right] dx = 45 \int_0^2 x^2 \left[ \frac{1}{3}(4-2x)^3 \right] dx
\]
\textbf{Step 5}: Integrate \(x\).
\[
15 \int_0^2 x^2 (4-2x)^3 dx = 120 \int_0^2 x^2 (2-x)^3 dx
\]
Expand
\[
(2-x)^3 = 8 - 12x + 6x^2 - x^3:
\]
\[
120 \int_0^2 (8x^2 - 12x^3 + 6x^4 - x^5) dx
\]
\[
120 \left[ \frac{8}{3}x^3 - 3x^4 + \frac{6}{5}x^5 - \frac{x^6}{6} \right]_0^2 = 120 \left( \frac{64}{3} - 48 + \frac{192}{5} - \frac{64}{6} \right) = 128
\]

\subsection{Worked Example 2: Integral of Divergence (PYQ
2019)}\label{worked-example-2-integral-of-divergence-pyq-2019}

\textbf{Problem}: Evaluate
\[
\iiint_V \nabla \cdot \vec{F} \, dV
\]
where
\[
\vec{F} = (2x^2 - 3z)\hat{i} - 2xy\hat{j} - 4x\hat{k}
\]
and \(V\) is bounded by \(x=0, y=0, z=0, 2x+2y+z=4\).

\textbf{Solution}:

\textbf{Step 1}: Find divergence.
\[
\nabla \cdot \vec{F} = \frac{\partial}{\partial x}(2x^2 - 3z) + \frac{\partial}{\partial y}(-2xy) + \frac{\partial}{\partial z}(-4x) = 4x - 2x + 0 = 2x
\]
\textbf{Step 2}: Find limits. \(z\) from \(0\) to \(4 - 2x - 2y\). \(y\)
from \(0\) to \(2 - x\). \(x\) from \(0\) to \(2\).

\textbf{Step 3}: Integrate.
\[
\int_0^2 \int_0^{2-x} \int_0^{4-2x-2y} 2x \, dz \, dy \, dx = \int_0^2 \int_0^{2-x} 2x(4 - 2x - 2y) dy dx
\]
\[
= \int_0^2 2x \left[ (4-2x)y - y^2 \right]_0^{2-x} dx = \int_0^2 2x \left[ 2(2-x)(2-x) - (2-x)^2 \right] dx
\]
\[
= \int_0^2 2x (2-x)^2 dx = \int_0^2 (8x - 8x^2 + 2x^3) dx = \left[ 4x^2 - \frac{8}{3}x^3 + \frac{1}{2}x^4 \right]_0^2 = \frac{8}{3}
\]

\subsection{Exam Patterns}\label{exam-patterns-13}

\begin{itemize}
\tightlist
\item
  Volume integrals in this course are mostly an exercise in polynomial
  integration and algebra. Make sure your expansions are careful.
\item
  Setting up limits from an intersecting plane like \(ax + by + cz = d\)
  is a recurring pattern. Always work from \(z \to y \to x\).
\item
  You will also compute volume integrals as the right-hand side of the
  Gauss Divergence Theorem.
\end{itemize}

\section{A16: Green's Theorem}\label{a16-greens-theorem}

\begin{quote}
\textbf{Exam Weight}: Tier 1 \textbar{} \textbf{Appeared in}: 2017
(twice), 2019, 2020, 2021, 2023, 2024 \textbar{} \textbf{Typical Marks}:
5--6
\end{quote}

Green's theorem relates a line integral around a closed boundary curve
\(C\) to a double integral over the plane region \(R\) bounded by \(C\).
It is one of the most frequently tested topics.

\subsection{Statement of the Theorem}\label{statement-of-the-theorem}

Let \(R\) be a closed, flat region bounded by a simple, closed curve
\(C\) in the \(xy\)-plane. If \(P(x,y)\) and \(Q(x,y)\) are continuous
functions with continuous first-order partial derivatives in \(R\),
then:
\[
\oint_C (P dx + Q dy) = \iint_R \left( \frac{\partial Q}{\partial x} - \frac{\partial P}{\partial y} \right) dx dy
\]
Here the curve \(C\) is traversed in the counterclockwise (positive)
direction.

\subsection{Must-Know Proof (PYQ 2017, 2021,
2024)}\label{must-know-proof-pyq-2017-2021-2024}

You must be able to prove Green's theorem. We split it into two parts:
1.
\[
\oint_C P dx = -\iint_R \frac{\partial P}{\partial y} dx dy
\]
\begin{enumerate}
\def\labelenumi{\arabic{enumi}.}
\setcounter{enumi}{1}
\tightlist
\item
\end{enumerate}
\[
\oint_C Q dy = \iint_R \frac{\partial Q}{\partial x} dx dy
\]
\textbf{Proof of Part 1}: Let region \(R\) be bounded by curves
\[
y = y_1(x)
\]
(lower) and
\[
y = y_2(x)
\]
(upper) from \(x=a\) to \(x=b\).

Evaluate the double integral:
\[
\iint_R \frac{\partial P}{\partial y} dx dy = \int_a^b \left[ \int_{y_1(x)}^{y_2(x)} \frac{\partial P}{\partial y} dy \right] dx = \int_a^b \left[ P(x, y_2(x)) - P(x, y_1(x)) \right] dx \quad \text{--- (Eq 1)}
\]
Now evaluate the line integral \(\oint_C P dx\). The curve \(C\) has two
paths: - \(C_1\): Along \(y_1(x)\) from \(a\) to \(b\). - \(C_2\): Along
\(y_2(x)\) from \(b\) to \(a\).
\[
\oint_C P dx = \int_a^b P(x, y_1(x)) dx + \int_b^a P(x, y_2(x)) dx
\]
Reverse limits on the second integral:
\[
\oint_C P dx = \int_a^b P(x, y_1(x)) dx - \int_a^b P(x, y_2(x)) dx = -\int_a^b \left[ P(x, y_2(x)) - P(x, y_1(x)) \right] dx \quad \text{--- (Eq 2)}
\]
Comparing (Eq 1) and (Eq 2) proves Part 1. By symmetry (bounding \(x\)
with functions of \(y\)), Part 2 is proven. Adding them yields Green's
theorem.

\subsection{Worked Example 1: Evaluation (PYQ
2017)}\label{worked-example-1-evaluation-pyq-2017}

\textbf{Problem}: Evaluate
\[
\oint_C (2x + y^2)dx + (3y - 4x)dy
\]
where \(C\) is the boundary of the region between
\[
y = x^2
\]
and
\[
y^2 = x.
\]
\textbf{Solution}: Identify \(P\) and \(Q\):
\[
P = 2x + y^2 \implies \frac{\partial P}{\partial y} = 2y
\]
\[
Q = 3y - 4x \implies \frac{\partial Q}{\partial x} = -4
\]
Apply Green's theorem:
\[
\iint_R (-4 - 2y) dy dx
\]
Limits: The curves intersect at \((0,0)\) and \((1,1)\). For a given
\(x\), \(y\) goes from the lower curve
\[
y=x^2
\]
to the upper curve \(y=\sqrt{x}\).
\[
\int_0^1 \int_{x^2}^{\sqrt{x}} (-4 - 2y) dy dx = \int_0^1 \left[ -4y - y^2 \right]_{x^2}^{\sqrt{x}} dx
\]
\[
= \int_0^1 \left[ (-4\sqrt{x} - x) - (-4x^2 - x^4) \right] dx = \int_0^1 (-4x^{1/2} - x + 4x^2 + x^4) dx
\]
\[
= \left[ -\frac{8}{3}x^{3/2} - \frac{1}{2}x^2 + \frac{4}{3}x^3 + \frac{1}{5}x^5 \right]_0^1 = -\frac{8}{3} - \frac{1}{2} + \frac{4}{3} + \frac{1}{5} = -\frac{49}{30}
\]

\subsection{Worked Example 2: Verification (PYQ 2020,
2023)}\label{worked-example-2-verification-pyq-2020-2023}

\textbf{Problem}: Verify Green's theorem for
\(\oint_C (xy + y^2)dx + x^2 dy\) where \(C\) is bounded by \(y = x\)
and
\[
y = x^2.
\]
\textbf{Solution}:

\textbf{Step 1: Evaluate via Double Integral (Green's Theorem)}
\[
P = xy + y^2 \implies \frac{\partial P}{\partial y} = x + 2y
\]
\[
Q = x^2 \implies \frac{\partial Q}{\partial x} = 2x
\]
\[
\frac{\partial Q}{\partial x} - \frac{\partial P}{\partial y} = 2x - (x + 2y) = x - 2y
\]
Limits: \(x\) from 0 to 1, \(y\) from \(x^2\) to \(x\).
\[
\int_0^1 \int_{x^2}^x (x - 2y) dy dx = \int_0^1 \left[ xy - y^2 \right]_{x^2}^x dx = \int_0^1 (0 - (x^3 - x^4)) dx
\]
\[
= \int_0^1 (x^4 - x^3) dx = \left[ \frac{x^5}{5} - \frac{x^4}{4} \right]_0^1 = \frac{1}{5} - \frac{1}{4} = -\frac{1}{20}
\]
\textbf{Step 2: Evaluate via Line Integral} Path 1 (
\[
y = x^2, dy = 2xdx
\]
from 0 to 1):
\[
\int_0^1 \left[ (x^3 + x^4)dx + x^2(2x dx) \right] = \int_0^1 (3x^3 + x^4) dx = \left[ \frac{3}{4}x^4 + \frac{x^5}{5} \right]_0^1 = \frac{19}{20}
\]
Path 2 (\(y = x, dy = dx\) from 1 to 0):
\[
\int_1^0 (x^2 + x^2)dx + x^2 dx = \int_1^0 3x^2 dx = [x^3]_1^0 = -1
\]
Total line integral =
\[
\frac{19}{20} - 1 = -\frac{1}{20}.
\]
Since both equal \(-\frac{1}{20}\), the theorem is verified.

\subsection{Exam Patterns}\label{exam-patterns-14}

\begin{itemize}
\tightlist
\item
  ``Verify Green's theorem'' means you must compute BOTH the line
  integral (summing along all paths) and the double integral
  independently to show they match. This takes time, be careful with
  signs.
\item
  The proof of Green's theorem is a high-value question. Memorize the
  algebraic steps for Part 1.
\item
  Sometimes evaluating the double integral just yields the Area of the
  region (e.g., if
\end{itemize}
\[
\frac{\partial Q}{\partial x} - \frac{\partial P}{\partial y} = -1
\]
). Don't overcomplicate it if it's a known shape like a triangle.

\section{A17: Stokes' Theorem}\label{a17-stokes-theorem}

\begin{quote}
\textbf{Exam Weight}: Tier 2 \textbar{} \textbf{Appeared in}: 2023, 2024
\textbar{} \textbf{Typical Marks}: 5
\end{quote}

Stokes' theorem is the 3D extension of Green's theorem. It relates a
line integral around a closed curve in 3D space to a surface integral
over any surface bounded by that curve.

\subsection{Statement of the Theorem}\label{statement-of-the-theorem-1}

Let \(S\) be an open, two-sided surface bounded by a closed,
non-self-intersecting curve \(C\). If \(\vec{A}\) is a vector field with
continuous first-order partial derivatives, then:
\[
\oint_C \vec{A} \cdot d\vec{r} = \iint_S (\nabla \times \vec{A}) \cdot \hat{n} dS
\]
Here the boundary curve \(C\) is traversed in the positive direction
(determined by the right-hand rule with respect to the normal vector
\(\hat{n}\)).

\subsection{Must-Know Proof (PYQ 2023)}\label{must-know-proof-pyq-2023}

You must be able to prove Stokes' theorem by projecting the surface onto
a coordinate plane and using Green's theorem.

Let the surface equation be \(z = f(x, y)\) over a region \(R\) in the
\(xy\)-plane. Let
\[
\vec{A} = A_1\hat{i} + A_2\hat{j} + A_3\hat{k}.
\]
We prove it for \(A_1\) first:
\[
\oint_C A_1 dx = \iint_S [\nabla \times (A_1\hat{i})] \cdot \hat{n} dS.
\]
\textbf{Step 1: The Line Integral} Project the curve \(C\) onto the
\(xy\)-plane as \(C'\).
\[
\oint_C A_1(x, y, z) dx = \oint_{C'} A_1(x, y, f(x, y)) dx
\]
Apply Green's theorem to \(C'\):
\[
= \iint_R -\frac{\partial}{\partial y} [A_1(x, y, f(x, y))] dA
\]
Using the chain rule:
\[
\oint_C A_1 dx = \iint_R \left( -\frac{\partial A_1}{\partial y} - \frac{\partial A_1}{\partial z}\frac{\partial z}{\partial y} \right) dx dy \quad \text{--- (Eq 1)}
\]
\textbf{Step 2: The Surface Integral} Calculate the curl of
\(A_1\hat{i}\):
\[
\nabla \times (A_1\hat{i}) = \frac{\partial A_1}{\partial z}\hat{j} - \frac{\partial A_1}{\partial y}\hat{k}
\]
For surface \(z - f(x, y) = 0\), the normal vector gives:
\[
\hat{n} dS = \left( -\frac{\partial z}{\partial x}\hat{i} - \frac{\partial z}{\partial y}\hat{j} + \hat{k} \right) dx dy
\]
Calculate the dot product:
\[
[\nabla \times (A_1\hat{i})] \cdot \hat{n} dS = \left( -\frac{\partial A_1}{\partial z}\frac{\partial z}{\partial y} - \frac{\partial A_1}{\partial y} \right) dx dy \quad \text{--- (Eq 2)}
\]
Comparing (Eq 1) and (Eq 2) proves the theorem for \(A_1\). The process
is identical for \(A_2\) and \(A_3\). Adding them together yields
Stokes' theorem.

\subsection{Worked Example: Verification (PYQ
2024)}\label{worked-example-verification-pyq-2024}

\textbf{Problem}: Verify Stokes' theorem for
\[
\vec{A} = (y - z + 2)\hat{i} + (yz + 4)\hat{j} - xz\hat{k}
\]
where \(S\) is the surface of the cube \(x=0, y=0, z=0, x=2, y=2, z=2\)
above the \(xy\) plane (i.e., open at the bottom).

\textbf{Solution}:

\textbf{Step 1: The Line Integral} The boundary curve \(C\) is the
square in the \(z=0\) plane from \(x=0\) to \(x=2\) and \(y=0\) to
\(y=2\). On \(z=0, dz=0\), the vector field is
\[
\vec{A} = (y + 2)\hat{i} + 4\hat{j}.
\]
\[
\oint_C \vec{A} \cdot d\vec{r} = \oint_C (y + 2) dx + 4 dy
\]
Path 1 (\(y=0, dy=0, x\) from 0 to 2):
\[
\int_0^2 2 dx = 4
\]
Path 2 (\(x=2, dx=0, y\) from 0 to 2):
\[
\int_0^2 4 dy = 8
\]
Path 3 (\(y=2, dy=0, x\) from 2 to 0):
\[
\int_2^0 4 dx = -8
\]
Path 4 (\(x=0, dx=0, y\) from 2 to 0):
\[
\int_2^0 4 dy = -8
\]
Total Line Integral = \(4 + 8 - 8 - 8 = -4\).

\textbf{Step 2: The Surface Integral} Calculate Curl:
\[
\nabla \times \vec{A} =
\begin{vmatrix}
\hat{i} & \hat{j} & \hat{k} \\
\frac{\partial}{\partial x} & \frac{\partial}{\partial y} & \frac{\partial}{\partial z} \\
y-z+2 & yz+4 & -xz
\end{vmatrix} = -y\hat{i} + (z - 1)\hat{j} - \hat{k}
\]
Evaluate over the 5 faces of the cube above \(z=0\): 1. **Top (
\[
z=2, \hat{n}=\hat{k}
\]
)**:
\[
(\nabla \times \vec{A}) \cdot \hat{n} = -1
\]
. Integral:
\[
\iint (-1) dxdy = -4.
\]
\begin{enumerate}
\def\labelenumi{\arabic{enumi}.}
\setcounter{enumi}{1}
\tightlist
\item
  **Front (
\end{enumerate}
\[
x=2, \hat{n}=\hat{i}
\]
)**:
\[
(\nabla \times \vec{A}) \cdot \hat{n} = -y
\]
. Integral:
\[
\int_0^2\int_0^2 -y dydz = 2\left[-\frac{y^2}{2}\right]_0^2 = -4.
\]
\begin{enumerate}
\def\labelenumi{\arabic{enumi}.}
\setcounter{enumi}{2}
\tightlist
\item
  **Back (
\end{enumerate}
\[
x=0, \hat{n}=-\hat{i}
\]
)**:
\[
(\nabla \times \vec{A}) \cdot \hat{n} = y
\]
. Integral:
\[
\int_0^2\int_0^2 y dydz = 2\left[\frac{y^2}{2}\right]_0^2 = 4.
\]
\begin{enumerate}
\def\labelenumi{\arabic{enumi}.}
\setcounter{enumi}{3}
\tightlist
\item
  **Right (
\end{enumerate}
\[
y=2, \hat{n}=\hat{j}
\]
)**:
\[
(\nabla \times \vec{A}) \cdot \hat{n} = z-1
\]
. Integral:
\[
\int_0^2\int_0^2 (z-1) dxdz = 2\left[\frac{z^2}{2}-z\right]_0^2 = 0.
\]
\begin{enumerate}
\def\labelenumi{\arabic{enumi}.}
\setcounter{enumi}{4}
\tightlist
\item
  **Left (
\end{enumerate}
\[
y=0, \hat{n}=-\hat{j}
\]
)**:
\[
(\nabla \times \vec{A}) \cdot \hat{n} = -(z-1)
\]
. Integral:
\[
\int_0^2\int_0^2 (1-z) dxdz = 2\left[z-\frac{z^2}{2}\right]_0^2 = 0.
\]
Total Surface Integral = \(-4 - 4 + 4 + 0 + 0 = -4\).

Both integrals yield \(-4\). Stokes' theorem is verified.

\subsection{Exam Patterns}\label{exam-patterns-15}

\begin{itemize}
\tightlist
\item
  Verifying Stokes' theorem over a cube involves 5 surface integrals (if
  open at the bottom) or 6 surface integrals. Be extremely meticulous
  with your normal vectors
\end{itemize}
\[
\hat{i}, -\hat{i}, \hat{j}, -\hat{j}, \hat{k}, -\hat{k}
\]
For each face. - When asked to ``evaluate using Stokes' theorem'', you
can choose whichever side of the equation is easier. Usually, converting
a complex line integral into a surface integral via curl is the intended
path.

\section{A18: Gauss Divergence
Theorem}\label{a18-gauss-divergence-theorem}

\begin{quote}
\textbf{Exam Weight}: Tier 1 \textbar{} \textbf{Appeared in}: 2018, 2021
\textbar{} \textbf{Typical Marks}: 6
\end{quote}

The Gauss Divergence Theorem relates the flux of a vector field across a
closed surface to the volume integral of the divergence inside that
surface. It is extremely powerful for converting difficult surface
integrals into simple volume integrals.

\subsection{Statement of the Theorem}\label{statement-of-the-theorem-2}

Let \(V\) be a volume bounded by a closed surface \(S\). If \(\vec{A}\)
is a vector field with continuous first-order partial derivatives, then:
\[
\iint_S \vec{A} \cdot \hat{n} dS = \iiint_V (\nabla \cdot \vec{A}) dV
\]
where \(\hat{n}\) is the unit outward normal vector to the surface
\(S\).

\emph{(Note: The theorem requires a strictly CLOSED surface. If a
surface is open, like a hemisphere without its base, you must add the
base to close it before applying the theorem).}

\subsection{Worked Example: Verification (PYQ 2018,
2021)}\label{worked-example-verification-pyq-2018-2021}

\textbf{Problem}: Verify the divergence theorem for
\[
\vec{A} = 4x\hat{i} - 2y^2\hat{j} + z^2\hat{k}
\]
taken over the region bounded by
\[
x^2 + y^2 = 4, z = 0, z = 3.
\]
\textbf{Solution}:

\textbf{Step 1: Evaluate the Volume Integral} Calculate Divergence:
\[
\nabla \cdot \vec{A} = \frac{\partial}{\partial x}(4x) + \frac{\partial}{\partial y}(-2y^2) + \frac{\partial}{\partial z}(z^2) = 4 - 4y + 2z
\]
Set up cylindrical coordinates
(\(x = r\cos\phi, y = r\sin\phi, z = z, dV = r dr d\phi dz\)): Limits:
\(r \in [0, 2]\), \(\phi \in [0, 2\pi]\), \(z \in [0, 3]\).
\[
\iiint_V (\nabla \cdot \vec{A}) dV = \int_0^3 \int_0^{2\pi} \int_0^2 (4 - 4r\sin\phi + 2z) r \, dr \, d\phi \, dz
\]
Integrate w.r.t \(\phi\) first (the \(\sin\phi\) term evaluates to 0
over \(2\pi\)):
\[
\int_0^{2\pi} (4r - 4r^2\sin\phi + 2rz) d\phi = 2\pi(4r + 2rz)
\]
Integrate w.r.t \(r\):
\[
2\pi \int_0^2 (4r + 2rz) dr = 2\pi [2r^2 + r^2 z]_0^2 = 2\pi(8 + 4z)
\]
Integrate w.r.t \(z\):
\[
2\pi \int_0^3 (8 + 4z) dz = 2\pi [8z + 2z^2]_0^3 = 2\pi(24 + 18) = 84\pi
\]
\textbf{Step 2: Evaluate the Surface Integral} The closed cylinder has 3
surfaces: Top (\(S_1\)), Bottom (\(S_2\)), and Curved Wall (\(S_3\)).

\begin{itemize}
\tightlist
\item
  **Top Surface \(S_1\) (
\end{itemize}
\[
z=3, \hat{n}=\hat{k}
\]
)**:
\[
\vec{A} \cdot \hat{k} = z^2 = 3^2 = 9.
\]
\[
\iint_{S_1} 9 dS = 9 \times (\text{Area of circle}) = 9 \times \pi(2^2) = 36\pi.
\]
\begin{itemize}
\tightlist
\item
  **Bottom Surface \(S_2\) (
\end{itemize}
\[
z=0, \hat{n}=-\hat{k}
\]
)**:
\[
\vec{A} \cdot (-\hat{k}) = -z^2 = -0^2 = 0.
\]
\[
\iint_{S_2} 0 dS = 0.
\]
\begin{itemize}
\tightlist
\item
  **Curved Wall \(S_3\)
\end{itemize}
\[
x^2+y^2=4, \hat{n}=\frac{x\hat{i}+y\hat{j}}{2}
\]
**:
\[
\vec{A} \cdot \hat{n} = (4x\hat{i} - 2y^2\hat{j} + z^2\hat{k}) \cdot \left( \frac{x\hat{i} + y\hat{j}}{2} \right) = 2x^2 - y^3.
\]
In cylindrical (\(x=2\cos\phi, y=2\sin\phi, dS = 2 d\phi dz\)):
\[
\iint_{S_3} (2x^2 - y^3) dS = \int_0^3 \int_0^{2\pi} \left[ 2(4\cos^2\phi) - 8\sin^3\phi \right] 2 d\phi dz
\]
We know
\[
\int_0^{2\pi} \cos^2\phi d\phi = \pi
\]
and
\[
\int_0^{2\pi} \sin^3\phi d\phi = 0.
\]
\[
= \int_0^3 2(8\pi - 0) dz = \int_0^3 16\pi dz = 48\pi
\]
Total Surface Integral = \(36\pi + 0 + 48\pi = 84\pi\).

Both integrals equal \(84\pi\). The divergence theorem is verified.

\subsection{The ``Big Three'' Summary}\label{the-big-three-summary}

Vector Calculus revolves around three major integration theorems.
Understanding when to use which is critical:

\begin{longtable}[]{@{}
  >{\raggedright\arraybackslash}p{(\linewidth - 4\tabcolsep) * \real{0.3333}}
  >{\raggedright\arraybackslash}p{(\linewidth - 4\tabcolsep) * \real{0.3333}}
  >{\raggedright\arraybackslash}p{(\linewidth - 4\tabcolsep) * \real{0.3333}}@{}}
\toprule\noalign{}
\begin{minipage}[b]{\linewidth}\raggedright
Theorem
\end{minipage} & \begin{minipage}[b]{\linewidth}\raggedright
Connects\ldots{}
\end{minipage} & \begin{minipage}[b]{\linewidth}\raggedright
Common Use Case
\end{minipage} \\
\midrule\noalign{}
\endhead
\bottomrule\noalign{}
\endlastfoot
\textbf{Green's} & 1D Closed Line \(\iff\) 2D Plane Area & Evaluating
planar work integrals. \\
\textbf{Stokes'} & 1D Closed Line \(\iff\) 3D Open Surface & Evaluating
work in 3D space by calculating flux of curl. \\
\textbf{Gauss's} & 2D Closed Surface \(\iff\) 3D Volume & Evaluating
outward flux through a closed solid by integrating divergence. \\
\end{longtable}

\subsection{Exam Patterns}\label{exam-patterns-16}

\begin{itemize}
\tightlist
\item
  Like Green's and Stokes', ``verify'' means you must compute the volume
  integral and the surface integral independently.
\item
  The cylinder problem is a classic. Always split the cylinder into top,
  bottom, and curved surfaces, treating them entirely separately before
  summing them.
\item
  Cylindrical coordinates (\(r, \phi, z\)) are absolutely essential
  here.
\end{itemize}

\section{B01: Matrix Basics and
Operations}\label{b01-matrix-basics-and-operations}

\begin{quote}
\textbf{Exam Weight}: Tier 3 \textbar{} \textbf{Typical Marks}: 2--3
(definition questions)
\end{quote}

These fundamentals rarely appear as standalone problems. But they are
the foundation for everything else. Knowing them helps you avoid
careless mistakes.

\subsection{What is a Matrix?}\label{what-is-a-matrix}

A matrix is a rectangular array of numbers arranged in rows and columns.
A matrix with \(m\) rows and \(n\) columns has order \(m \times n\).
\[
A =
\begin{bmatrix}
a_{11} & a_{12} & \dots & a_{1n} \\
a_{21} & a_{22} & \dots & a_{2n} \\
\vdots & \vdots & \ddots & \vdots \\
a_{m1} & a_{m2} & \dots & a_{mn}
\end{bmatrix}
\]
The element \(a_{ij}\) sits in row \(i\) and column \(j\).

\subsection{Types of Matrices}\label{types-of-matrices}

\begin{longtable}[]{@{}lll@{}}
\toprule\noalign{}
Type & Condition & Example \\
\midrule\noalign{}
\endhead
\bottomrule\noalign{}
\endlastfoot
\textbf{Square} & \(m = n\) & 3x3 matrix \\
\textbf{Row} & \(m = 1\) & \([1 \; 2 \; 3]\) \\
\textbf{Column} & \(n = 1\) & \([1, 2]^T\) \\
\textbf{Diagonal} & & \\
\end{longtable}
\[
a_{ij} = 0
\]
for \(i \neq j\) \textbar{} Only diagonal entries non-zero \textbar{}
\textbar{} \textbf{Scalar} \textbar{} Diagonal with all diagonal entries
equal \textbar{} \(kI\) \textbar{} \textbar{} \textbf{Identity}
\textbar{} Diagonal with all 1s \textbar{} \(I_n\) \textbar{} \textbar{}
\textbf{Null/Zero} \textbar{} All entries are 0 \textbar{} \(O\)
\textbar{} \textbar{} \textbf{Upper triangular} \textbar{}
\[
a_{ij} = 0
\]
for \(i > j\) \textbar{} Zeros below diagonal \textbar{} \textbar{}
\textbf{Lower triangular} \textbar{}
\[
a_{ij} = 0
\]
for \(i < j\) \textbar{} Zeros above diagonal \textbar{}

\subsection{Matrix Operations}\label{matrix-operations}

\subsubsection{Addition and Subtraction}\label{addition-and-subtraction}

Add or subtract matrices of the same order by adding or subtracting
corresponding entries:
\[
(A + B)_{ij} = a_{ij} + b_{ij}
\]
Addition is commutative: \(A + B = B + A\).

\subsubsection{Scalar Multiplication}\label{scalar-multiplication}

Multiply every entry by the scalar:
\[
(kA)_{ij} = k \cdot a_{ij}
\]
\subsubsection{Matrix Multiplication}\label{matrix-multiplication}

\(A\) (\(m \times n\)) times \(B\) (\(n \times p\)) gives \(C\)
(\(m \times p\)). The \((i,j)\)-th entry of \(C\) is:
\[
c_{ij} = \sum_{k=1}^{n} a_{ik} b_{kj}
\]
Multiplication is NOT commutative: \(AB \neq BA\) in general.

Multiplication IS associative: \((AB)C = A(BC)\).

Multiplication IS distributive: \(A(B + C) = AB + AC\).

\subsubsection{Important: Size
Compatibility}\label{important-size-compatibility}

You can only multiply \(A_{m \times n}\) by \(B_{n \times p}\). The
inner dimensions must match. The result is \(m \times p\).

\subsection{Determinant (Quick
Reference)}\label{determinant-quick-reference}

For a 2x2 matrix:
\[
|A| =
\begin{vmatrix}
a & b \\
c & d
\end{vmatrix} = ad - bc
\]
For a 3x3 matrix, expand along row 1:
\[
|A| = a_{11}(a_{22}a_{33} - a_{23}a_{32}) - a_{12}(a_{21}a_{33} - a_{23}a_{31}) + a_{13}(a_{21}a_{32} - a_{22}a_{31})
\]
Properties: - If a row or column is all zeros,
\[
\lvert A \rvert = 0.
\]
\begin{itemize}
\tightlist
\item
  Swapping two rows changes the sign of the determinant.
\item
\end{itemize}
\[
\lvert AB \rvert = \lvert A \rvert \cdot \lvert B \rvert.
\]
\begin{itemize}
\tightlist
\item
\end{itemize}
\[
\lvert kA \rvert = k^n \lvert A \rvert
\]
for an \(n \times n\) matrix.

\section{B02: Transpose, Symmetric, and Skew-Symmetric
Matrices}\label{b02-transpose-symmetric-and-skew-symmetric-matrices}

\begin{quote}
\textbf{Exam Weight}: Tier 2 \textbar{} \textbf{Appeared in}: 2018,
2020, 2021 \textbar{} \textbf{Typical Marks}: 3--6
\end{quote}

These definitions are easy marks. The decomposition theorem (symmetric +
skew-symmetric) has appeared twice.

\subsection{Transpose of a Matrix}\label{transpose-of-a-matrix}

\subsubsection{Definition}\label{definition-5}

The transpose of matrix \(A\) (written \(A^T\) or \(A'\)) is obtained by
swapping rows and columns. If \(A\) is \(m \times n\), then \(A^T\) is
\(n \times m\).

The \((i,j)\)-th element of \(A^T\) equals the \((j,i)\)-th element of
\(A\):
\[
(A^T)_{ij} = a_{ji}
\]
\subsubsection{Properties of Transpose}\label{properties-of-transpose}

\begin{enumerate}
\def\labelenumi{\arabic{enumi}.}
\tightlist
\item
\end{enumerate}
\[
(A^T)^T = A
\]
\begin{enumerate}
\def\labelenumi{\arabic{enumi}.}
\setcounter{enumi}{1}
\tightlist
\item
\end{enumerate}
\[
(A + B)^T = A^T + B^T
\]
\begin{enumerate}
\def\labelenumi{\arabic{enumi}.}
\setcounter{enumi}{2}
\tightlist
\item
\end{enumerate}
\[
(kA)^T = kA^T
\]
\begin{enumerate}
\def\labelenumi{\arabic{enumi}.}
\setcounter{enumi}{3}
\tightlist
\item
\end{enumerate}
\[
(AB)^T = B^T A^T
\]
(order reverses) 5.
\[
(A^{-1})^T = (A^T)^{-1}
\]

\subsection{Symmetric Matrix}\label{symmetric-matrix}

\subsubsection{Definition}\label{definition-6}

A square matrix \(A\) is symmetric if it equals its transpose:
\[
A^T = A
\]
This means
\[
a_{ij} = a_{ji}
\]
for all \(i, j\). The matrix is a mirror image across its main diagonal.

\subsubsection{Example}\label{example}
\[
A =
\begin{bmatrix}
1 & 2 & 3 \\
2 & 5 & 6 \\
3 & 6 & 9
\end{bmatrix}
\]
Notice: position \((1,2)\) and \((2,1)\) both have 2. Position \((1,3)\)
and \((3,1)\) both have 3.

\subsubsection{Key Properties}\label{key-properties}

\begin{itemize}
\tightlist
\item
  Symmetric matrices have only real eigenvalues.
\item
  They can always be diagonalized.
\item
  They are very common in physics and engineering.
\end{itemize}

\subsection{Skew-Symmetric Matrix}\label{skew-symmetric-matrix}

\subsubsection{Definition}\label{definition-7}

A square matrix \(A\) is skew-symmetric if:
\[
A^T = -A
\]
This means
\[
a_{ij} = -a_{ji}
\]
for all \(i, j\). The diagonal entries must be zero (since
\[
a_{ii} = -a_{ii}
\]
forces
\[
a_{ii} = 0
\]
).

\subsubsection{Example}\label{example-1}
\[
A =
\begin{bmatrix}
0 & -3 & 5 \\
3 & 0 & -7 \\
-5 & 7 & 0
\end{bmatrix}
\]
All diagonal entries are 0. Each pair of symmetric positions has
opposite signs.

\subsection{Decomposition Theorem (PYQ 2020, 2021 ---
Repeated)}\label{decomposition-theorem-pyq-2020-2021-repeated}

\subsubsection{Statement}\label{statement-1}

Every square matrix can be uniquely expressed as the sum of a symmetric
matrix and a skew-symmetric matrix.

\subsubsection{Proof}\label{proof-1}

Let \(A\) be any square matrix. Define:
\[
P = \frac{1}{2}(A + A^T) \quad \text{and} \quad Q = \frac{1}{2}(A - A^T)
\]
\textbf{Step 1}: Show \(A = P + Q\).
\[
P + Q = \frac{1}{2}(A + A^T) + \frac{1}{2}(A - A^T) = \frac{1}{2}(2A) = A \quad \checkmark
\]
\textbf{Step 2}: Show \(P\) is symmetric.
\[
P^T = \left[\frac{1}{2}(A + A^T)\right]^T = \frac{1}{2}(A^T + (A^T)^T) = \frac{1}{2}(A^T + A) = P \quad \checkmark
\]
\textbf{Step 3}: Show \(Q\) is skew-symmetric.
\[
Q^T = \left[\frac{1}{2}(A - A^T)\right]^T = \frac{1}{2}(A^T - A) = -\frac{1}{2}(A - A^T) = -Q \quad \checkmark
\]
\textbf{Step 4}: Prove uniqueness. Suppose \(A = P' + Q'\) where \(P'\)
is symmetric and \(Q'\) is skew-symmetric.

Taking transpose:
\[
A^T = P' - Q'.
\]
Adding the two equations:
\[
A + A^T = 2P'
\]
, so
\[
P' = \frac{1}{2}(A + A^T) = P.
\]
Subtracting:
\[
A - A^T = 2Q'
\]
, so
\[
Q' = \frac{1}{2}(A - A^T) = Q.
\]
The decomposition is unique.

\subsubsection{Worked Example}\label{worked-example}

Express
\[
A =
\begin{bmatrix}
1 & 2 \\
3 & 4
\end{bmatrix}
\]
as symmetric + skew-symmetric.

Symmetric part:
\[
P = \frac{1}{2}(A + A^T) = \frac{1}{2}\left(
\begin{bmatrix}
1 & 2 \\
3 & 4
\end{bmatrix} + \begin{bmatrix} 1 & 3 \\ 2 & 4 \end{bmatrix}\right) = \frac{1}{2}\begin{bmatrix} 2 & 5 \\ 5 & 8 \end{bmatrix} = \begin{bmatrix} 1 & 5/2 \\ 5/2 & 4 \end{bmatrix}
\]
Skew-symmetric part:
\[
Q = \frac{1}{2}(A - A^T) = \frac{1}{2}\left(
\begin{bmatrix}
1 & 2 \\
3 & 4
\end{bmatrix} - \begin{bmatrix} 1 & 3 \\ 2 & 4 \end{bmatrix}\right) = \frac{1}{2}\begin{bmatrix} 0 & -1 \\ 1 & 0 \end{bmatrix} = \begin{bmatrix} 0 & -1/2 \\ 1/2 & 0 \end{bmatrix}
\]
Verify:
\[
P + Q =
\begin{bmatrix}
1 & 2 \\
3 & 4
\end{bmatrix} = A.
\]
Correct.

\subsection{Exam Patterns}\label{exam-patterns-17}

\begin{itemize}
\tightlist
\item
  ``Define symmetric, skew-symmetric, Hermitian matrices'' is a 2-3 mark
  definition question. Write the formula and one example each.
\item
  The decomposition proof (symmetric + skew-symmetric) appeared in 2020
  and 2021. Memorize the formulas
\end{itemize}
\[
P = \frac{1}{2}(A + A^T)
\]
and
\[
Q = \frac{1}{2}(A - A^T).
\]
\begin{itemize}
\tightlist
\item
  The uniqueness part adds 1-2 extra marks if you include it.
\end{itemize}

\section{B03: Hermitian, Skew-Hermitian, and Special
Matrices}\label{b03-hermitian-skew-hermitian-and-special-matrices}

\begin{quote}
\textbf{Exam Weight}: Tier 2 \textbar{} \textbf{Appeared in}: 2018, 2023
\textbar{} \textbf{Typical Marks}: 2--5
\end{quote}

Definition questions about matrix types are easy marks. The Hermitian
decomposition proof appeared in 2023.

\subsection{Hermitian Matrix}\label{hermitian-matrix}

\subsubsection{Definition}\label{definition-8}

A square matrix \(A\) with complex entries is Hermitian if it equals its
conjugate transpose:
\[
A^\dagger = (\bar{A})^T = A
\]
Here \(\bar{A}\) is the complex conjugate (replace \(i\) with \(-i\) in
every entry). Then \(A^\dagger\) transposes the result.

\subsubsection{Properties}\label{properties-1}

\begin{itemize}
\tightlist
\item
  Diagonal entries must be real (since
\end{itemize}
\[
a_{ii} = \overline{a_{ii}}
\]
implies \(a_{ii} \in \mathbb{R}\)). - Off-diagonal entries satisfy
\[
a_{ij} = \overline{a_{ji}}.
\]
\begin{itemize}
\tightlist
\item
  All eigenvalues of a Hermitian matrix are real.
\end{itemize}

\subsubsection{Example}\label{example-2}
\[
A =
\begin{pmatrix}
2 & 1+i \\
1-i & 3
\end{pmatrix}
\]
Check:
\[
\bar{A} =
\begin{pmatrix}
2 & 1-i \\
1+i & 3
\end{pmatrix}, \quad (\bar{A})^T = \begin{pmatrix} 2 & 1+i \\ 1-i & 3 \end{pmatrix} = A
\]
Hermitian.

\subsection{Skew-Hermitian Matrix}\label{skew-hermitian-matrix}

\subsubsection{Definition}\label{definition-9}

A square matrix \(A\) is skew-Hermitian if:
\[
A^\dagger = -A
\]
\subsubsection{Properties}\label{properties-2}

\begin{itemize}
\tightlist
\item
  Diagonal entries must be pure imaginary or zero (since
\end{itemize}
\[
\overline{a_{ii}} = -a_{ii}
\]
). - Off-diagonal entries satisfy
\[
a_{ij} = -\overline{a_{ji}}.
\]
\subsubsection{Example}\label{example-3}
\[
A =
\begin{pmatrix}
2i & 1+i \\
-1+i & 0
\end{pmatrix}
\]

\subsection{Hermitian Decomposition Theorem (PYQ
2023)}\label{hermitian-decomposition-theorem-pyq-2023}

\subsubsection{Statement}\label{statement-2}

Every square matrix can be expressed as the sum of a Hermitian matrix
and a skew-Hermitian matrix.

\subsubsection{Proof}\label{proof-2}

Let \(A\) be any square matrix. Define:
\[
P = \frac{1}{2}(A + A^\dagger) \quad \text{and} \quad Q = \frac{1}{2}(A - A^\dagger)
\]
\textbf{Show \(P\) is Hermitian}:
\[
P^\dagger = \frac{1}{2}(A^\dagger + (A^\dagger)^\dagger) = \frac{1}{2}(A^\dagger + A) = P
\]
\textbf{Show \(Q\) is skew-Hermitian}:
\[
Q^\dagger = \frac{1}{2}(A^\dagger - (A^\dagger)^\dagger) = \frac{1}{2}(A^\dagger - A) = -Q
\]
And \(P + Q = A\).

This is the same structure as the symmetric decomposition. Replace
transpose with conjugate transpose.

\subsection{Special Matrix Types}\label{special-matrix-types}

\subsubsection{Nilpotent Matrix}\label{nilpotent-matrix}

A square matrix \(A\) is nilpotent if there exists a positive integer
\(k\) such that:
\[
A^k = O
\]
where \(O\) is the zero matrix. The smallest such \(k\) is the index of
nilpotency.

\textbf{Example}:
\[
A =
\begin{bmatrix}
0 & 1 \\
0 & 0
\end{bmatrix}. \quad A^2 = \begin{bmatrix} 0 & 0 \\ 0 & 0 \end{bmatrix} = O.
\]
Nilpotent with index 2.

All eigenvalues of a nilpotent matrix are zero.

\subsubsection{Involutory Matrix}\label{involutory-matrix}

A square matrix \(A\) is involutory if:
\[
A^2 = I
\]
This means \(A\) is its own inverse:
\[
A^{-1} = A.
\]
\textbf{Example}:
\[
A =
\begin{bmatrix}
1 & 0 \\
0 & -1
\end{bmatrix}. \quad A^2 = I.
\]
\subsubsection{Idempotent Matrix}\label{idempotent-matrix}

A square matrix \(A\) is idempotent if:
\[
A^2 = A
\]
Applying the transformation twice gives the same result as applying it
once. Projection matrices are idempotent.

\textbf{Example}:
\[
A =
\begin{bmatrix}
1 & 0 \\
0 & 0
\end{bmatrix}. \quad A^2 = A.
\]
The eigenvalues of an idempotent matrix are only 0 or 1.

\subsubsection{Periodic Matrix}\label{periodic-matrix}

A square matrix \(A\) is periodic with period \(k\) if:
\[
A^{k+1} = A
\]
An idempotent matrix is periodic with period 1.

\subsubsection{Orthogonal Matrix}\label{orthogonal-matrix}

A square matrix \(A\) is orthogonal if:
\[
A^T A = AA^T = I
\]
This means
\[
A^{-1} = A^T
\]
. The columns form an orthonormal set. The determinant is \(\pm 1\).

\textbf{Example}: Rotation matrices are orthogonal.

\subsubsection{Unitary Matrix}\label{unitary-matrix}

A square matrix \(A\) is unitary if:
\[
A^\dagger A = AA^\dagger = I
\]
This is the complex version of orthogonal.
\[
A^{-1} = A^\dagger.
\]

\subsection{Exam Patterns}\label{exam-patterns-18}

\begin{itemize}
\tightlist
\item
  ``Define commutative, idempotent, involutory matrix'' appeared in 2023
  for 2 marks. Just write the definition and one example.
\item
  ``Define symmetric, skew-symmetric, Hermitian'' appeared in 2018 for 3
  marks.
\item
  The Hermitian decomposition proof appeared in 2023 for 5 marks.
\item
  These are easy marks. Memorize the one-line definitions.
\end{itemize}

\section{B04: Matrix Proofs}\label{b04-matrix-proofs}

\begin{quote}
\textbf{Exam Weight}: Tier 2 \textbar{} \textbf{Appeared in}: 2017,
2018, 2020, 2023 \textbar{} \textbf{Typical Marks}: 3--5
\end{quote}

These are standard proofs that appear regularly. They test whether you
understand matrix algebra rules. Each proof is short and worth learning
by heart.

\subsection{Proof 1: Transpose of a Product (PYQ
2020)}\label{proof-1-transpose-of-a-product-pyq-2020}

\subsubsection{Statement}\label{statement-3}

For matrices \(A\) (of size \(m \times n\)) and \(B\) (of size
\(n \times p\)):
\[
(AB)^T = B^T A^T
\]
The order reverses when you transpose a product.

\subsubsection{Proof}\label{proof-3}

Let \(C = AB\). The \((i,j)\)-th element of \(C\) is:
\[
c_{ij} = \sum_{k=1}^{n} a_{ik} b_{kj}
\]
The \((i,j)\)-th element of \(C^T\) is \(c_{ji}\):
\[
(C^T)_{ij} = c_{ji} = \sum_{k=1}^{n} a_{jk} b_{ki}
\]
Now consider \(B^T A^T\). Its \((i,j)\)-th element is:
\[
(B^T A^T)_{ij} = \sum_{k=1}^{n} (B^T)_{ik}(A^T)_{kj} = \sum_{k=1}^{n} b_{ki} \cdot a_{jk} = \sum_{k=1}^{n} a_{jk} b_{ki}
\]
Since
\[
(C^T)_{ij} = (B^T A^T)_{ij}
\]
for all \(i,j\):
\[
(AB)^T = B^T A^T \qquad \blacksquare
\]

\subsection{Proof 2: Inverse of a Product (PYQ 2018, 2023 ---
Repeated)}\label{proof-2-inverse-of-a-product-pyq-2018-2023-repeated}

\subsubsection{Statement}\label{statement-4}

If \(A\) and \(B\) are non-singular matrices of the same order:
\[
(AB)^{-1} = B^{-1}A^{-1}
\]
\subsubsection{Proof}\label{proof-4}

We show that \(B^{-1}A^{-1}\) satisfies the definition of the inverse of
\(AB\).

\textbf{Left product}:
\[
(AB)(B^{-1}A^{-1}) = A(BB^{-1})A^{-1} = AIA^{-1} = AA^{-1} = I
\]
\textbf{Right product}:
\[
(B^{-1}A^{-1})(AB) = B^{-1}(A^{-1}A)B = B^{-1}IB = B^{-1}B = I
\]
Since
\[
(AB)(B^{-1}A^{-1}) = I
\]
and
\[
(B^{-1}A^{-1})(AB) = I:
\]
\[
(AB)^{-1} = B^{-1}A^{-1} \qquad \blacksquare
\]

\subsection{Proof 3: Power of Inverse (PYQ
2018)}\label{proof-3-power-of-inverse-pyq-2018}

\subsubsection{Statement}\label{statement-5}

For any invertible matrix \(A\) and positive integer \(n\):
\[
(A^{-1})^n = (A^n)^{-1}
\]
\subsubsection{Proof (by induction)}\label{proof-by-induction}

\textbf{Base case} (\(n = 1\)):
\[
(A^{-1})^1 = (A^1)^{-1}
\]
. True.

\textbf{Inductive step}: Assume
\[
(A^{-1})^k = (A^k)^{-1}
\]
holds for some \(k\).

For \(n = k + 1\):
\[
(A^{k+1})^{-1} = (A^k \cdot A)^{-1} = A^{-1} \cdot (A^k)^{-1}
\]
using Proof 2. By the inductive hypothesis,
\[
(A^k)^{-1} = (A^{-1})^k:
\]
\[
A^{-1} \cdot (A^{-1})^k = (A^{-1})^{k+1}
\]
The statement holds for \(k + 1\). By induction, it holds for all
positive integers \(n\).
\[
\blacksquare
\]

\subsection{Proof 4: Rotation Matrix Property (PYQ
2017)}\label{proof-4-rotation-matrix-property-pyq-2017}

\subsubsection{Statement}\label{statement-6}

If
\[
A_\alpha =
\begin{pmatrix}
\cos\alpha & \sin\alpha \\
-\sin\alpha & \cos\alpha
\end{pmatrix}
\]
then
\[
A_\alpha \cdot A_\beta = A_{\alpha+\beta} = A_\beta \cdot A_\alpha.
\]
\subsubsection{Proof}\label{proof-5}

Multiply \(A_\alpha \cdot A_\beta\):
\[
A_\alpha \cdot A_\beta =
\begin{pmatrix}
\cos\alpha\cos\beta - \sin\alpha\sin\beta & \cos\alpha\sin\beta + \sin\alpha\cos\beta \\
-\sin\alpha\cos\beta - \cos\alpha\sin\beta & -\sin\alpha\sin\beta + \cos\alpha\cos\beta
\end{pmatrix}
\]
Apply the angle addition formulas:
\[
\cos(\alpha+\beta) = \cos\alpha\cos\beta - \sin\alpha\sin\beta
\]
\[
\sin(\alpha+\beta) = \sin\alpha\cos\beta + \cos\alpha\sin\beta
\]
Substitute:
\[
A_\alpha \cdot A_\beta =
\begin{pmatrix}
\cos(\alpha+\beta) & \sin(\alpha+\beta) \\
-\sin(\alpha+\beta) & \cos(\alpha+\beta)
\end{pmatrix} = A_{\alpha+\beta}
\]
Since \(\alpha + \beta = \beta + \alpha\):
\[
A_{\alpha+\beta} = A_{\beta+\alpha} = A_\beta \cdot A_\alpha \qquad \blacksquare
\]

\subsection{Proof 5: Singular Matrix
Definition}\label{proof-5-singular-matrix-definition}

\subsubsection{Statement}\label{statement-7}

A square matrix \(A\) is singular if
\[
\lvert A \rvert = 0
\]
. It is non-singular if \(\lvert A \rvert \neq 0\).

A singular matrix has no inverse. A non-singular matrix is invertible
and
\[
A^{-1} = \frac{1}{\lvert A \rvert}\text{adj}(A).
\]
This definition is often asked as a 1-2 mark opener before a proof
question.

\subsection{Exam Patterns}\label{exam-patterns-19}

\begin{itemize}
\tightlist
\item
  The
\end{itemize}
\[
(AB)^{-1} = B^{-1}A^{-1}
\]
proof has appeared in 2018 and 2023. This is the most important proof. -
The
\[
(AB)^T = B^T A^T
\]
proof appeared in 2020. - The rotation matrix property appeared in 2017.
- These proofs are short (5-8 lines each). Write them neatly with clear
steps. - Always state what you need to prove before starting.

\section{B05: Row Echelon Form and Rank of a
Matrix}\label{b05-row-echelon-form-and-rank-of-a-matrix}

\begin{quote}
\textbf{Exam Weight}: Tier 1 \textbar{} \textbf{Appeared in}: 2017,
2018, 2019, 2020, 2021, 2023, 2024 (all 7 papers) + CT \textbar{}
\textbf{Typical Marks}: 3--6
\end{quote}

This topic has appeared in every single exam paper. It is the most
consistent topic in the entire course. You cannot afford to skip it.

\subsection{What is the Rank of a
Matrix?}\label{what-is-the-rank-of-a-matrix}

\subsubsection{Definition}\label{definition-10}

The rank of a matrix \(A\) is the maximum number of linearly independent
row vectors in it. It also equals the maximum number of linearly
independent column vectors. We write it as \(\rho(A)\) or
\(\text{rank}(A)\).

\subsubsection{Key Idea}\label{key-idea}

Think of rank as ``how many rows actually carry new information.'' Some
rows in a matrix might just be multiples or combinations of other rows.
Those rows add no new information. The rank counts only the rows that
are truly independent from each other.

A zero matrix has rank 0. A non-zero matrix has rank at least 1.

For an \(m \times n\) matrix, the rank can be at most \(\min(m, n)\). So
a \(3 \times 4\) matrix has rank at most 3. A \(4 \times 3\) matrix also
has rank at most 3.

\subsubsection{Why Rank Matters}\label{why-rank-matters}

Rank shows up everywhere in this course: - It decides if a system of
equations has a solution (rank conditions) - It tells you the dimension
of the column space and row space - It connects to the number of free
variables in a system - It is the order of the largest non-zero minor

\subsection{Elementary Row Operations}\label{elementary-row-operations}

Before we find rank, we need row operations. These are tools that change
a matrix into a simpler form without changing its rank.

Three types of elementary row operations exist:

\begin{enumerate}
\def\labelenumi{\arabic{enumi}.}
\tightlist
\item
  \textbf{Swap two rows}: \(R_i \leftrightarrow R_j\)
\item
  \textbf{Multiply a row by a non-zero scalar}: \(R_i \to kR_i\) where
  \(k \neq 0\)
\item
  \textbf{Add a scalar multiple of one row to another}:
  \(R_i \to R_i + kR_j\)
\end{enumerate}

Two matrices are called \textbf{equivalent} (written \(A \sim B\)) if
one can be reached from the other through row operations. Equivalent
matrices always have the same rank.

\subsection{Row Echelon Form}\label{row-echelon-form}

\subsubsection{Definition}\label{definition-11}

A matrix is in row echelon form if it satisfies these three rules:

\begin{enumerate}
\def\labelenumi{\arabic{enumi}.}
\tightlist
\item
  All zero rows are at the bottom.
\item
  The first non-zero entry in each row (called the \textbf{pivot}) is
  strictly to the right of the pivot in the row above it.
\item
  All entries below each pivot are zero.
\end{enumerate}

\subsubsection{Example of Echelon Form}\label{example-of-echelon-form}

This matrix is in echelon form:
\[
\begin{bmatrix}
1 & 2 & 3 & 4 \\
0 & 1 & 2 & 3 \\
0 & 0 & 1 & 1 \\
0 & 0 & 0 & 0
\end{bmatrix}
\]
Notice how the pivots (the leading 1s) staircase to the right. The zero
row sits at the bottom.

This matrix is NOT in echelon form:
\[
\begin{bmatrix}
1 & 0 & 1 & 1 \\
0 & 1 & 1 & 1 \\
0 & 0 & 0 & 0 \\
0 & 0 & 0 & 2
\end{bmatrix}
\]
The zero row is not at the bottom. The last row's pivot is not to the
right of the row above it.

\subsection{Finding Rank by Echelon
Form}\label{finding-rank-by-echelon-form}

\subsubsection{The Method}\label{the-method-1}

\begin{enumerate}
\def\labelenumi{\arabic{enumi}.}
\tightlist
\item
  Use elementary row operations to convert the matrix to echelon form.
\item
  Count the number of non-zero rows.
\item
  That count is the rank.
\end{enumerate}

This is the fastest and most exam-friendly method for finding rank.
Every exam question on rank uses this approach.

\subsubsection{Worked Example 1: Rank of a 3x4
Matrix}\label{worked-example-1-rank-of-a-3x4-matrix}

\textbf{Problem}: Find the rank of the matrix (PYQ 2020):
\[
A =
\begin{bmatrix}
6 & 1 & 8 & 3 \\
2 & 1 & 0 & 2 \\
4 & -1 & -8 & -3
\end{bmatrix}
\]
\textbf{Solution}:

Swap \(R_1\) and \(R_2\) to get a small leading entry:
\[
\begin{bmatrix}
2 & 1 & 0 & 2 \\
6 & 1 & 8 & 3 \\
4 & -1 & -8 & -3
\end{bmatrix}
\]
Clear the first column. Apply \(R_2 \to R_2 - 3R_1\) and
\(R_3 \to R_3 - 2R_1\):
\[
\begin{bmatrix}
2 & 1 & 0 & 2 \\
0 & -2 & 8 & -3 \\
0 & -3 & -8 & -7
\end{bmatrix}
\]
Clear the second column in row 3. Apply \(R_3 \to 2R_3 - 3R_2\):
\[
\begin{bmatrix}
2 & 1 & 0 & 2 \\
0 & -2 & 8 & -3 \\
0 & 0 & -40 & -5
\end{bmatrix}
\]
All three rows are non-zero. So:
\[
\rho(A) = 3
\]

\subsubsection{Worked Example 2: Rank of a 4x4
Matrix}\label{worked-example-2-rank-of-a-4x4-matrix}

\textbf{Problem}: Find the rank of the matrix (PYQ 2018):
\[
A =
\begin{pmatrix}
1 & -2 & 1 & -1 \\
1 & 1 & -2 & 3 \\
4 & 1 & -5 & 8 \\
5 & -7 & 2 & -1
\end{pmatrix}
\]
\textbf{Solution}:

Clear the first column. Apply \(R_2 \to R_2 - R_1\),
\(R_3 \to R_3 - 4R_1\), \(R_4 \to R_4 - 5R_1\):
\[
\begin{pmatrix}
1 & -2 & 1 & -1 \\
0 & 3 & -3 & 4 \\
0 & 9 & -9 & 12 \\
0 & 3 & -3 & 4
\end{pmatrix}
\]
Clear the second column. Apply \(R_3 \to R_3 - 3R_2\) and
\(R_4 \to R_4 - R_2\):
\[
\begin{pmatrix}
1 & -2 & 1 & -1 \\
0 & 3 & -3 & 4 \\
0 & 0 & 0 & 0 \\
0 & 0 & 0 & 0
\end{pmatrix}
\]
Two non-zero rows remain. So:
\[
\rho(A) = 2
\]
Notice that rows 3 and 4 were multiples of row 2 after the first set of
operations. So they carried no new information.

\subsubsection{Worked Example 3: Rank of a 4x5
Matrix}\label{worked-example-3-rank-of-a-4x5-matrix}

\textbf{Problem}: Find the rank of the matrix (PYQ 2017):
\[
A =
\begin{bmatrix}
0 & 0 & 1 & -3 & -2 \\
0 & 1 & 2 & 6 & 0 \\
0 & 2 & 3 & 9 & 2 \\
0 & 1 & 1 & 3 & 2
\end{bmatrix}
\]
\textbf{Solution}:

The first column is all zeros. Swap \(R_1 \leftrightarrow R_2\) to put a
non-zero entry first:
\[
\begin{bmatrix}
0 & 1 & 2 & 6 & 0 \\
0 & 0 & 1 & -3 & -2 \\
0 & 2 & 3 & 9 & 2 \\
0 & 1 & 1 & 3 & 2
\end{bmatrix}
\]
Clear below the first pivot. Apply \(R_3 \to R_3 - 2R_1\) and
\(R_4 \to R_4 - R_1\):
\[
\begin{bmatrix}
0 & 1 & 2 & 6 & 0 \\
0 & 0 & 1 & -3 & -2 \\
0 & 0 & -1 & -3 & 2 \\
0 & 0 & -1 & -3 & 2
\end{bmatrix}
\]
Clear below the second pivot. Apply \(R_3 \to R_3 + R_2\) and
\(R_4 \to R_4 + R_2\):
\[
\begin{bmatrix}
0 & 1 & 2 & 6 & 0 \\
0 & 0 & 1 & -3 & -2 \\
0 & 0 & 0 & -6 & 0 \\
0 & 0 & 0 & -6 & 0
\end{bmatrix}
\]
Apply \(R_4 \to R_4 - R_3\):
\[
\begin{bmatrix}
0 & 1 & 2 & 6 & 0 \\
0 & 0 & 1 & -3 & -2 \\
0 & 0 & 0 & -6 & 0 \\
0 & 0 & 0 & 0 & 0
\end{bmatrix}
\]
Three non-zero rows remain. So:
\[
\rho(A) = 3
\]

\subsection{Finding Rank by Minor
Method}\label{finding-rank-by-minor-method}

\subsubsection{Definition of a Minor}\label{definition-of-a-minor}

Take any \(r\) rows and \(r\) columns from an \(m \times n\) matrix. The
determinant of the resulting \(r \times r\) submatrix is a minor of
order \(r\).

\subsubsection{The Method}\label{the-method-2}

\begin{enumerate}
\def\labelenumi{\arabic{enumi}.}
\tightlist
\item
  Start with the highest possible order of a square submatrix.
\item
  Compute its determinant. If it is non-zero, the rank equals this
  order.
\item
  If all minors of that order are zero, drop down one order and check
  again.
\item
  The rank is the order of the largest non-zero minor.
\end{enumerate}

\subsubsection{When to Use This}\label{when-to-use-this}

The minor method is slower than echelon form for large matrices. But it
is useful for small matrices or when the problem specifically asks for
``rank by minors.'' In exams, the echelon method is almost always
preferred.

\subsubsection{Worked Example: Minor
Method}\label{worked-example-minor-method}

\textbf{Problem}: Find the rank of:
\[
A =
\begin{bmatrix}
1 & 2 & 3 \\
2 & 3 & 4 \\
3 & 4 & 5
\end{bmatrix}
\]
\textbf{Solution}:

The matrix is \(3 \times 3\). Start by checking if the full
\(3 \times 3\) determinant is non-zero:
\[
|A| = 1(15 - 16) - 2(10 - 12) + 3(8 - 9) = -1 + 4 - 3 = 0
\]
The \(3 \times 3\) minor is zero. So rank is not 3.

Check a \(2 \times 2\) minor. Take the top-left \(2 \times 2\)
submatrix:
\[
\begin{vmatrix}
1 & 2 \\
2 & 3
\end{vmatrix} = 3 - 4 = -1 \neq 0
\]
A non-zero \(2 \times 2\) minor exists. So:
\[
\rho(A) = 2
\]

\subsection{Nullity}\label{nullity}

\subsubsection{Definition}\label{definition-12}

For a square matrix of order \(n\), the nullity is defined as:
\[
N(A) = n - \rho(A)
\]
Nullity tells you the dimension of the null space (the set of all
solutions to \(AX = 0\)). It equals the number of free variables in the
system.

\subsubsection{Example}\label{example-4}

If \(A\) is a \(4 \times 4\) matrix with \(\rho(A) = 3\), then:
\[
N(A) = 4 - 3 = 1
\]
This means the system \(AX = 0\) has one free variable and its solution
space is one-dimensional.

\subsection{Important Properties of
Rank}\label{important-properties-of-rank}

\begin{enumerate}
\def\labelenumi{\arabic{enumi}.}
\tightlist
\item
  \(\rho(A) = 0\) if and only if \(A\) is the zero matrix.
\item
  For an \(m \times n\) matrix, \(\rho(A) \leq \min(m, n)\).
\item
\end{enumerate}
\[
\rho(A) = \rho(A^T)
\]
. Row rank equals column rank. 4. Row operations do not change the rank.
5. If \(A\) is an \(n \times n\) matrix, then \(\rho(A) = n\) if and
only if \(A\) is invertible (non-singular). 6. If \(\rho(A) < n\) for a
square matrix of order \(n\), then
\[
\lvert A \rvert = 0.
\]

\subsection{Exam Patterns}\label{exam-patterns-20}

\begin{itemize}
\tightlist
\item
  \textbf{Every paper asks rank.} The typical phrasing is: ``Define rank
  of a matrix. Find the rank of\ldots{}'' This gives you 2 marks for the
  definition and 3--4 marks for the computation.
\item
  Matrices in exams are usually \(3 \times 4\), \(4 \times 4\), or
  \(4 \times 5\).
\item
  The definition they expect is: ``The rank of a matrix is the number of
  non-zero rows in its row echelon form'' or ``the maximum number of
  linearly independent rows.''
\item
  Always swap rows first to get a 1 (or the smallest non-zero number) in
  the top-left corner. This makes arithmetic easier.
\item
  Watch for rows that become identical during reduction. They collapse
  to zeros and lower the rank.
\end{itemize}

\subsection{Must-Know Problems}\label{must-know-problems}

\subsubsection{Problem 4: Rank of a 4x4 Matrix (PYQ
2024)}\label{problem-4-rank-of-a-4x4-matrix-pyq-2024}

\textbf{Question}: Define rank. Find the rank of:
\[
A =
\begin{bmatrix}
1 & 1 & 1 & 1 \\
1 & 3 & -2 & 1 \\
2 & 0 & -3 & 2 \\
3 & 3 & -3 & 3
\end{bmatrix}
\]
\textbf{Solution}:

Apply \(R_2 \to R_2 - R_1\), \(R_3 \to R_3 - 2R_1\),
\(R_4 \to R_4 - 3R_1\):
\[
\begin{bmatrix}
1 & 1 & 1 & 1 \\
0 & 2 & -3 & 0 \\
0 & -2 & -5 & 0 \\
0 & 0 & -6 & 0
\end{bmatrix}
\]
Apply \(R_3 \to R_3 + R_2\):
\[
\begin{bmatrix}
1 & 1 & 1 & 1 \\
0 & 2 & -3 & 0 \\
0 & 0 & -8 & 0 \\
0 & 0 & -6 & 0
\end{bmatrix}
\]
Apply \(R_4 \to 4R_4 - 3R_3\):
\[
\begin{bmatrix}
1 & 1 & 1 & 1 \\
0 & 2 & -3 & 0 \\
0 & 0 & -8 & 0 \\
0 & 0 & 0 & 0
\end{bmatrix}
\]
Three non-zero rows. So \(\rho(A) = 3\).

\subsubsection{Problem 5: Rank of a 3x4 Matrix (PYQ
2023)}\label{problem-5-rank-of-a-3x4-matrix-pyq-2023}

\textbf{Question}: Find the rank of:
\[
A =
\begin{bmatrix}
1 & 2 & 3 & 2 \\
2 & 3 & 5 & 1 \\
1 & 3 & 4 & 5
\end{bmatrix}
\]
\textbf{Solution}:

Apply \(R_2 \to R_2 - 2R_1\) and \(R_3 \to R_3 - R_1\):
\[
\begin{bmatrix}
1 & 2 & 3 & 2 \\
0 & -1 & -1 & -3 \\
0 & 1 & 1 & 3
\end{bmatrix}
\]
Apply \(R_3 \to R_3 + R_2\):
\[
\begin{bmatrix}
1 & 2 & 3 & 2 \\
0 & -1 & -1 & -3 \\
0 & 0 & 0 & 0
\end{bmatrix}
\]
Two non-zero rows. So \(\rho(A) = 2\).

\subsubsection{Problem 6: Echelon Form and Canonical Form
(CT)}\label{problem-6-echelon-form-and-canonical-form-ct}

\textbf{Question}: Find the echelon form, row canonical form, and rank
of:
\[
A =
\begin{bmatrix}
6 & 3 & -4 \\
-4 & 1 & -6 \\
1 & 2 & -5
\end{bmatrix}
\]
\textbf{Solution}:

\textbf{Step 1: Echelon form.}

Swap \(R_1 \leftrightarrow R_3\):
\[
\begin{bmatrix}
1 & 2 & -5 \\
-4 & 1 & -6 \\
6 & 3 & -4
\end{bmatrix}
\]
Apply \(R_2 \to R_2 + 4R_1\) and \(R_3 \to R_3 - 6R_1\):
\[
\begin{bmatrix}
1 & 2 & -5 \\
0 & 9 & -26 \\
0 & -9 & 26
\end{bmatrix}
\]
Apply \(R_3 \to R_3 + R_2\):
\[
\begin{bmatrix}
1 & 2 & -5 \\
0 & 9 & -26 \\
0 & 0 & 0
\end{bmatrix}
\]
Scale
\[
R_2 \to \frac{1}{9}R_2:
\]
\[
\begin{bmatrix}
1 & 2 & -5 \\
0 & 1 & -\frac{26}{9} \\
0 & 0 & 0
\end{bmatrix}
\]
This is the echelon form.

\textbf{Step 2: Row canonical form.}

Clear above the second pivot. Apply \(R_1 \to R_1 - 2R_2\):
\[
R_1 = \begin{bmatrix} 1 & 0 & -5 + \frac{52}{9} \end{bmatrix} = \begin{bmatrix} 1 & 0 & \frac{7}{9} \end{bmatrix}
\]
The row canonical form is:
\[
\begin{bmatrix}
1 & 0 & \frac{7}{9} \\
0 & 1 & -\frac{26}{9} \\
0 & 0 & 0
\end{bmatrix}
\]
\textbf{Step 3: Rank.}

Two non-zero rows in echelon form. So \(\rho(A) = 2\).

\section{B06: Normal (Canonical) Form of a
Matrix}\label{b06-normal-canonical-form-of-a-matrix}

\begin{quote}
\textbf{Exam Weight}: Tier 1-2 \textbar{} \textbf{Appeared in}: 2019,
2021, 2023 \textbar{} \textbf{Typical Marks}: 3--6
\end{quote}

The normal form reduces a matrix using both row AND column operations to
reach the identity block form. The rank is then read directly from the
size of the identity block.

\subsection{Definition of Normal Form}\label{definition-of-normal-form}

A non-zero matrix \(A\) of rank \(r\) can be reduced to one of these
forms using row and column operations:
\[
[I_r], \quad [I_r \mid 0], \quad
\begin{bmatrix}
I_r \\
0
\end{bmatrix}, \quad \begin{bmatrix} I_r & 0 \\ 0 & 0 \end{bmatrix}
\]
where \(I_r\) is the identity matrix of order \(r\). This is the
\textbf{normal form} (or \textbf{canonical form}).

The rank equals \(r\) (the order of \(I_r\) in the normal form).

\subsection{Difference from Echelon
Form}\label{difference-from-echelon-form}

\begin{longtable}[]{@{}
  >{\raggedright\arraybackslash}p{(\linewidth - 4\tabcolsep) * \real{0.2571}}
  >{\raggedright\arraybackslash}p{(\linewidth - 4\tabcolsep) * \real{0.3714}}
  >{\raggedright\arraybackslash}p{(\linewidth - 4\tabcolsep) * \real{0.3714}}@{}}
\toprule\noalign{}
\begin{minipage}[b]{\linewidth}\raggedright
Feature
\end{minipage} & \begin{minipage}[b]{\linewidth}\raggedright
Echelon Form
\end{minipage} & \begin{minipage}[b]{\linewidth}\raggedright
Normal Form
\end{minipage} \\
\midrule\noalign{}
\endhead
\bottomrule\noalign{}
\endlastfoot
Operations used & Row operations only & Row AND column operations \\
Result & Staircase pattern & Identity block \\
Off-diagonal entries & May be non-zero above pivots & All zero except
\(I_r\) \\
Rank reading & Count non-zero rows & Read order of \(I_r\) \\
\end{longtable}

\subsection{The Method}\label{the-method-3}

\begin{enumerate}
\def\labelenumi{\arabic{enumi}.}
\tightlist
\item
  Start with the given matrix.
\item
  Use row operations to get echelon form (upper triangular with pivots).
\item
  Scale each pivot to 1.
\item
  Use column operations to clear all non-zero entries above and beside
  each pivot.
\item
  If needed, swap columns to group the 1s together into \(I_r\).
\end{enumerate}

\subsection{Worked Example 1: 4x4 Matrix to Normal Form (PYQ
2019)}\label{worked-example-1-4x4-matrix-to-normal-form-pyq-2019}

\textbf{Problem}: Find the rank by reducing to normal form:
\[
A =
\begin{pmatrix}
3 & -2 & 0 & -7 \\
0 & 2 & 1 & -5 \\
1 & -2 & -2 & 1 \\
0 & 1 & 1 & -6
\end{pmatrix}
\]
\textbf{Solution}:

Swap \(R_1 \leftrightarrow R_3\) for leading 1. Then
\(R_3 \to R_3 - 3R_1\):
\[
\begin{pmatrix}
1 & -2 & -2 & 1 \\
0 & 2 & 1 & -5 \\
0 & 4 & 6 & -10 \\
0 & 1 & 1 & -6
\end{pmatrix}
\]
Swap \(R_2 \leftrightarrow R_4\). Then \(R_3 \to R_3 - 4R_2\) and
\(R_4 \to R_4 - 2R_2\):
\[
\begin{pmatrix}
1 & -2 & -2 & 1 \\
0 & 1 & 1 & -6 \\
0 & 0 & 2 & 14 \\
0 & 0 & -1 & 7
\end{pmatrix}
\]
Scale
\[
R_3 \to \frac{1}{2}R_3
\]
. Then \(R_4 \to R_4 + R_3\):
\[
\begin{pmatrix}
1 & -2 & -2 & 1 \\
0 & 1 & 1 & -6 \\
0 & 0 & 1 & 7 \\
0 & 0 & 0 & 14
\end{pmatrix}
\]
Scale
\[
R_4 \to \frac{1}{14}R_4
\]
. Now clear upward using column operations.

\(C_2 \to C_2 + 2C_1\), \(C_3 \to C_3 + 2C_1\), \(C_4 \to C_4 - C_1\):
\[
\begin{pmatrix}
1 & 0 & 0 & 0 \\
0 & 1 & 1 & -6 \\
0 & 0 & 1 & 7 \\
0 & 0 & 0 & 1
\end{pmatrix}
\]
\(C_3 \to C_3 - C_2\), \(C_4 \to C_4 + 6C_2\):
\[
\begin{pmatrix}
1 & 0 & 0 & 0 \\
0 & 1 & 0 & 0 \\
0 & 0 & 1 & 7 \\
0 & 0 & 0 & 1
\end{pmatrix}
\]
\(C_4 \to C_4 - 7C_3\):
\[
\begin{pmatrix}
1 & 0 & 0 & 0 \\
0 & 1 & 0 & 0 \\
0 & 0 & 1 & 0 \\
0 & 0 & 0 & 1
\end{pmatrix} = I_4
\]
Normal form is \(I_4\). Rank = 4.

\subsection{Worked Example 2: 4x4 to Normal Form (PYQ
2021)}\label{worked-example-2-4x4-to-normal-form-pyq-2021}

\textbf{Problem}: Reduce to canonical form and find rank:
\[
A =
\begin{pmatrix}
2 & 3 & -1 & -1 \\
1 & -1 & -2 & -4 \\
3 & 1 & 3 & -2 \\
6 & 3 & 0 & -7
\end{pmatrix}
\]
\textbf{Solution}: After row operations, the echelon form has 3 non-zero
rows. After column operations to clear off-diagonal entries:
\[
\begin{pmatrix}
1 & 0 & 0 & 0 \\
0 & 1 & 0 & 0 \\
0 & 0 & 1 & 0 \\
0 & 0 & 0 & 0
\end{pmatrix} = \begin{bmatrix} I_3 & 0 \\ 0 & 0 \end{bmatrix}
\]
Normal form is \([I_3 \mid 0]\)-type. Rank = 3.

\subsection{Exam Patterns}\label{exam-patterns-21}

\begin{itemize}
\tightlist
\item
  ``Reduce to normal/canonical form and find rank'' appears regularly.
\item
  Always show row operations first to reach echelon form. Then show
  column operations to reach the identity block.
\item
  Write each operation clearly (\(R_i \to ...\), \(C_j \to ...\)).
\item
  State the final rank explicitly.
\end{itemize}

\section{B07: System of Linear
Equations}\label{b07-system-of-linear-equations}

\begin{quote}
\textbf{Exam Weight}: Tier 1 \textbar{} \textbf{Appeared in}: 2017,
2018, 2019, 2020, 2021, 2023, 2024 (all 7 papers) + CT \textbar{}
\textbf{Typical Marks}: 5--7
\end{quote}

This topic has appeared in every single exam paper. Together with rank
and eigenvalues, it forms the trio of guaranteed exam topics.

\subsection{Matrix Form of a System}\label{matrix-form-of-a-system}

\subsubsection{Setting Up the System}\label{setting-up-the-system}

Any system of linear equations can be written in matrix form as:
\[
AX = B
\]
Here: - \(A\) is the coefficient matrix (contains the coefficients of
the variables) - \(X\) is the column vector of unknowns - \(B\) is the
column vector of constants

For example, the system:
\[
\begin{aligned}
x - 2y - 5z &= 2 \\
-x + y + 3z &= -2 \\
2x + y + z &= 3
\end{aligned}
\]
becomes:
\[
\begin{bmatrix}
1 & -2 & -5 \\
-1 & 1 & 3 \\
2 & 1 & 1
\end{bmatrix} \begin{bmatrix} x \\ y \\ z \end{bmatrix} = \begin{bmatrix} 2 \\ -2 \\ 3 \end{bmatrix}
\]
\subsubsection{The Augmented Matrix}\label{the-augmented-matrix}

We combine \(A\) and \(B\) into one matrix called the augmented matrix.
We write it as \([A : B]\) or \([A \lvert  B]\):
\[
[A : B] =
\begin{bmatrix}
1 & -2 & -5 & | & 2 \\
-1 & 1 & 3 & | & -2 \\
2 & 1 & 1 & | & 3
\end{bmatrix}
\]
The vertical line separates the coefficients from the constants. We
apply row operations to this augmented matrix to solve the system.

\subsection{Consistency Conditions}\label{consistency-conditions}

\subsubsection{The Rank Test}\label{the-rank-test}

A system \(AX = B\) is either consistent (has at least one solution) or
inconsistent (has no solution). The rank decides which case applies.

Let \(n\) = number of unknowns, \(\rho(A)\) = rank of coefficient
matrix, \(\rho[A:B]\) = rank of augmented matrix.

\begin{longtable}[]{@{}
  >{\raggedright\arraybackslash}p{(\linewidth - 2\tabcolsep) * \real{0.5789}}
  >{\raggedright\arraybackslash}p{(\linewidth - 2\tabcolsep) * \real{0.4211}}@{}}
\toprule\noalign{}
\begin{minipage}[b]{\linewidth}\raggedright
Condition
\end{minipage} & \begin{minipage}[b]{\linewidth}\raggedright
Result
\end{minipage} \\
\midrule\noalign{}
\endhead
\bottomrule\noalign{}
\endlastfoot
\(\rho(A) \neq \rho[A:B]\) & \textbf{No solution} (inconsistent) \\
\(\rho(A) = \rho[A:B] = n\) & \textbf{Unique solution} \\
\(\rho(A) = \rho[A:B] < n\) & \textbf{Infinitely many solutions} (with
\(n - \rho(A)\) free variables) \\
\end{longtable}

\subsubsection{Key Idea}\label{key-idea-1}

When we reduce \([A:B]\) to echelon form, a contradiction row looks like
\([0 \; 0 \; 0 \; \lvert  \; c]\) where \(c \neq 0\). This says
\(0 = c\), which is impossible. Such a row means the system has no
solution.

If no contradiction row appears, the system is consistent. Count the
non-zero rows in the coefficient part. That is \(\rho(A)\). If it equals
the number of variables, you get a unique solution. If it is less, you
get infinitely many solutions.

\subsection{Solving by Row Reduction}\label{solving-by-row-reduction}

\subsubsection{The Method}\label{the-method-4}

\begin{enumerate}
\def\labelenumi{\arabic{enumi}.}
\tightlist
\item
  Write the augmented matrix \([A:B]\).
\item
  Use elementary row operations to reach echelon form.
\item
  Check for contradiction rows.
\item
  Back-substitute from the bottom row upward.
\end{enumerate}

\subsubsection{Worked Example 1: Unique Solution (PYQ
2017)}\label{worked-example-1-unique-solution-pyq-2017}

\textbf{Problem}: Solve by elementary transformation:
\[
\begin{aligned}
2x_1 + 2x_2 + x_3 &= 2 \\
3x_1 + x_2 - 2x_3 &= 1 \\
4x_1 - 3x_2 - x_3 &= 3
\end{aligned}
\]
\textbf{Solution}:

Write the augmented matrix:
\[
\begin{bmatrix}
2 & 2 & 1 & | & 2 \\
3 & 1 & -2 & | & 1 \\
4 & -3 & -1 & | & 3
\end{bmatrix}
\]
Apply \(R_2 \to 2R_2 - 3R_1\) and \(R_3 \to R_3 - 2R_1\):
\[
\begin{bmatrix}
2 & 2 & 1 & | & 2 \\
0 & -4 & -7 & | & -4 \\
0 & -7 & -3 & | & -1
\end{bmatrix}
\]
Apply \(R_3 \to 4R_3 - 7R_2\):
\[
\begin{bmatrix}
2 & 2 & 1 & | & 2 \\
0 & 4 & 7 & | & 4 \\
0 & 0 & 37 & | & 24
\end{bmatrix}
\]
Back-substitute. From row 3:
\[
37x_3 = 24 \implies x_3 = \frac{24}{37}
\]
From row 2:
\[
4x_2 + 7 \cdot \frac{24}{37} = 4 \implies 4x_2 = 4 - \frac{168}{37} = \frac{-20}{37} \implies x_2 = -\frac{5}{37}
\]
From row 1:
\[
2x_1 + 2 \cdot \left(-\frac{5}{37}\right) + \frac{24}{37} = 2 \implies 2x_1 + \frac{14}{37} = 2 \implies x_1 = \frac{30}{37}
\]
The unique solution is:
\[
x_1 = \frac{30}{37}, \quad x_2 = -\frac{5}{37}, \quad x_3 = \frac{24}{37}
\]

\subsubsection{Worked Example 2: Infinite Solutions with Free Variables
(PYQ
2019)}\label{worked-example-2-infinite-solutions-with-free-variables-pyq-2019}

\textbf{Problem}: Discuss consistency and find the solution set:
\[
\begin{aligned}
x + 2y + 2z &= 1 \\
2x + y + z &= 2 \\
3x + 2y + 2z &= 3 \\
y + z &= 0
\end{aligned}
\]
\textbf{Solution}:

Write the augmented matrix:
\[
\begin{bmatrix}
1 & 2 & 2 & | & 1 \\
2 & 1 & 1 & | & 2 \\
3 & 2 & 2 & | & 3 \\
0 & 1 & 1 & | & 0
\end{bmatrix}
\]
Apply \(R_2 \to R_2 - 2R_1\) and \(R_3 \to R_3 - 3R_1\):
\[
\begin{bmatrix}
1 & 2 & 2 & | & 1 \\
0 & -3 & -3 & | & 0 \\
0 & -4 & -4 & | & 0 \\
0 & 1 & 1 & | & 0
\end{bmatrix}
\]
Scale:
\[
R_2 \to -\frac{1}{3}R_2
\]
and
\[
R_3 \to -\frac{1}{4}R_3:
\]
\[
\begin{bmatrix}
1 & 2 & 2 & | & 1 \\
0 & 1 & 1 & | & 0 \\
0 & 1 & 1 & | & 0 \\
0 & 1 & 1 & | & 0
\end{bmatrix}
\]
Apply \(R_3 \to R_3 - R_2\) and \(R_4 \to R_4 - R_2\):
\[
\begin{bmatrix}
1 & 2 & 2 & | & 1 \\
0 & 1 & 1 & | & 0 \\
0 & 0 & 0 & | & 0 \\
0 & 0 & 0 & | & 0
\end{bmatrix}
\]
\textbf{Check consistency}: \(\rho(A) = 2\) and \(\rho[A:B] = 2\). So
the system is consistent.

Since there are 3 variables and rank is 2, we have \(3 - 2 = 1\) free
variable.

From row 2: \(y + z = 0 \implies y = -z\).

From row 1: \(x + 2(-z) + 2z = 1 \implies x = 1\).

Let \(z = k\) where \(k \in \mathbb{R}\). The solution set is:
\[
x = 1, \quad y = -k, \quad z = k
\]

\subsubsection{Worked Example 3: Solving a System (PYQ
2023)}\label{worked-example-3-solving-a-system-pyq-2023}

\textbf{Problem}: Solve:
\[
\begin{aligned}
x + 2y - z &= 2 \\
2x + y + z &= 1 \\
x + 5y - 4z &= 5
\end{aligned}
\]
\textbf{Solution}:

Write the augmented matrix:
\[
\begin{bmatrix}
1 & 2 & -1 & | & 2 \\
2 & 1 & 1 & | & 1 \\
1 & 5 & -4 & | & 5
\end{bmatrix}
\]
Apply \(R_2 \to R_2 - 2R_1\) and \(R_3 \to R_3 - R_1\):
\[
\begin{bmatrix}
1 & 2 & -1 & | & 2 \\
0 & -3 & 3 & | & -3 \\
0 & 3 & -3 & | & 3
\end{bmatrix}
\]
Apply \(R_3 \to R_3 + R_2\):
\[
\begin{bmatrix}
1 & 2 & -1 & | & 2 \\
0 & -3 & 3 & | & -3 \\
0 & 0 & 0 & | & 0
\end{bmatrix}
\]
Rank is 2 with 3 variables. So there is 1 free variable.

From row 2: \(-3y + 3z = -3 \implies y - z = 1 \implies y = z + 1\).

From row 1:
\(x + 2(z + 1) - z = 2 \implies x + z + 2 = 2 \implies x = -z\).

Let \(z = t\). The general solution is:
\[
x = -t, \quad y = t + 1, \quad z = t
\]

\subsection{Finding a System with Parameters (PYQ
2018)}\label{finding-a-system-with-parameters-pyq-2018}

Sometimes the system contains a parameter \(\lambda\) and you must find
the values of \(\lambda\) that allow solutions.

\subsubsection{Worked Example 4: Parameter in the
System}\label{worked-example-4-parameter-in-the-system}

\textbf{Problem}: Find the value of \(\lambda\) so that the system has a
solution. Solve completely in each case.
\[
\begin{aligned}
x + y + z &= 1 \\
x + 2y + 4z &= \lambda \\
x + 4y + 10z &= \lambda^2
\end{aligned}
\]
\textbf{Solution}:

Write the augmented matrix:
\[
\begin{bmatrix}
1 & 1 & 1 & | & 1 \\
1 & 2 & 4 & | & \lambda \\
1 & 4 & 10 & | & \lambda^2
\end{bmatrix}
\]
Apply \(R_2 \to R_2 - R_1\) and \(R_3 \to R_3 - R_1\):
\[
\begin{bmatrix}
1 & 1 & 1 & | & 1 \\
0 & 1 & 3 & | & \lambda - 1 \\
0 & 3 & 9 & | & \lambda^2 - 1
\end{bmatrix}
\]
Apply \(R_3 \to R_3 - 3R_2\):
\[
\begin{bmatrix}
1 & 1 & 1 & | & 1 \\
0 & 1 & 3 & | & \lambda - 1 \\
0 & 0 & 0 & | & \lambda^2 - 3\lambda + 2
\end{bmatrix}
\]
For consistency, the last row must not be a contradiction. So we need:
\[
\lambda^2 - 3\lambda + 2 = 0 \implies (\lambda - 1)(\lambda - 2) = 0
\]
So \(\lambda = 1\) or \(\lambda = 2\).

\textbf{Case 1}: \(\lambda = 1\).

The augmented matrix becomes:
\[
\begin{bmatrix}
1 & 1 & 1 & | & 1 \\
0 & 1 & 3 & | & 0 \\
0 & 0 & 0 & | & 0
\end{bmatrix}
\]
Let \(z = t\). From row 2: \(y = -3t\). From row 1:
\(x = 1 - y - z = 1 + 3t - t = 1 + 2t\).
\[
x = 1 + 2t, \quad y = -3t, \quad z = t
\]
\textbf{Case 2}: \(\lambda = 2\).

The augmented matrix becomes:
\[
\begin{bmatrix}
1 & 1 & 1 & | & 1 \\
0 & 1 & 3 & | & 1 \\
0 & 0 & 0 & | & 0
\end{bmatrix}
\]
Let \(z = t\). From row 2: \(y = 1 - 3t\). From row 1:
\(x = 1 - y - z = 1 - (1 - 3t) - t = 2t\).
\[
x = 2t, \quad y = 1 - 3t, \quad z = t
\]

\subsection{Key Concepts to Remember}\label{key-concepts-to-remember}

\subsubsection{Consistent vs
Inconsistent}\label{consistent-vs-inconsistent}

A system is \textbf{consistent} if it has at least one solution. It is
\textbf{inconsistent} if it has no solution at all.

\subsubsection{Free Variables}\label{free-variables}

When the rank is less than the number of unknowns, the difference gives
the number of free variables. You assign parameters (\(t\), \(k\), etc.)
to these free variables and express all other variables in terms of
them.

\subsubsection{Choosing Free Variables}\label{choosing-free-variables}

In echelon form, the columns without pivots correspond to free
variables. The columns with pivots correspond to basic (determined)
variables.

\subsection{Exam Patterns}\label{exam-patterns-22}

\begin{itemize}
\tightlist
\item
  \textbf{This topic appears every year.} Usually as ``Solve the
  following system by elementary transformation'' or ``Discuss the
  consistency and find the solution.''
\item
  Marks range from 5 to 7 per question.
\item
  The most common format is 3 equations in 3 unknowns. Sometimes 4
  equations appear.
\item
  \textbf{Parameter problems} (with \(\lambda\)) appear frequently. See
  B08 for more on those.
\item
  Always show the augmented matrix and every row operation step. This is
  where marks are given.
\item
  Write the final answer clearly with a box or line. State the solution
  as \((x, y, z) = (a, b, c)\) or in parametric form.
\end{itemize}

\subsection{Must-Know Problems}\label{must-know-problems-1}

\subsubsection{Problem 5: Class Note
System}\label{problem-5-class-note-system}

\textbf{Problem}: Solve:
\[
\begin{aligned}
x - 2y - 5z &= 2 \\
-x + y + 3z &= -2 \\
2x + y + z &= 3
\end{aligned}
\]
\textbf{Solution}:

Augmented matrix:
\[
\begin{bmatrix}
1 & -2 & -5 & | & 2 \\
-1 & 1 & 3 & | & -2 \\
2 & 1 & 1 & | & 3
\end{bmatrix}
\]
Apply \(R_2 \to R_1 + R_2\) and \(R_3 \to 2R_1 - R_3\):
\[
\begin{bmatrix}
1 & -2 & -5 & | & 2 \\
0 & -1 & -2 & | & 0 \\
0 & -5 & -11 & | & 1
\end{bmatrix}
\]
Apply \(R_3 \to 5R_2 - R_3\):
\[
\begin{bmatrix}
1 & -2 & -5 & | & 2 \\
0 & -1 & -2 & | & 0 \\
0 & 0 & 1 & | & -1
\end{bmatrix}
\]
Back-substitute. From row 3: \(z = -1\).

From row 2: \(-y - 2(-1) = 0 \implies y = 2\).

From row 1: \(x - 2(2) - 5(-1) = 2 \implies x = 1\).
\[
(x, y, z) = (1, 2, -1)
\]

\subsubsection{Problem 6: Homogeneous System (PYQ
2018)}\label{problem-6-homogeneous-system-pyq-2018}

\textbf{Problem}: Solve:
\[
\begin{aligned}
2x_1 - x_2 + x_3 &= 0 \\
3x_1 + 2x_2 + x_3 &= 0 \\
x_1 - 3x_2 + 5x_3 &= 0
\end{aligned}
\]
\textbf{Solution}:

For a homogeneous system \(AX = 0\), the trivial solution
\[
x_1 = x_2 = x_3 = 0
\]
always exists.

Write the coefficient matrix and reduce:
\[
\begin{bmatrix}
1 & -3 & 5 \\
3 & 2 & 1 \\
2 & -1 & 1
\end{bmatrix}
\]
(We swapped \(R_1 \leftrightarrow R_3\) to get a leading 1.)

Apply \(R_2 \to R_2 - 3R_1\) and \(R_3 \to R_3 - 2R_1\):
\[
\begin{bmatrix}
1 & -3 & 5 \\
0 & 11 & -14 \\
0 & 5 & -9
\end{bmatrix}
\]
Apply \(R_3 \to 11R_3 - 5R_2\):
\[
\begin{bmatrix}
1 & -3 & 5 \\
0 & 11 & -14 \\
0 & 0 & -29
\end{bmatrix}
\]
Rank = 3 = number of variables. So the system has only the trivial
solution:
\[
x_1 = 0, \quad x_2 = 0, \quad x_3 = 0
\]

\subsubsection{Problem 7: Homogeneous System with Non-Trivial Solution
(PYQ
2019)}\label{problem-7-homogeneous-system-with-non-trivial-solution-pyq-2019}

\textbf{Problem}: Find the complete solution of:
\[
\begin{aligned}
x + y + z + w &= 0 \\
x + 3y - 2z + 4w &= 0 \\
2x + y - z - w &= 0
\end{aligned}
\]
\textbf{Solution}:

Coefficient matrix:
\[
\begin{bmatrix}
1 & 1 & 1 & 1 \\
1 & 3 & -2 & 4 \\
2 & 1 & -1 & -1
\end{bmatrix}
\]
Apply \(R_2 \to R_2 - R_1\) and \(R_3 \to R_3 - 2R_1\):
\[
\begin{bmatrix}
1 & 1 & 1 & 1 \\
0 & 2 & -3 & 3 \\
0 & -1 & -3 & -3
\end{bmatrix}
\]
Apply \(R_3 \to 2R_3 + R_2\):
\[
\begin{bmatrix}
1 & 1 & 1 & 1 \\
0 & 2 & -3 & 3 \\
0 & 0 & -9 & -3
\end{bmatrix}
\]
Rank = 3. Number of variables = 4. So there is \(4 - 3 = 1\) free
variable.

From row 3: \(-9z - 3w = 0 \implies w = -3z\).

From row 2: \(2y - 3z + 3(-3z) = 0 \implies 2y = 12z \implies y = 6z\).

From row 1: \(x + 6z + z + (-3z) = 0 \implies x = -4z\).

Let \(z = k\):
\[
x = -4k, \quad y = 6k, \quad z = k, \quad w = -3k
\]
The trivial solution is \(k = 0\). Any \(k \neq 0\) gives a non-trivial
solution.

\section{B08: Parameter Analysis for Systems of
Equations}\label{b08-parameter-analysis-for-systems-of-equations}

\begin{quote}
\textbf{Exam Weight}: Tier 1 \textbar{} \textbf{Appeared in}: 2018,
2020, 2021, 2024 + CT \textbar{} \textbf{Typical Marks}: 5--7
\end{quote}

Parameter analysis problems are a favourite of the examiners. The
classic question asks: ``For what values of \(\lambda\) and \(\mu\) does
the system have no solution, a unique solution, or infinitely many
solutions?'' One specific problem has appeared word-for-word in multiple
years.

\subsection{The General Approach}\label{the-general-approach}

\subsubsection{What the Problem Looks
Like}\label{what-the-problem-looks-like}

You are given a system of equations where one or more coefficients are
replaced by parameters (\(\lambda\), \(\mu\), \(k\), etc.). You must
find which values of these parameters lead to each of the three cases:

\begin{enumerate}
\def\labelenumi{\arabic{enumi}.}
\tightlist
\item
  \textbf{Unique solution}: The system behaves normally. One answer for
  each variable.
\item
  \textbf{No solution}: The system is contradictory. No set of values
  can satisfy all equations.
\item
  \textbf{Infinitely many solutions}: The system is underdetermined.
  Whole families of solutions exist.
\end{enumerate}

\subsubsection{The Method}\label{the-method-5}

\begin{enumerate}
\def\labelenumi{\arabic{enumi}.}
\tightlist
\item
  Write the augmented matrix \([A:B]\).
\item
  Row-reduce to echelon form. The parameter will survive in the last
  row.
\item
  The last row will look like
\end{enumerate}
\[
[0 \; 0 \; f(\lambda) \; \lvert  \; g(\lambda, \mu)].
\]
\begin{enumerate}
\def\labelenumi{\arabic{enumi}.}
\setcounter{enumi}{3}
\tightlist
\item
  Analyze three cases based on \(f(\lambda)\) and \(g(\lambda, \mu)\).
\end{enumerate}

\subsubsection{The Decision Table}\label{the-decision-table}

After row reduction, the critical row has the form:
\[
[0 \quad 0 \quad f(\lambda) \quad | \quad g(\lambda, \mu)]
\]
\begin{longtable}[]{@{}
  >{\raggedright\arraybackslash}p{(\linewidth - 2\tabcolsep) * \real{0.5789}}
  >{\raggedright\arraybackslash}p{(\linewidth - 2\tabcolsep) * \real{0.4211}}@{}}
\toprule\noalign{}
\begin{minipage}[b]{\linewidth}\raggedright
Condition
\end{minipage} & \begin{minipage}[b]{\linewidth}\raggedright
Result
\end{minipage} \\
\midrule\noalign{}
\endhead
\bottomrule\noalign{}
\endlastfoot
\(f(\lambda) \neq 0\) & Unique solution (for any \(\mu\)) \\
\(f(\lambda) = 0\) and \(g(\lambda, \mu) \neq 0\) & No solution
(contradiction: \(0 = c\)) \\
\(f(\lambda) = 0\) and \(g(\lambda, \mu) = 0\) & Infinitely many
solutions \\
\end{longtable}

\subsection{Must-Know Problem 1: The Classic (PYQ 2020, 2024 ---
Repeated)}\label{must-know-problem-1-the-classic-pyq-2020-2024-repeated}

\textbf{Problem}: Investigate for what values of \(\lambda\) and \(\mu\)
the system has (i) no solution, (ii) unique solution, (iii) infinitely
many solutions:
\[
\begin{aligned}
x + y + z &= 6 \\
x + 2y + 3z &= 10 \\
x + 2y + \lambda z &= \mu
\end{aligned}
\]
This exact problem appeared in 2020 (7 marks) and 2024 (5 marks).

\textbf{Solution}:

Write the augmented matrix:
\[
\begin{bmatrix}
1 & 1 & 1 & | & 6 \\
1 & 2 & 3 & | & 10 \\
1 & 2 & \lambda & | & \mu
\end{bmatrix}
\]
Apply \(R_2 \to R_2 - R_1\) and \(R_3 \to R_3 - R_1\):
\[
\begin{bmatrix}
1 & 1 & 1 & | & 6 \\
0 & 1 & 2 & | & 4 \\
0 & 1 & \lambda - 1 & | & \mu - 6
\end{bmatrix}
\]
Apply \(R_3 \to R_3 - R_2\):
\[
\begin{bmatrix}
1 & 1 & 1 & | & 6 \\
0 & 1 & 2 & | & 4 \\
0 & 0 & \lambda - 3 & | & \mu - 10
\end{bmatrix}
\]
Now analyze the last row:

\textbf{(i) No solution}: Need \(\lambda - 3 = 0\) and
\(\mu - 10 \neq 0\):
\[
\lambda = 3 \quad \text{and} \quad \mu \neq 10
\]
The last row becomes
\[
[0 \; 0 \; 0 \; \lvert  \; \mu - 10]
\]
, which says \(0 = \mu - 10\). This is a contradiction since
\(\mu \neq 10\).

\textbf{(ii) Unique solution}: Need \(\lambda - 3 \neq 0\):
\[
\lambda \neq 3 \quad \text{(for any value of } \mu\text{)}
\]
All three rows have pivots. Rank = 3 = number of unknowns.

\textbf{(iii) Infinitely many solutions}: Need \(\lambda - 3 = 0\) and
\(\mu - 10 = 0\):
\[
\lambda = 3 \quad \text{and} \quad \mu = 10
\]
The last row becomes all zeros. Rank = 2 \textless{} 3 unknowns. One
free variable.

\subsection{Must-Know Problem 2: Single Parameter (PYQ
2021)}\label{must-know-problem-2-single-parameter-pyq-2021}

\textbf{Problem}: For what values of \(\lambda\) does the system have
(i) no solution, (ii) unique solution, (iii) infinitely many solutions?
\[
\begin{aligned}
\lambda x + y + z &= 1 \\
x + \lambda y + z &= \lambda \\
x + y + \lambda z &= \lambda^2
\end{aligned}
\]
\textbf{Solution}:

This system has \(\lambda\) in the coefficient matrix AND the constant
vector. We use the determinant approach.

The coefficient matrix determinant is:
\[
D =
\begin{vmatrix}
\lambda & 1 & 1 \\
1 & \lambda & 1 \\
1 & 1 & \lambda
\end{vmatrix}
\]
Add all three rows to \(R_1\):
\[
R_1 \to [\lambda + 2, \lambda + 2, \lambda + 2]
\]
Factor out \((\lambda + 2)\):
\[
D = (\lambda + 2)
\begin{vmatrix}
1 & 1 & 1 \\
1 & \lambda & 1 \\
1 & 1 & \lambda
\end{vmatrix}
\]
Apply \(C_2 \to C_2 - C_1\) and \(C_3 \to C_3 - C_1\):
\[
D = (\lambda + 2)
\begin{vmatrix}
1 & 0 & 0 \\
1 & \lambda - 1 & 0 \\
1 & 0 & \lambda - 1
\end{vmatrix} = (\lambda + 2)(\lambda - 1)^2
\]
\textbf{Analysis}:

\textbf{(ii) Unique solution}: \(D \neq 0\), which means
\(\lambda \neq 1\) and \(\lambda \neq -2\).

\textbf{(iii) Infinitely many solutions}: When \(\lambda = 1\), all
three equations become \(x + y + z = 1\). They are identical. Infinitely
many solutions.

\textbf{(i) No solution}: When \(\lambda = -2\), the system becomes:
\[
\begin{aligned}
-2x + y + z &= 1 \\
x - 2y + z &= -2 \\
x + y - 2z &= 4
\end{aligned}
\]
Add all three equations: \(0 = 3\). Contradiction. No solution.

\subsection{Must-Know Problem 3: CT
Problem}\label{must-know-problem-3-ct-problem}

\textbf{Problem}: Find the values of \(k\) for which the system has (i)
unique solution, (ii) no solution, (iii) infinitely many solutions:
\[
\begin{aligned}
x + y + z &= 1 \\
x + 2y + 3z &= k \\
x + 5y + 9z &= k^2
\end{aligned}
\]
\textbf{Solution}:

Augmented matrix:
\[
\begin{bmatrix}
1 & 1 & 1 & | & 1 \\
1 & 2 & 3 & | & k \\
1 & 5 & 9 & | & k^2
\end{bmatrix}
\]
Apply \(R_2 \to R_2 - R_1\) and \(R_3 \to R_3 - R_1\):
\[
\begin{bmatrix}
1 & 1 & 1 & | & 1 \\
0 & 1 & 2 & | & k - 1 \\
0 & 4 & 8 & | & k^2 - 1
\end{bmatrix}
\]
Apply \(R_3 \to R_3 - 4R_2\):
\[
\begin{bmatrix}
1 & 1 & 1 & | & 1 \\
0 & 1 & 2 & | & k - 1 \\
0 & 0 & 0 & | & k^2 - 4k + 3
\end{bmatrix}
\]
Factor the last entry:
\[
k^2 - 4k + 3 = (k - 1)(k - 3).
\]
\textbf{(i) Unique solution}: The coefficient part has rank 2 (not 3)
because the last row of coefficients is all zeros. So this system can
NEVER have a unique solution.

Wait. Let me re-examine. The coefficient matrix after reduction is:
\[
\begin{bmatrix}
1 & 1 & 1 \\
0 & 1 & 2 \\
0 & 0 & 0
\end{bmatrix}
\]
Rank of coefficient matrix = 2. Since rank \textless{} 3 (number of
unknowns), a unique solution is impossible regardless of \(k\).

\textbf{(ii) No solution}: \((k - 1)(k - 3) \neq 0 \implies k \neq 1\)
and \(k \neq 3\).

\textbf{(iii) Infinitely many solutions}:
\((k - 1)(k - 3) = 0 \implies k = 1\) or \(k = 3\).

When \(k = 1\): Let \(z = t\). From row 2: \(y = -2t\). From row 1:
\(x = 1 - y - z = 1 + 2t - t = 1 + t\).
\[
x = 1 + t, \quad y = -2t, \quad z = t
\]
When \(k = 3\): Let \(z = t\). From row 2: \(y = 2 - 2t\). From row 1:
\(x = 1 - y - z = 1 - (2 - 2t) - t = -1 + t\).
\[
x = -1 + t, \quad y = 2 - 2t, \quad z = t
\]

\subsection{Key Observations}\label{key-observations}

\subsubsection{When Two Parameters
Appear}\label{when-two-parameters-appear}

If the system has both \(\lambda\) (in the coefficient matrix) and
\(\mu\) (in the constant vector), you get a 2D analysis. The answer
typically looks like: - Unique: \(\lambda \neq\) some value, \(\mu\) =
anything - No solution: \(\lambda =\) that value, \(\mu \neq\) some
other value - Infinite: \(\lambda\) and \(\mu\) both equal specific
values

\subsubsection{When One Parameter
Appears}\label{when-one-parameter-appears}

If \(\lambda\) appears only in the coefficient matrix, the analysis is
simpler. You compute the determinant and find where it is zero. At those
zero-points, check if the system is consistent or not.

\subsubsection{The Determinant Shortcut}\label{the-determinant-shortcut}

For a 3x3 system with parameter \(\lambda\): - Compute
\[
D = \lvert A \rvert
\]
as a function of \(\lambda\). - \(D \neq 0\) means unique solution. -
\(D = 0\) means either no solution or infinite solutions. You must check
further.

\subsection{Exam Patterns}\label{exam-patterns-23}

\begin{itemize}
\tightlist
\item
  \textbf{This appears about every other year.} It has shown up in 2018,
  2020, 2021, 2024, and CTs.
\item
  The \(x+y+z=6\), \(x+2y+3z=10\), \(x+2y+\lambda z=\mu\) problem is the
  most repeated. Memorize it.
\item
  Always show the complete row reduction. Marks are given for the steps.
\item
  State all three cases clearly: (i), (ii), (iii) with the exact
  conditions on \(\lambda\) and \(\mu\).
\item
  If the problem asks to ``solve completely,'' you must also give the
  general solution in the infinite-solutions case.
\end{itemize}

\section{B09: Homogeneous Systems of
Equations}\label{b09-homogeneous-systems-of-equations}

\begin{quote}
\textbf{Exam Weight}: Tier 2 \textbar{} \textbf{Appeared in}: 2018, 2019
\textbar{} \textbf{Typical Marks}: 5--6
\end{quote}

A homogeneous system always has the trivial solution. The real question
is whether non-trivial solutions exist.

\subsection{What is a Homogeneous
System?}\label{what-is-a-homogeneous-system}

A system \(AX = 0\) where the constant vector is all zeros:
\[
\begin{aligned}
a_{11}x_1 + a_{12}x_2 + \dots + a_{1n}x_n &= 0 \\
a_{21}x_1 + a_{22}x_2 + \dots + a_{2n}x_n &= 0 \\
\vdots \\
a_{m1}x_1 + a_{m2}x_2 + \dots + a_{mn}x_n &= 0
\end{aligned}
\]
The trivial solution
\[
x_1 = x_2 = \dots = x_n = 0
\]
always exists. It always satisfies every equation.

\subsection{When Do Non-Trivial Solutions
Exist?}\label{when-do-non-trivial-solutions-exist}

Let \(\rho(A)\) be the rank of the coefficient matrix and \(n\) be the
number of unknowns.

\begin{longtable}[]{@{}
  >{\raggedright\arraybackslash}p{(\linewidth - 2\tabcolsep) * \real{0.5789}}
  >{\raggedright\arraybackslash}p{(\linewidth - 2\tabcolsep) * \real{0.4211}}@{}}
\toprule\noalign{}
\begin{minipage}[b]{\linewidth}\raggedright
Condition
\end{minipage} & \begin{minipage}[b]{\linewidth}\raggedright
Result
\end{minipage} \\
\midrule\noalign{}
\endhead
\bottomrule\noalign{}
\endlastfoot
\(\rho(A) = n\) & Only the trivial solution (zero solution) \\
\(\rho(A) < n\) & Non-trivial solutions exist. The solution space has
dimension \(n - \rho(A)\) \\
\end{longtable}

For a square system (\(m = n\)): - \(\lvert A \rvert \neq 0\) means only
the trivial solution. -
\[
\lvert A \rvert = 0
\]
means non-trivial solutions exist.

\subsection{The Solution Space}\label{the-solution-space}

The set of all solutions to \(AX = 0\) forms a vector subspace called
the \textbf{null space} of \(A\). Its dimension is the \textbf{nullity}
= \(n - \rho(A)\).

When non-trivial solutions exist, you express the solution in terms of
free variables (parameters). Each free variable gives one dimension.

\subsection{Worked Example 1: Only Trivial Solution (PYQ
2018)}\label{worked-example-1-only-trivial-solution-pyq-2018}

\textbf{Problem}: Solve the homogeneous system:
\[
\begin{aligned}
2x_1 - x_2 + x_3 &= 0 \\
3x_1 + 2x_2 + x_3 &= 0 \\
x_1 - 3x_2 + 5x_3 &= 0
\end{aligned}
\]
\textbf{Solution}:

Swap \(R_1 \leftrightarrow R_3\):
\[
\begin{bmatrix}
1 & -3 & 5 \\
3 & 2 & 1 \\
2 & -1 & 1
\end{bmatrix}
\]
Apply \(R_2 \to R_2 - 3R_1\) and \(R_3 \to R_3 - 2R_1\):
\[
\begin{bmatrix}
1 & -3 & 5 \\
0 & 11 & -14 \\
0 & 5 & -9
\end{bmatrix}
\]
Apply \(R_3 \to 11R_3 - 5R_2\):
\[
\begin{bmatrix}
1 & -3 & 5 \\
0 & 11 & -14 \\
0 & 0 & -29
\end{bmatrix}
\]
Rank = 3 = number of unknowns. Only the trivial solution:
\[
x_1 = 0, \quad x_2 = 0, \quad x_3 = 0
\]

\subsection{Worked Example 2: Non-Trivial Solutions (PYQ
2019)}\label{worked-example-2-non-trivial-solutions-pyq-2019}

\textbf{Problem}: Find the complete solution of:
\[
\begin{aligned}
x + y + z + w &= 0 \\
x + 3y - 2z + 4w &= 0 \\
2x + y - z - w &= 0
\end{aligned}
\]
\textbf{Solution}:

After row reduction:
\[
\begin{bmatrix}
1 & 1 & 1 & 1 \\
0 & 1 & 3 & 3 \\
0 & 0 & 3 & 1
\end{bmatrix}
\]
Rank = 3. Variables = 4. So \(4 - 3 = 1\) free variable.

From row 3: \(3z + w = 0 \implies w = -3z\).

From row 2: \(y + 3z + 3(-3z) = 0 \implies y = 6z\).

From row 1: \(x + 6z + z - 3z = 0 \implies x = -4z\).

Let \(z = k\). The general solution is:
\[
x = -4k, \quad y = 6k, \quad z = k, \quad w = -3k
\]
In vector form:
\[
\begin{bmatrix}
x \\
y \\
z \\
w
\end{bmatrix} = k \begin{bmatrix} -4 \\ 6 \\ 1 \\ -3 \end{bmatrix}
\]
The null space is one-dimensional, spanned by \((-4, 6, 1, -3)^T\).

\subsection{Exam Patterns}\label{exam-patterns-24}

\begin{itemize}
\tightlist
\item
  Homogeneous systems appear as standalone problems or as part of
  eigenvalue problems (when solving \((A - \lambda I)X = 0\)).
\item
  The key test: compare rank with number of unknowns.
\item
  If the problem says ``find the trivial solution,'' the mathematical
  intent is usually to find ALL solutions.
\item
  Express solutions in vector form when possible.
\end{itemize}

\section{B10: Eigenvalues and
Eigenvectors}\label{b10-eigenvalues-and-eigenvectors}

\begin{quote}
\textbf{Exam Weight}: Tier 1 \textbar{} \textbf{Appeared in}: 2017,
2018, 2019, 2020, 2021, 2023, 2024 (all 7 papers) \textbar{}
\textbf{Typical Marks}: 5--12
\end{quote}

Eigenvalues and eigenvectors appear in every exam. The same 3x3 matrix
has even been reused across years. Master the method and you guarantee
yourself 5-12 marks.

\subsection{Core Definitions}\label{core-definitions}

\subsubsection{Eigenvalue and
Eigenvector}\label{eigenvalue-and-eigenvector}

Let \(A\) be a square matrix of order \(n\). If a non-zero vector \(X\)
and a scalar \(\lambda\) satisfy:
\[
AX = \lambda X
\]
then \(\lambda\) is an \textbf{eigenvalue} of \(A\) and \(X\) is the
corresponding \textbf{eigenvector}.

The eigenvector does not change direction under the matrix
transformation. It only gets scaled by the factor \(\lambda\).

\subsubsection{Characteristic Matrix}\label{characteristic-matrix}

The matrix \(A - \lambda I\) is called the characteristic matrix of
\(A\). Here \(I\) is the identity matrix of the same size as \(A\).

\subsubsection{Characteristic
Polynomial}\label{characteristic-polynomial}

The determinant \(\lvert A - \lambda I \rvert\) is a polynomial in
\(\lambda\). This polynomial is called the characteristic polynomial.

\subsubsection{Characteristic Equation}\label{characteristic-equation}

Setting the characteristic polynomial equal to zero gives the
characteristic equation:
\[
|A - \lambda I| = 0
\]
The roots of this equation are the eigenvalues.

\subsection{The Working Rule}\label{the-working-rule}

Follow these three steps for any eigenvalue/eigenvector problem:

\textbf{Step 1: Form the characteristic equation.} Write down
\(A - \lambda I\) and compute its determinant. Set it equal to zero.

\textbf{Step 2: Solve the polynomial.} Find all roots
\[
\lambda_1, \lambda_2, \dots, \lambda_n
\]
. These are the eigenvalues.

\textbf{Step 3: Find eigenvectors.} For each \(\lambda_i\), substitute
it into
\[
(A - \lambda_i I)X = 0
\]
. Solve this homogeneous system by row reduction. The non-zero solutions
are the eigenvectors.

\subsection{Properties of Eigenvalues}\label{properties-of-eigenvalues}

These properties help you check your answers and sometimes shortcut the
calculation:

\begin{enumerate}
\def\labelenumi{\arabic{enumi}.}
\tightlist
\item
  The sum of all eigenvalues equals the trace (sum of diagonal elements)
  of \(A\).
\end{enumerate}
\[
\lambda_1 + \lambda_2 + \dots + \lambda_n = \text{tr}(A) = a_{11} + a_{22} + \dots + a_{nn}
\]
\begin{enumerate}
\def\labelenumi{\arabic{enumi}.}
\setcounter{enumi}{1}
\tightlist
\item
  The product of all eigenvalues equals the determinant of \(A\).
\end{enumerate}
\[
\lambda_1 \cdot \lambda_2 \cdots \lambda_n = |A|
\]
\begin{enumerate}
\def\labelenumi{\arabic{enumi}.}
\setcounter{enumi}{2}
\item
  For a triangular matrix (upper or lower), the eigenvalues are the
  diagonal entries.
\item
  If \(\lambda\) is an eigenvalue of \(A\), then \(\lambda^k\) is an
  eigenvalue of \(A^k\).
\item
  If \(A\) is invertible and \(\lambda\) is an eigenvalue, then
  \(\frac{1}{\lambda}\) is an eigenvalue of \(A^{-1}\).
\item
  An \(n \times n\) matrix has exactly \(n\) eigenvalues (counting
  multiplicity). They may be real or complex.
\item
  Eigenvectors corresponding to distinct eigenvalues are always linearly
  independent.
\end{enumerate}

\subsection{Worked Example 1: 2x2 Matrix (PYQ
2021)}\label{worked-example-1-2x2-matrix-pyq-2021}

\textbf{Problem}: Find the eigenvalues and eigenvectors of:
\[
A =
\begin{bmatrix}
2 & 3 \\
1 & 4
\end{bmatrix}
\]
\textbf{Solution}:

\textbf{Step 1}: The characteristic equation is
\[
\lvert A - \lambda I \rvert = 0:
\]
\[
\begin{vmatrix}
2 - \lambda & 3 \\
1 & 4 - \lambda
\end{vmatrix} = 0
\]
\[
(2 - \lambda)(4 - \lambda) - 3 = 0 \implies \lambda^2 - 6\lambda + 5 = 0
\]
\textbf{Step 2}: Factor: \((\lambda - 1)(\lambda - 5) = 0\). So
\[
\lambda_1 = 1
\]
and
\[
\lambda_2 = 5.
\]
\textbf{Check}: Trace = \(2 + 4 = 6 = 1 + 5\). Determinant =
\(8 - 3 = 5 = 1 \times 5\). Correct.

\textbf{Step 3a}: For \(\lambda = 1\), solve \((A - I)X = 0\):
\[
\begin{bmatrix}
1 & 3 \\
1 & 3
\end{bmatrix} \begin{bmatrix} x \\ y \end{bmatrix} = \begin{bmatrix} 0 \\ 0 \end{bmatrix}
\]
This gives \(x + 3y = 0 \implies x = -3y\). Let \(y = 1\):
\[
X_1 =
\begin{bmatrix}
-3 \\
1
\end{bmatrix}
\]
\textbf{Step 3b}: For \(\lambda = 5\), solve \((A - 5I)X = 0\):
\[
\begin{bmatrix}
-3 & 3 \\
1 & -1
\end{bmatrix} \begin{bmatrix} x \\ y \end{bmatrix} = \begin{bmatrix} 0 \\ 0 \end{bmatrix}
\]
This gives \(x - y = 0 \implies x = y\). Let \(y = 1\):
\[
X_2 =
\begin{bmatrix}
1 \\
1
\end{bmatrix}
\]

\subsection{Worked Example 2: 3x3 Matrix (PYQ 2017, 2018 ---
Repeated)}\label{worked-example-2-3x3-matrix-pyq-2017-2018-repeated}

\textbf{Problem}: Find the eigenvalues and eigenvectors of:
\[
A =
\begin{bmatrix}
2 & 2 & 1 \\
1 & 3 & 1 \\
1 & 2 & 2
\end{bmatrix}
\]
This exact matrix appeared in both 2017 (12 marks) and 2018 (6 marks).

\textbf{Solution}:

\textbf{Step 1}: Compute
\[
\lvert A - \lambda I \rvert = 0:
\]
\[
\begin{vmatrix}
2-\lambda & 2 & 1 \\
1 & 3-\lambda & 1 \\
1 & 2 & 2-\lambda
\end{vmatrix} = 0
\]
Expand the determinant:
\[
(2-\lambda)[(3-\lambda)(2-\lambda) - 2] - 2[(2-\lambda) - 1] + 1[2 - (3-\lambda)] = 0
\]
\[
(2-\lambda)[\lambda^2 - 5\lambda + 4] - 2[1-\lambda] + [\lambda - 1] = 0
\]
After simplification:
\[
\lambda^3 - 7\lambda^2 + 11\lambda - 5 = 0
\]
\textbf{Step 2}: Test \(\lambda = 1\): \(1 - 7 + 11 - 5 = 0\). Yes, it
is a root.

Divide by \((\lambda - 1)\):
\[
(\lambda - 1)(\lambda^2 - 6\lambda + 5) = 0 \implies (\lambda - 1)(\lambda - 1)(\lambda - 5) = 0.
\]
Eigenvalues: \(\lambda = 1, 1, 5\).

\textbf{Check}: Trace = \(2 + 3 + 2 = 7 = 1 + 1 + 5\). Correct.

\textbf{Step 3a}: For \(\lambda = 5\), solve \((A - 5I)X = 0\):
\[
\begin{bmatrix}
-3 & 2 & 1 \\
1 & -2 & 1 \\
1 & 2 & -3
\end{bmatrix}
\]
Swap \(R_1 \leftrightarrow R_2\), then \(R_2 \to R_2 + 3R_1\),
\(R_3 \to R_3 - R_1\):
\[
\begin{bmatrix}
1 & -2 & 1 \\
0 & -4 & 4 \\
0 & 4 & -4
\end{bmatrix}
\]
Apply \(R_3 \to R_3 + R_2\): row 3 becomes all zeros. From row 2:
\(y = z\). From row 1: \(x = 2y - z = z\).

So \(x = y = z\). Let \(z = 1\):
\[
X_1 =
\begin{bmatrix}
1 \\
1 \\
1
\end{bmatrix}
\]
\textbf{Step 3b}: For \(\lambda = 1\) (repeated), solve
\((A - I)X = 0\):
\[
\begin{bmatrix}
1 & 2 & 1 \\
1 & 2 & 1 \\
1 & 2 & 1
\end{bmatrix}
\]
All three rows are identical. This gives one equation:
\(x + 2y + z = 0 \implies x = -2y - z\).

Two free variables (\(y\) and \(z\)). Let \(y = s\), \(z = t\):
\[
X = s
\begin{bmatrix}
-2 \\
1 \\
0
\end{bmatrix} + t \begin{bmatrix} -1 \\ 0 \\ 1 \end{bmatrix}
\]
Two independent eigenvectors:
\[
X_2 =
\begin{bmatrix}
-2 \\
1 \\
0
\end{bmatrix}, \quad X_3 = \begin{bmatrix} -1 \\ 0 \\ 1 \end{bmatrix}
\]

\subsection{Worked Example 3: 3x3 Matrix (PYQ
2020)}\label{worked-example-3-3x3-matrix-pyq-2020}

\textbf{Problem}: Find eigenvalues and eigenvectors of:
\[
A =
\begin{bmatrix}
1 & 1 & -2 \\
-1 & 2 & 1 \\
0 & 1 & -1
\end{bmatrix}
\]
\textbf{Solution}:

\textbf{Step 1-2}: The characteristic equation simplifies to:
\[
\lambda^3 - 2\lambda^2 - \lambda + 2 = 0 \implies (\lambda - 1)(\lambda + 1)(\lambda - 2) = 0
\]
Eigenvalues: \(\lambda = 1, -1, 2\).

\textbf{Step 3a}: For \(\lambda = 1\):
\[
\begin{bmatrix}
0 & 1 & -2 \\
-1 & 1 & 1 \\
0 & 1 & -2
\end{bmatrix}
\]
From row 1: \(y = 2z\). From row 2: \(x = y + z = 3z\). Let \(z = 1\):
\[
X_1 =
\begin{bmatrix}
3 \\
2 \\
1
\end{bmatrix}
\]
\textbf{Step 3b}: For \(\lambda = -1\):
\[
\begin{bmatrix}
2 & 1 & -2 \\
-1 & 3 & 1 \\
0 & 1 & 0
\end{bmatrix}
\]
From row 3: \(y = 0\). From row 1: \(2x - 2z = 0 \implies x = z\). Let
\(z = 1\):
\[
X_2 =
\begin{bmatrix}
1 \\
0 \\
1
\end{bmatrix}
\]
\textbf{Step 3c}: For \(\lambda = 2\):
\[
\begin{bmatrix}
-1 & 1 & -2 \\
-1 & 0 & 1 \\
0 & 1 & -3
\end{bmatrix}
\]
From row 2: \(x = z\). From row 3: \(y = 3z\). Let \(z = 1\):
\[
X_3 =
\begin{bmatrix}
1 \\
3 \\
1
\end{bmatrix}
\]

\subsection{Worked Example 4: 3x3 Matrix (PYQ
2023)}\label{worked-example-4-3x3-matrix-pyq-2023}

\textbf{Problem}: Find eigenvalues and eigenvectors of:
\[
A =
\begin{bmatrix}
2 & 1 & 1 \\
-1 & 2 & -1 \\
1 & -1 & 2
\end{bmatrix}
\]
\textbf{Solution}:

The characteristic equation gives:
\[
(2 - \lambda)(\lambda - 1)(\lambda - 3) = 0
\]
Eigenvalues: \(\lambda = 1, 2, 3\).

For \(\lambda = 1\): Solving \((A - I)X = 0\) gives \(y = 0\),
\(x = -z\). Eigenvector:
\[
X_1 =
\begin{bmatrix}
-1 \\
0 \\
1
\end{bmatrix}
\]
For \(\lambda = 2\): Solving \((A - 2I)X = 0\) gives \(x = -z\),
\(y = -z\). Eigenvector:
\[
X_2 =
\begin{bmatrix}
-1 \\
-1 \\
1
\end{bmatrix}
\]
For \(\lambda = 3\): Solving \((A - 3I)X = 0\) gives \(x = 0\),
\(y = -z\). Eigenvector:
\[
X_3 =
\begin{bmatrix}
0 \\
-1 \\
1
\end{bmatrix}
\]

\subsection{Special Cases}\label{special-cases}

\subsubsection{Triangular Matrices}\label{triangular-matrices}

For any upper or lower triangular matrix, the eigenvalues are the
diagonal entries. No computation needed.

\textbf{Example} (PYQ 2017): The eigenvalues of:
\[
\begin{pmatrix}
1 & 2 & 3 \\
0 & -4 & 2 \\
0 & 0 & 7
\end{pmatrix}
\]
are simply \(\lambda = 1, -4, 7\).

\subsubsection{Repeated Eigenvalues}\label{repeated-eigenvalues}

When an eigenvalue has multiplicity \(m\) (it appears \(m\) times as a
root), the number of independent eigenvectors can range from 1 to \(m\).
If you get fewer than \(m\) independent eigenvectors, the matrix cannot
be diagonalized. This is important for the diagonalization topic (see
B13).

\subsubsection{Complex Eigenvalues}\label{complex-eigenvalues}

If the characteristic equation has complex roots, the eigenvectors will
also have complex entries. Complex eigenvalues always come in conjugate
pairs for real matrices. See the class notes for the worked example with
\[
\lambda = \frac{3 \pm i\sqrt{3}}{2}.
\]

\subsection{Exam Patterns}\label{exam-patterns-25}

\begin{itemize}
\tightlist
\item
  \textbf{Every paper asks eigenvalues.} The phrasing is always ``Find
  the eigenvalues and eigenvectors of\ldots{}''
\item
  Usually a 3x3 matrix. Marks range from 5 to 12.
\item
  The same matrix has appeared multiple years (the
  \([2,2,1; 1,3,1; 1,2,2]\) matrix appeared in 2017 and 2018).
\item
  \textbf{Always verify} using trace and determinant. The trace must
  equal the sum of eigenvalues. The determinant must equal their
  product.
\item
  For 3x3 matrices, the characteristic polynomial is a cubic. Try
  integer values (0, 1, 2, -1) first. One root is usually an integer,
  and you can factor out.
\item
  When solving \((A - \lambda I)X = 0\), you solve a homogeneous system.
  Use echelon form and express variables in terms of free variables.
\item
  Write eigenvectors with integer entries. Choose the free variable
  value to avoid fractions.
\end{itemize}

\subsection{Eigenvalue Quick-Check
Table}\label{eigenvalue-quick-check-table}

For a 3x3 matrix with eigenvalues
\[
\lambda_1, \lambda_2, \lambda_3:
\]
\begin{longtable}[]{@{}ll@{}}
\toprule\noalign{}
Property & Formula \\
\midrule\noalign{}
\endhead
\bottomrule\noalign{}
\endlastfoot
Sum & \\
\end{longtable}
\[
\lambda_1 + \lambda_2 + \lambda_3 = a_{11} + a_{22} + a_{33}
\]
\hfill\break
Product \textbar{}
\[
\lambda_1 \cdot \lambda_2 \cdot \lambda_3 = \lvert A \rvert
\]
\hfill\break
Sum of products of pairs \textbar{}
\[
\lambda_1 \lambda_2 + \lambda_1 \lambda_3 + \lambda_2 \lambda_3 =
\]
sum of \(2 \times 2\) cofactors \textbar{}

\section{B11: Cayley-Hamilton
Theorem}\label{b11-cayley-hamilton-theorem}

\begin{quote}
\textbf{Exam Weight}: Tier 2 \textbar{} \textbf{Appeared in}: 2018,
2019, 2024 \textbar{} \textbf{Typical Marks}: 5--6
\end{quote}

The Cayley-Hamilton theorem connects a matrix to its own characteristic
equation. The exam typically asks you to state the theorem, verify it
for a given matrix, and then use it to find the inverse.

\subsection{Statement}\label{statement-8}

Every square matrix satisfies its own characteristic equation.

If the characteristic equation of matrix \(A\) of order \(n\) is:
\[
\lambda^n + c_{n-1}\lambda^{n-1} + \dots + c_1\lambda + c_0 = 0
\]
then replacing \(\lambda\) with \(A\) and the constant term with
\(c_0 I\) gives:
\[
A^n + c_{n-1}A^{n-1} + \dots + c_1 A + c_0 I = O
\]
where \(O\) is the zero matrix.

\subsubsection{Key Idea}\label{key-idea-2}

The characteristic equation is a polynomial equation in the scalar
\(\lambda\). The Cayley-Hamilton theorem says the same polynomial, when
you plug in the matrix \(A\) instead of \(\lambda\), gives the zero
matrix. Constants become scalar multiples of the identity matrix \(I\).

\subsection{Full Proof}\label{full-proof}

Let \(A\) be an \(n \times n\) square matrix. Its characteristic
polynomial is:
\[
p(\lambda) = |A - \lambda I| = c_0 + c_1\lambda + c_2\lambda^2 + \dots + (-1)^n\lambda^n
\]
The adjoint of the characteristic matrix \((A - \lambda I)\) can be
written as a polynomial in \(\lambda\):
\[
\text{adj}(A - \lambda I) = B_{n-1}\lambda^{n-1} + B_{n-2}\lambda^{n-2} + \dots + B_0
\]
where \(B_0, B_1, \dots, B_{n-1}\) are constant \(n \times n\) matrices.
This works because each cofactor of \((A - \lambda I)\) is a polynomial
in \(\lambda\) of degree at most \(n - 1\).

By the matrix identity
\[
D \cdot \text{adj}(D) = \lvert D \rvert \cdot I
\]
applied to \(D = A - \lambda I\):
\[
(A - \lambda I) \cdot \text{adj}(A - \lambda I) = |A - \lambda I| \cdot I
\]
Substitute both polynomial expressions:
\[
(A - \lambda I)(B_{n-1}\lambda^{n-1} + B_{n-2}\lambda^{n-2} + \dots + B_0) = (c_0 + c_1\lambda + \dots + (-1)^n\lambda^n)I
\]
Expand the left side and equate coefficients of each power of
\(\lambda\):

For \(\lambda^0\):
\[
AB_0 = c_0 I
\]
For \(\lambda^1\):
\[
AB_1 - B_0 = c_1 I
\]
For \(\lambda^2\):
\[
AB_2 - B_1 = c_2 I
\]
\(\vdots\)

For \(\lambda^{n-1}\):
\[
AB_{n-1} - B_{n-2} = c_{n-1}I
\]
For \(\lambda^n\):
\[
-B_{n-1} = (-1)^n I
\]
Now pre-multiply these equations by \(I, A, A^2, \dots, A^n\)
respectively:
\[
AB_0 = c_0 I
\]
\[
A^2 B_1 - AB_0 = c_1 A
\]
\[
A^3 B_2 - A^2 B_1 = c_2 A^2
\]
\[
\vdots
\]
\[
A^n B_{n-1} - A^{n-1}B_{n-2} = c_{n-1}A^{n-1}
\]
\[
-A^n B_{n-1} = (-1)^n A^n
\]
Add all these equations. The left side telescopes (each \(A^k B_{k-1}\)
appears with \(+\) in one equation and \(-\) in the next). Everything
cancels:
\[
0 = c_0 I + c_1 A + c_2 A^2 + \dots + (-1)^n A^n = p(A)
\]
So \(p(A) = O\). The matrix satisfies its own characteristic equation.

\subsection{How to Verify the Theorem}\label{how-to-verify-the-theorem}

The verification problem is a standard exam question. Follow these
steps:

\begin{enumerate}
\def\labelenumi{\arabic{enumi}.}
\tightlist
\item
  Find the characteristic equation
\end{enumerate}
\[
\lvert A - \lambda I \rvert = 0.
\]
\begin{enumerate}
\def\labelenumi{\arabic{enumi}.}
\setcounter{enumi}{1}
\tightlist
\item
  Write the matrix version: replace \(\lambda^k\) with \(A^k\) and
  constants with \(cI\).
\item
  Compute \(A^2\) by multiplying \(A \cdot A\).
\item
  Compute \(A^3\) by multiplying \(A^2 \cdot A\) (for 3x3 matrices).
\item
  Substitute into the equation and verify every entry is zero.
\end{enumerate}

\subsection{Using Cayley-Hamilton to Find the
Inverse}\label{using-cayley-hamilton-to-find-the-inverse}

This is the most important application. If the characteristic equation
is:
\[
A^n + c_{n-1}A^{n-1} + \dots + c_1 A + c_0 I = O
\]
Multiply both sides by \(A^{-1}\):
\[
A^{n-1} + c_{n-1}A^{n-2} + \dots + c_1 I + c_0 A^{-1} = O
\]
Rearrange:
\[
A^{-1} = -\frac{1}{c_0}\left(A^{n-1} + c_{n-1}A^{n-2} + \dots + c_1 I\right)
\]
This requires \(c_0 \neq 0\) (which means \(\lvert A \rvert \neq 0\), so
\(A\) must be invertible).

For a 3x3 matrix with equation
\[
A^3 + aA^2 + bA + cI = O:
\]
\[
A^{-1} = -\frac{1}{c}(A^2 + aA + bI)
\]
For a 2x2 matrix with equation
\[
A^2 + aA + bI = O:
\]
\[
A^{-1} = -\frac{1}{b}(A + aI)
\]

\subsection{Worked Example 1: 2x2 Matrix
(Lecture)}\label{worked-example-1-2x2-matrix-lecture}

\textbf{Problem}: Verify Cayley-Hamilton and find \(A^{-1}\) for:
\[
A =
\begin{bmatrix}
1 & 2 \\
-1 & 3
\end{bmatrix}
\]
\textbf{Solution}:

\textbf{Step 1}: Characteristic equation.
\[
|A - \lambda I| = (1 - \lambda)(3 - \lambda) + 2 = \lambda^2 - 4\lambda + 5 = 0
\]
\textbf{Step 2}: Matrix equation:
\[
A^2 - 4A + 5I = O.
\]
\textbf{Step 3}: Compute \(A^2\):
\[
A^2 =
\begin{bmatrix}
1 & 2 \\
-1 & 3
\end{bmatrix}\begin{bmatrix} 1 & 2 \\ -1 & 3 \end{bmatrix} = \begin{bmatrix} -1 & 8 \\ -4 & 7 \end{bmatrix}
\]
\textbf{Step 4}: Verify:
\[
A^2 - 4A + 5I =
\begin{bmatrix}
-1 & 8 \\
-4 & 7
\end{bmatrix} - \begin{bmatrix} 4 & 8 \\ -4 & 12 \end{bmatrix} + \begin{bmatrix} 5 & 0 \\ 0 & 5 \end{bmatrix} = \begin{bmatrix} 0 & 0 \\ 0 & 0 \end{bmatrix} = O \quad \checkmark
\]
\textbf{Step 5}: Find \(A^{-1}\). Multiply by \(A^{-1}\):
\[
A - 4I + 5A^{-1} = O \implies A^{-1} = \frac{1}{5}(4I - A) = \frac{1}{5}
\begin{bmatrix}
3 & -2 \\
1 & 1
\end{bmatrix}
\]

\subsection{Worked Example 2: 3x3 Matrix (PYQ
2019)}\label{worked-example-2-3x3-matrix-pyq-2019}

\textbf{Problem}: Verify Cayley-Hamilton for the following matrix and
find \(A^{-1}\):
\[
A =
\begin{pmatrix}
2 & -1 & 1 \\
-1 & 2 & -1 \\
1 & -1 & 2
\end{pmatrix}
\]
\textbf{Solution}:

\textbf{Step 1}: Characteristic equation.
\[
|A - \lambda I| = 0 \implies \lambda^3 - 6\lambda^2 + 9\lambda - 4 = 0
\]
\textbf{Step 2}: Matrix equation:
\[
A^3 - 6A^2 + 9A - 4I = O.
\]
\textbf{Step 3}: Compute \(A^2\):
\[
A^2 =
\begin{pmatrix}
6 & -5 & 5 \\
-5 & 6 & -5 \\
5 & -5 & 6
\end{pmatrix}
\]
Compute
\[
A^3 = A^2 \cdot A:
\]
\[
A^3 =
\begin{pmatrix}
22 & -21 & 21 \\
-21 & 22 & -21 \\
21 & -21 & 22
\end{pmatrix}
\]
\textbf{Step 4}: Verify. Check diagonal entry \((1,1)\):
\(22 - 6(6) + 9(2) - 4 = 22 - 36 + 18 - 4 = 0\). Check off-diagonal
entry \((1,2)\): \(-21 - 6(-5) + 9(-1) - 0 = -21 + 30 - 9 = 0\). All
entries verify to zero.

\textbf{Step 5}: Find \(A^{-1}\). Multiply by \(A^{-1}\):
\[
A^2 - 6A + 9I = 4A^{-1}
\]
\[
4A^{-1} =
\begin{pmatrix}
6 & -5 & 5 \\
-5 & 6 & -5 \\
5 & -5 & 6
\end{pmatrix} - \begin{pmatrix} 12 & -6 & 6 \\ -6 & 12 & -6 \\ 6 & -6 & 12 \end{pmatrix} + \begin{pmatrix} 9 & 0 & 0 \\ 0 & 9 & 0 \\ 0 & 0 & 9 \end{pmatrix} = \begin{pmatrix} 3 & 1 & -1 \\ 1 & 3 & 1 \\ -1 & 1 & 3 \end{pmatrix}
\]
\[
A^{-1} = \frac{1}{4}
\begin{pmatrix}
3 & 1 & -1 \\
1 & 3 & 1 \\
-1 & 1 & 3
\end{pmatrix}
\]

\subsection{Worked Example 3: 3x3 Matrix (PYQ
2024)}\label{worked-example-3-3x3-matrix-pyq-2024}

\textbf{Problem}: Verify Cayley-Hamilton for \(A\) and compute
\(A^{-1}\):
\[
A =
\begin{bmatrix}
2 & 2 & 1 \\
1 & 3 & 1 \\
1 & 2 & 2
\end{bmatrix}
\]
\textbf{Solution}:

\textbf{Step 1}: Characteristic equation:
\[
\lambda^3 - 7\lambda^2 + 11\lambda - 5 = 0.
\]
\textbf{Step 2}: Matrix equation:
\[
A^3 - 7A^2 + 11A - 5I = O.
\]
\textbf{Step 3}: Compute:
\[
A^2 =
\begin{bmatrix}
7 & 12 & 6 \\
6 & 13 & 6 \\
6 & 12 & 7
\end{bmatrix}, \quad A^3 = \begin{bmatrix} 32 & 62 & 31 \\ 31 & 63 & 31 \\ 31 & 62 & 32 \end{bmatrix}
\]
\textbf{Step 4}: Verify. Entry \((1,1)\): \(32 - 49 + 22 - 5 = 0\).
Entry \((1,2)\): \(62 - 84 + 22 + 0 = 0\). All entries zero.

\textbf{Step 5}: Find \(A^{-1}\):
\[
5A^{-1} = A^2 - 7A + 11I =
\begin{bmatrix}
7 & 12 & 6 \\
6 & 13 & 6 \\
6 & 12 & 7
\end{bmatrix} - \begin{bmatrix} 14 & 14 & 7 \\ 7 & 21 & 7 \\ 7 & 14 & 14 \end{bmatrix} + \begin{bmatrix} 11 & 0 & 0 \\ 0 & 11 & 0 \\ 0 & 0 & 11 \end{bmatrix} = \begin{bmatrix} 4 & -2 & -1 \\ -1 & 3 & -1 \\ -1 & -2 & 4 \end{bmatrix}
\]
\[
A^{-1} = \frac{1}{5}
\begin{bmatrix}
4 & -2 & -1 \\
-1 & 3 & -1 \\
-1 & -2 & 4
\end{bmatrix}
\]

\subsection{Exam Patterns}\label{exam-patterns-26}

\begin{itemize}
\tightlist
\item
  The typical question is ``State and prove Cayley-Hamilton theorem'' (6
  marks) or ``Verify Cayley-Hamilton for matrix \(A\) and find
  \(A^{-1}\)'' (5-6 marks).
\item
  The proof was asked in 2018. Verification was asked in 2019 and 2024.
\item
  When verifying, show \(A^2\) computation fully. You can abbreviate
  \(A^3\) by showing a few entries if space is short.
\item
  The inverse step is almost always attached. Do not skip it.
\item
  The matrix from PYQ 2024 (\([2,2,1; 1,3,1; 1,2,2]\)) is the same
  matrix used in the eigenvalue problem. Examiners pair these topics.
\end{itemize}

\section{B12: Adjoint and Inverse of a
Matrix}\label{b12-adjoint-and-inverse-of-a-matrix}

\begin{quote}
\textbf{Exam Weight}: Tier 1-2 \textbar{} \textbf{Appeared in}: 2017,
2019, 2020, 2021, 2023 \textbar{} \textbf{Typical Marks}: 5--6
\end{quote}

Finding the inverse of a matrix is one of the most common exam
questions. Two methods exist: the adjoint method and the row operations
method. You must know both.

\subsection{Minors and Cofactors}\label{minors-and-cofactors}

\subsubsection{Minor}\label{minor}

The minor \(M_{ij}\) of element \(a_{ij}\) in a matrix \(A\) is the
determinant of the submatrix obtained by deleting row \(i\) and column
\(j\) from \(A\).

\subsubsection{Cofactor}\label{cofactor}

The cofactor \(C_{ij}\) of element \(a_{ij}\) is the signed minor:
\[
C_{ij} = (-1)^{i+j} M_{ij}
\]
The sign follows a checkerboard pattern:
\[
\begin{bmatrix}
+ & - & + \\
- & + & - \\
+ & - & +
\end{bmatrix}
\]
\subsubsection{Example}\label{example-5}

For matrix:
\[
A =
\begin{pmatrix}
0 & 1 & 2 \\
1 & 2 & 3 \\
3 & 1 & 1
\end{pmatrix}
\]
The minor \(M_{11}\) is obtained by deleting row 1 and column 1:
\[
M_{11} =
\begin{vmatrix}
2 & 3 \\
1 & 1
\end{vmatrix} = 2 - 3 = -1
\]
The cofactor
\[
C_{11} = (-1)^{1+1}(-1) = -1.
\]

\subsection{Adjoint Matrix}\label{adjoint-matrix}

\subsubsection{Definition}\label{definition-13}

The adjoint of a square matrix \(A\) (written \(\text{adj}(A)\)) is the
transpose of the cofactor matrix.

\textbf{Step 1}: Find all cofactors \(C_{ij}\).

\textbf{Step 2}: Form the cofactor matrix \(C\).

\textbf{Step 3}: Transpose it:
\[
\text{adj}(A) = C^T.
\]
\subsubsection{Key Property}\label{key-property}
\[
A \cdot \text{adj}(A) = |A| \cdot I
\]
This is the foundation of the adjoint method for finding the inverse.

\subsection{Inverse by Adjoint Method}\label{inverse-by-adjoint-method}

\subsubsection{Definition}\label{definition-14}

A square matrix \(A\) is invertible if there exists a matrix \(B\) such
that \(AB = BA = I\). This \(B\) is called \(A^{-1}\).

A matrix is invertible if and only if \(\lvert A \rvert \neq 0\).

\subsubsection{Formula}\label{formula}
\[
A^{-1} = \frac{1}{|A|} \text{adj}(A)
\]
\subsubsection{Worked Example 1: Adjoint Method (PYQ
2017)}\label{worked-example-1-adjoint-method-pyq-2017}

\textbf{Problem}: Find the adjoint and inverse of:
\[
A =
\begin{pmatrix}
0 & 1 & 2 \\
1 & 2 & 3 \\
3 & 1 & 1
\end{pmatrix}
\]
\textbf{Solution}:

\textbf{Step 1}: Find the determinant:
\[
|A| = 0(2-3) - 1(1-9) + 2(1-6) = 0 + 8 - 10 = -2
\]
Since \(\lvert A \rvert \neq 0\), the inverse exists.

\textbf{Step 2}: Find all cofactors:
\[
C_{11} = +(2-3) = -1, \quad C_{12} = -(1-9) = 8, \quad C_{13} = +(1-6) = -5
\]
\[
C_{21} = -(1-2) = 1, \quad C_{22} = +(0-6) = -6, \quad C_{23} = -(0-3) = 3
\]
\[
C_{31} = +(3-4) = -1, \quad C_{32} = -(0-2) = 2, \quad C_{33} = +(0-1) = -1
\]
\textbf{Step 3}: Adjoint = transpose of cofactor matrix:
\[
\text{adj}(A) =
\begin{pmatrix}
-1 & 1 & -1 \\
8 & -6 & 2 \\
-5 & 3 & -1
\end{pmatrix}
\]
\textbf{Step 4}: Inverse:
\[
A^{-1} = \frac{1}{-2}
\begin{pmatrix}
-1 & 1 & -1 \\
8 & -6 & 2 \\
-5 & 3 & -1
\end{pmatrix} = \begin{pmatrix} 1/2 & -1/2 & 1/2 \\ -4 & 3 & -1 \\ 5/2 & -3/2 & 1/2 \end{pmatrix}
\]

\subsubsection{Worked Example 2: Rotation Matrix Inverse (PYQ
2020)}\label{worked-example-2-rotation-matrix-inverse-pyq-2020}

\textbf{Problem}: Find the adjoint and inverse of:
\[
A =
\begin{bmatrix}
\cos\theta & -\sin\theta & 0 \\
\sin\theta & \cos\theta & 0 \\
0 & 0 & 1
\end{bmatrix}
\]
\textbf{Solution}:

Determinant:
\[
\lvert A \rvert = \cos^2\theta + \sin^2\theta = 1.
\]
Cofactors:
\[
C_{11} = \cos\theta, \quad C_{12} = -\sin\theta, \quad C_{13} = 0
\]
\[
C_{21} = \sin\theta, \quad C_{22} = \cos\theta, \quad C_{23} = 0
\]
\[
C_{31} = 0, \quad C_{32} = 0, \quad C_{33} = 1
\]
Adjoint (transpose of cofactor matrix):
\[
\text{adj}(A) =
\begin{bmatrix}
\cos\theta & \sin\theta & 0 \\
-\sin\theta & \cos\theta & 0 \\
0 & 0 & 1
\end{bmatrix}
\]
Since
\[
\lvert A \rvert = 1:
\]
\[
A^{-1} =
\begin{bmatrix}
\cos\theta & \sin\theta & 0 \\
-\sin\theta & \cos\theta & 0 \\
0 & 0 & 1
\end{bmatrix}
\]
Notice:
\[
A^{-1} = A^T
\]
. This is because \(A\) is an orthogonal matrix (rotation matrix).

\subsection{Inverse by Row Operations}\label{inverse-by-row-operations}

\subsubsection{The Method}\label{the-method-6}

\begin{enumerate}
\def\labelenumi{\arabic{enumi}.}
\tightlist
\item
  Write the augmented matrix \([A \; \lvert  \; I]\).
\item
  Apply elementary row operations to transform \(A\) into \(I\).
\item
  The same operations transform \(I\) into \(A^{-1}\).
\item
  When the left side becomes \(I\), the right side is \(A^{-1}\).
\end{enumerate}

\subsubsection{Worked Example 3: Row Operations (PYQ
2019)}\label{worked-example-3-row-operations-pyq-2019}

\textbf{Problem}: Find the inverse of:
\[
A =
\begin{pmatrix}
1 & 3 & 3 \\
1 & 4 & 3 \\
1 & 3 & 4
\end{pmatrix}
\]
\textbf{Solution}:

Set up \([A \lvert  I]\):
\[
\begin{bmatrix}
1 & 3 & 3 & | & 1 & 0 & 0 \\
1 & 4 & 3 & | & 0 & 1 & 0 \\
1 & 3 & 4 & | & 0 & 0 & 1
\end{bmatrix}
\]
Apply \(R_2 \to R_2 - R_1\) and \(R_3 \to R_3 - R_1\):
\[
\begin{bmatrix}
1 & 3 & 3 & | & 1 & 0 & 0 \\
0 & 1 & 0 & | & -1 & 1 & 0 \\
0 & 0 & 1 & | & -1 & 0 & 1
\end{bmatrix}
\]
Apply \(R_1 \to R_1 - 3R_2\):
\[
\begin{bmatrix}
1 & 0 & 3 & | & 4 & -3 & 0 \\
0 & 1 & 0 & | & -1 & 1 & 0 \\
0 & 0 & 1 & | & -1 & 0 & 1
\end{bmatrix}
\]
Apply \(R_1 \to R_1 - 3R_3\):
\[
\begin{bmatrix}
1 & 0 & 0 & | & 7 & -3 & -3 \\
0 & 1 & 0 & | & -1 & 1 & 0 \\
0 & 0 & 1 & | & -1 & 0 & 1
\end{bmatrix}
\]
The inverse is:
\[
A^{-1} =
\begin{pmatrix}
7 & -3 & -3 \\
-1 & 1 & 0 \\
-1 & 0 & 1
\end{pmatrix}
\]

\subsubsection{Worked Example 4: Row Operations (PYQ
2021)}\label{worked-example-4-row-operations-pyq-2021}

\textbf{Problem}: Find the inverse by row operations:
\[
A =
\begin{pmatrix}
2 & 1 & 2 \\
2 & 2 & 1 \\
1 & 2 & 2
\end{pmatrix}
\]
\textbf{Solution}:

Set up \([A \lvert  I]\) and swap \(R_1 \leftrightarrow R_3\):
\[
\begin{bmatrix}
1 & 2 & 2 & | & 0 & 0 & 1 \\
2 & 2 & 1 & | & 0 & 1 & 0 \\
2 & 1 & 2 & | & 1 & 0 & 0
\end{bmatrix}
\]
Apply \(R_2 \to R_2 - 2R_1\) and \(R_3 \to R_3 - 2R_1\):
\[
\begin{bmatrix}
1 & 2 & 2 & | & 0 & 0 & 1 \\
0 & -2 & -3 & | & 0 & 1 & -2 \\
0 & -3 & -2 & | & 1 & 0 & -2
\end{bmatrix}
\]
Scale
\[
R_2 \to -\frac{1}{2}R_2:
\]
\[
\begin{bmatrix}
1 & 2 & 2 & | & 0 & 0 & 1 \\
0 & 1 & 3/2 & | & 0 & -1/2 & 1 \\
0 & -3 & -2 & | & 1 & 0 & -2
\end{bmatrix}
\]
Apply \(R_3 \to R_3 + 3R_2\):
\[
\begin{bmatrix}
1 & 2 & 2 & | & 0 & 0 & 1 \\
0 & 1 & 3/2 & | & 0 & -1/2 & 1 \\
0 & 0 & 5/2 & | & 1 & -3/2 & 1
\end{bmatrix}
\]
Scale
\[
R_3 \to \frac{2}{5}R_3
\]
, then back-substitute upward. After all operations:
\[
A^{-1} = \frac{1}{5}
\begin{pmatrix}
2 & 2 & -3 \\
-3 & 2 & 2 \\
2 & -3 & 2
\end{pmatrix}
\]

\subsubsection{Worked Example 5: Row Operations (PYQ
2023)}\label{worked-example-5-row-operations-pyq-2023}

\textbf{Problem}: Find the inverse by row transformation:
\[
A =
\begin{bmatrix}
1 & 3 & 5 \\
2 & 4 & 5 \\
3 & 7 & 6
\end{bmatrix}
\]
\textbf{Solution}:

Set up \([A \lvert  I]\). Apply \(R_2 \to R_2 - 2R_1\) and
\(R_3 \to R_3 - 3R_1\):
\[
\begin{bmatrix}
1 & 3 & 5 & | & 1 & 0 & 0 \\
0 & -2 & -5 & | & -2 & 1 & 0 \\
0 & -2 & -9 & | & -3 & 0 & 1
\end{bmatrix}
\]
Apply \(R_3 \to R_3 - R_2\):
\[
\begin{bmatrix}
1 & 3 & 5 & | & 1 & 0 & 0 \\
0 & -2 & -5 & | & -2 & 1 & 0 \\
0 & 0 & -4 & | & -1 & -1 & 1
\end{bmatrix}
\]
Scale
\[
R_2 \to -\frac{1}{2}R_2
\]
and
\[
R_3 \to -\frac{1}{4}R_3
\]
. Then back-substitute:
\[
A^{-1} = \frac{1}{8}
\begin{bmatrix}
-11 & 17 & -5 \\
3 & -9 & 5 \\
2 & 2 & -2
\end{bmatrix}
\]

\subsection{Which Method to Use?}\label{which-method-to-use}

\begin{longtable}[]{@{}
  >{\raggedright\arraybackslash}p{(\linewidth - 6\tabcolsep) * \real{0.2667}}
  >{\raggedright\arraybackslash}p{(\linewidth - 6\tabcolsep) * \real{0.3333}}
  >{\raggedright\arraybackslash}p{(\linewidth - 6\tabcolsep) * \real{0.2000}}
  >{\raggedright\arraybackslash}p{(\linewidth - 6\tabcolsep) * \real{0.2000}}@{}}
\toprule\noalign{}
\begin{minipage}[b]{\linewidth}\raggedright
Method
\end{minipage} & \begin{minipage}[b]{\linewidth}\raggedright
Best For
\end{minipage} & \begin{minipage}[b]{\linewidth}\raggedright
Pros
\end{minipage} & \begin{minipage}[b]{\linewidth}\raggedright
Cons
\end{minipage} \\
\midrule\noalign{}
\endhead
\bottomrule\noalign{}
\endlastfoot
Adjoint & Small matrices, when adjoint is also asked & Gives both
adjoint and inverse & Tedious for large matrices (9 cofactors for
3x3) \\
Row operations & Any size matrix, when only inverse is asked &
Systematic, fewer calculations & Does not give the adjoint \\
Cayley-Hamilton & When eigenvalues or characteristic equation is already
known & Elegant, uses known information & Requires \(A^2\)
computation \\
\end{longtable}

If the problem says ``find adjoint and inverse,'' use the adjoint
method. If the problem says ``find inverse by row operations,'' use row
operations. If the problem says ``verify Cayley-Hamilton and find
inverse,'' use Cayley-Hamilton (see B11).

\subsection{Important Proofs}\label{important-proofs}

\subsubsection{Proof: Inverse of a
Product}\label{proof-inverse-of-a-product}

\textbf{Statement}:
\[
(AB)^{-1} = B^{-1}A^{-1}
\]
(the order reverses).

\textbf{Proof}: We need to show that \(B^{-1}A^{-1}\) satisfies the
definition of the inverse of \(AB\).
\[
(AB)(B^{-1}A^{-1}) = A(BB^{-1})A^{-1} = AIA^{-1} = AA^{-1} = I
\]
\[
(B^{-1}A^{-1})(AB) = B^{-1}(A^{-1}A)B = B^{-1}IB = B^{-1}B = I
\]
Since
\[
(AB)(B^{-1}A^{-1}) = I
\]
and
\[
(B^{-1}A^{-1})(AB) = I
\]
, we have
\[
(AB)^{-1} = B^{-1}A^{-1}.
\]
This proof has appeared in 2018 and 2023.

\subsection{Exam Patterns}\label{exam-patterns-27}

\begin{itemize}
\tightlist
\item
  \textbf{Inverse appears in 5 out of 7 papers.} Usually ``Find the
  inverse by row operations'' or ``Find the adjoint and inverse.''
\item
  Row operations method is more common in recent papers (2019, 2021,
  2023).
\item
  The adjoint method was popular in older papers (2017, 2020).
\item
  The proof
\end{itemize}
\[
(AB)^{-1} = B^{-1}A^{-1}
\]
has appeared twice. Memorize it. - Always check:
\(\lvert A \rvert \neq 0\) before attempting inverse. State this
explicitly. - Verify your answer by multiplying \(AA^{-1}\) (at least
check one row) if time permits.

\section{B13: Diagonalization, Matrix Exponential, and Matrix
ODE}\label{b13-diagonalization-matrix-exponential-and-matrix-ode}

\begin{quote}
\textbf{Exam Weight}: Tier 2-3 \textbar{} \textbf{Appeared in}: 2017,
2020, 2021 \textbar{} \textbf{Typical Marks}: 6--12
\end{quote}

This topic combines eigenvalues with differential equations. The matrix
ODE problem appeared identically in 2020 and 2021.

\subsection{Diagonalization}\label{diagonalization}

\subsubsection{What is Diagonalization?}\label{what-is-diagonalization}

A square matrix \(A\) is diagonalizable if there exists an invertible
matrix \(P\) such that:
\[
P^{-1}AP = D
\]
where \(D\) is a diagonal matrix. The diagonal entries of \(D\) are the
eigenvalues of \(A\). The columns of \(P\) are the corresponding
eigenvectors.

\subsubsection{When is Diagonalization
Possible?}\label{when-is-diagonalization-possible}

An \(n \times n\) matrix is diagonalizable if and only if it has \(n\)
linearly independent eigenvectors.

\begin{itemize}
\tightlist
\item
  If all \(n\) eigenvalues are distinct, the matrix is always
  diagonalizable.
\item
  If eigenvalues are repeated, diagonalization depends on whether enough
  independent eigenvectors exist. If the geometric multiplicity (number
  of independent eigenvectors) equals the algebraic multiplicity (repeat
  count) for each eigenvalue, the matrix is diagonalizable.
\end{itemize}

\subsubsection{The Modal Matrix}\label{the-modal-matrix}

The matrix \(P\) whose columns are the eigenvectors is called the
\textbf{modal matrix}:
\[
P = [X_1 \; X_2 \; \dots \; X_n]
\]
Then
\[
A = PDP^{-1}.
\]

\subsection{When Diagonalization Fails (PYQ
2017)}\label{when-diagonalization-fails-pyq-2017}

The matrix
\[
A =
\begin{bmatrix}
1 & 1 \\
0 & 1
\end{bmatrix}
\]
has eigenvalue \(\lambda = 1\) with algebraic multiplicity 2.

Solving \((A - I)X = 0\):
\[
\begin{bmatrix}
0 & 1 \\
0 & 0
\end{bmatrix}X = 0
\]
This gives only one eigenvector:
\[
X =
\begin{bmatrix}
1 \\
0
\end{bmatrix}.
\]
Geometric multiplicity = 1 \textless{} algebraic multiplicity = 2. So
\(A\) is NOT diagonalizable.

\subsection{Matrix Exponential}\label{matrix-exponential}

\subsubsection{Definition}\label{definition-15}

The matrix exponential \(e^A\) is defined by the power series:
\[
e^A = I + A + \frac{A^2}{2!} + \frac{A^3}{3!} + \dots
\]
\subsubsection{Computation Using
Diagonalization}\label{computation-using-diagonalization}

If
\[
A = PDP^{-1}:
\]
\[
e^A = Pe^D P^{-1}
\]
where \(e^D\) is simply:
\[
e^D =
\begin{bmatrix}
e^{\lambda_1} & 0 & \dots \\
0 & e^{\lambda_2} & \dots \\
\vdots & \vdots & \ddots
\end{bmatrix}
\]
\subsubsection{Computation Using Decomposition (PYQ
2017)}\label{computation-using-decomposition-pyq-2017}

When diagonalization fails, decompose \(A = I + B\) where \(B\) is
nilpotent.

For
\[
A =
\begin{bmatrix}
1 & 1 \\
0 & 1
\end{bmatrix}:
\]
Set
\[
B =
\begin{bmatrix}
0 & 1 \\
0 & 0
\end{bmatrix}.
\]
Then
\[
B^2 = O.
\]
\[
e^A = e^{I+B} = e^I \cdot e^B = e(I + B) = e
\begin{bmatrix}
1 & 1 \\
0 & 1
\end{bmatrix} = \begin{bmatrix} e & e \\ 0 & e \end{bmatrix}
\]
For time-dependent version:
\[
e^{At} = e^t(I + Bt) = e^t
\begin{bmatrix}
1 & t \\
0 & 1
\end{bmatrix} = \begin{bmatrix} e^t & te^t \\ 0 & e^t \end{bmatrix}
\]

\subsection{Solving Matrix ODE}\label{solving-matrix-ode}

\subsubsection{The System}\label{the-system}

Given a system of first-order linear ODEs:
\[
\frac{dX}{dt} = MX
\]
where \(M\) is a constant matrix and \(X\) is the vector of unknowns.

\subsubsection{The Method}\label{the-method-7}

\begin{enumerate}
\def\labelenumi{\arabic{enumi}.}
\tightlist
\item
  Find eigenvalues of \(M\).
\item
  Find eigenvectors of \(M\).
\item
  The general solution is:
\end{enumerate}
\[
X(t) = C_1 e^{\lambda_1 t} v_1 + C_2 e^{\lambda_2 t} v_2 + \dots
\]
where \(\lambda_i\) are eigenvalues and \(v_i\) are corresponding
eigenvectors.

\subsection{Must-Know Problem: Matrix ODE (PYQ 2020, 2021 ---
Repeated)}\label{must-know-problem-matrix-ode-pyq-2020-2021-repeated}

\textbf{Problem}: Solve by matrix method:
\[
\frac{dx}{dt} = 6x - 3y, \quad \frac{dy}{dt} = 2x + y
\]
\textbf{Solution}:

Write in matrix form
\[
\frac{dX}{dt} = MX
\]
where:
\[
M =
\begin{bmatrix}
6 & -3 \\
2 & 1
\end{bmatrix}
\]
\textbf{Step 1}: Eigenvalues.
\[
\lvert M - \lambda I \rvert = 0:
\]
\[
(6-\lambda)(1-\lambda) + 6 = \lambda^2 - 7\lambda + 12 = (\lambda - 3)(\lambda - 4) = 0
\]
So
\[
\lambda_1 = 3
\]
and
\[
\lambda_2 = 4.
\]
\textbf{Step 2}: Eigenvectors.

For \(\lambda = 3\): \(3x - 3y = 0 \implies x = y\). Eigenvector:
\[
v_1 =
\begin{bmatrix}
1 \\
1
\end{bmatrix}.
\]
For \(\lambda = 4\):
\[
2x - 3y = 0 \implies x = \frac{3}{2}y
\]
. Eigenvector:
\[
v_2 =
\begin{bmatrix}
3 \\
2
\end{bmatrix}.
\]
\textbf{Step 3}: General solution:
\[
\begin{bmatrix}
x(t) \\
y(t)
\end{bmatrix} = C_1 e^{3t}\begin{bmatrix} 1 \\ 1 \end{bmatrix} + C_2 e^{4t}\begin{bmatrix} 3 \\ 2 \end{bmatrix}
\]
Component form:
\[
x(t) = C_1 e^{3t} + 3C_2 e^{4t}
\]
\[
y(t) = C_1 e^{3t} + 2C_2 e^{4t}
\]

\subsection{Exam Patterns}\label{exam-patterns-28}

\begin{itemize}
\tightlist
\item
  Matrix ODE (\(dx/dt = 6x - 3y\), \(dy/dt = 2x + y\)) appeared
  identically in 2020 and 2021. Memorize it.
\item
  The 2017 paper had a 12-mark diagonalization + \(e^A\) + ODE problem.
\item
  The method is always: eigenvalues then eigenvectors then write the
  general solution.
\item
  If the matrix is not diagonalizable, mention it and use the
  decomposition method for \(e^A\).
\end{itemize}

\section{C01: Vector Space Definition and
Properties}\label{c01-vector-space-definition-and-properties}

\begin{quote}
\textbf{Exam Weight}: Tier 1-2 \textbar{} \textbf{Appeared in}: 2017,
2018, 2024 \textbar{} \textbf{Typical Marks}: 2--8
\end{quote}

A vector space is the fundamental structure of linear algebra. It is a
set of elements (vectors) that can be scaled and added together
according to specific mathematical rules.

\subsection{Formal Definition}\label{formal-definition}

Let \(V\) be a non-empty set of objects called vectors. Let \(K\) be a
field of scalars (usually real numbers \(\mathbb{R}\) or complex numbers
\(\mathbb{C}\)).

The set \(V\) is a \textbf{vector space} over the field \(K\) if it
satisfies two closure conditions and eight axioms under vector addition
and scalar multiplication:

\textbf{Closure Properties}: 1. \textbf{Closure under Addition:} If
\(u, v \in V\), then \(u + v \in V\). 2. \textbf{Closure under Scalar
Multiplication:} If \(u \in V\) and \(k \in K\), then \(k u \in V\).

\textbf{Addition Axioms}: 3. \textbf{Commutativity:} \(u + v = v + u\)
4. \textbf{Associativity:} \((u + v) + w = u + (v + w)\) 5.
\textbf{Additive Identity:} There exists a zero vector \(0 \in V\) such
that \(u + 0 = u\). 6. \textbf{Additive Inverse:} For every \(u \in V\),
there exists \(-u \in V\) such that \(u + (-u) = 0\).

\textbf{Scalar Multiplication Axioms}: 7. \textbf{Distributivity (Scalar
over Vector Addition):} \(k(u + v) = ku + kv\) 8. \textbf{Distributivity
(Vector over Scalar Addition):} \((a + b)u = au + bu\) 9.
\textbf{Associativity of Multiplication:} \(a(bu) = (ab)u\) 10.
\textbf{Unitary Identity:} \(1u = u\), where \(1\) is the multiplicative
identity in \(K\).

\subsection{Must-Know Proofs: Basic Properties (PYQ
2017)}\label{must-know-proofs-basic-properties-pyq-2017}

You must be able to prove four basic properties of vector spaces using
\emph{only} the axioms above. This is an 8-mark question.

\subsubsection{Proof 1}\label{proof-1-1}

We know \(0 + 0 = 0\). Multiply by scalar \(k\): \(k(0 + 0) = k0\). By
distributivity: \(k0 + k0 = k0\). Add the additive inverse \(-(k0)\) to
both sides: \((k0 + k0) + (-(k0)) = k0 + (-(k0))\)
\(k0 + (k0 - k0) = 0\) \(k0 + 0 = 0 \implies k0 = 0\). \(\blacksquare\)

\subsubsection{Proof 2}\label{proof-2-1}

In the scalar field, \(0 + 0 = 0\). Multiply by vector \(u\):
\((0 + 0)u = 0u\). By distributivity: \(0u + 0u = 0u\). Add the additive
inverse \(-(0u)\) to both sides: \((0u + 0u) + (-(0u)) = 0u + (-(0u))\)
\(0u + (0u - 0u) = 0\) \(0u + 0 = 0 \implies 0u = 0\). \(\blacksquare\)

\subsubsection{Proof 3: If, then or}\label{proof-3-if-then-or}

Assume \(ku = 0\). If \(k = 0\), the statement holds. If \(k \neq 0\),
then its multiplicative inverse \(k^{-1}\) exists in the field \(K\).
Multiply both sides by \(k^{-1}\):
\[
k^{-1}(ku) = k^{-1}(0)
\]
By associativity and Proof 1:
\[
(k^{-1}k)u = 0
\]
\(1u = 0\) By the unitary axiom: \(u = 0\). Thus, either \(k=0\) or
\(u=0\). \(\blacksquare\)

\subsubsection{Proof 4}\label{proof-4-1}

The vector \(-ku\) is the unique additive inverse of \(ku\), meaning
\(ku + (-ku) = 0\). Show that \((-k)u\) is the additive inverse:
\(ku + (-k)u = (k + (-k))u = 0u = 0\) (by Proof 2). Thus,
\((-k)u = -ku\).

Show that \(k(-u)\) is the additive inverse:
\(ku + k(-u) = k(u + (-u)) = k(0) = 0\) (by Proof 1). Thus,
\(k(-u) = -ku\). Therefore, \((-k)u = k(-u) = -ku\). \(\blacksquare\)

\subsection{Exam Patterns}\label{exam-patterns-29}

\begin{itemize}
\tightlist
\item
  When asked to ``define a vector space'', you must list the properties
  clearly. Bullet points are fine, but be sure to mention the field
  \(K\), the closure properties, and the 8 axioms.
\item
  The 4-part proof from 2017 is a classic test of abstract algebra
  logic. Always explicitly state which axiom you are using at each step
  (e.g., ``by distributivity'').
\end{itemize}

\section{C02: Vector Subspaces}\label{c02-vector-subspaces}

\begin{quote}
\textbf{Exam Weight}: Tier 2 \textbar{} \textbf{Appeared in}: 2017,
2018, 2024 \textbar{} \textbf{Typical Marks}: 3--4
\end{quote}

A subspace is a smaller vector space entirely contained within a larger
one.

\subsection{Definition of a Subspace}\label{definition-of-a-subspace}

Let \(W\) be a non-empty subset of a vector space \(V\) over a field
\(K\). \(W\) is called a \textbf{subspace} of \(V\) if \(W\) is itself a
vector space under the same vector addition and scalar multiplication
defined on \(V\).

\subsubsection{The Subspace Test}\label{the-subspace-test}

To prove that a subset \(W\) is a subspace, you do \emph{not} need to
check all 10 vector space axioms. You only need to verify three things:

\begin{enumerate}
\def\labelenumi{\arabic{enumi}.}
\tightlist
\item
  \textbf{Non-empty / Zero Vector}: The zero vector of \(V\) must be in
  \(W\) (\(0 \in W\)).
\item
  \textbf{Closed under Addition}: For any vectors \(u, v \in W\), their
  sum must also be in \(W\) (\(u + v \in W\)).
\item
  \textbf{Closed under Scalar Multiplication}: For any vector
  \(u \in W\) and any scalar \(k \in K\), the scaled vector must also be
  in \(W\) (\(ku \in W\)).
\end{enumerate}

\emph{Note: Conditions 2 and 3 are sometimes combined into a single
condition: For any scalars \(a,b\) and vectors \(u,v\),
\(au + bv \in W\).}

\subsection{Worked Example: Subspace Verification (PYQ
2024)}\label{worked-example-subspace-verification-pyq-2024}

\textbf{Problem}: Let \(U\) consist of all vectors in \(\mathbb{R}^3\)
whose entries are equal, that is
\(U = \lbrace(a, b, c) \mid a = b = c\rbrace\). Prove that \(U\) is a
vector space over \(\mathbb{R}\).

\textbf{Solution}: Since \(U\) is a subset of the known vector space
\(\mathbb{R}^3\), we only need to show it is a subspace. We apply the
subspace test.

\textbf{Step 1: Contains the Zero Vector} The zero vector of
\(\mathbb{R}^3\) is \((0, 0, 0)\). Since \(0 = 0 = 0\) satisfies the
condition \(a=b=c\), the zero vector is in \(U\).

\textbf{Step 2: Closed Under Addition} Let
\[
u_1 = (a_1, a_1, a_1) \in U
\]
and
\[
u_2 = (a_2, a_2, a_2) \in U.
\]
We add these vectors:
\[
u_1 + u_2 = (a_1 + a_2, a_1 + a_2, a_1 + a_2)
\]
Since all three components of the sum vector are identical (equal to
\(a_1 + a_2\)), the sum vector is in \(U\).

\textbf{Step 3: Closed Under Scalar Multiplication} Let
\(u = (a, a, a) \in U\) and \(k \in \mathbb{R}\) be a scalar. We
multiply the vector by the scalar:
\[
k u = (ka, ka, ka)
\]
Since all three components are identical (equal to \(ka\)), this vector
is in \(U\).

Since \(U\) satisfies all three conditions, it is a subspace of
\(\mathbb{R}^3\). Because a subspace is itself a vector space, \(U\) is
a vector space over \(\mathbb{R}\). \(\blacksquare\)

\subsection{Common Examples of
Subspaces}\label{common-examples-of-subspaces}

\begin{itemize}
\tightlist
\item
  \textbf{Trivial Subspaces}: For any vector space \(V\), the set
  containing only the zero vector \(\lbrace 0\rbrace\) and the entire
  space \(V\) itself are both subspaces.
\item
  \textbf{Lines and Planes}: In \(\mathbb{R}^3\), any line that passes
  through the origin is a subspace. Any flat plane that passes through
  the origin (e.g., \(z=0\)) is a subspace. \emph{Warning: If a line or
  plane does NOT pass through the origin (e.g., \(z=5\)), it is NOT a
  subspace because it fails the zero vector test.}
\end{itemize}

\subsection{Exam Patterns}\label{exam-patterns-30}

\begin{itemize}
\tightlist
\item
  When asked to prove something is a vector space, usually it is a
  subset of a known space (like \(\mathbb{R}^n\) or the space of
  polynomials). Do not list all 10 axioms; just use the 3-step subspace
  test.
\item
  Always check the zero vector first. It is the easiest way to prove a
  set is \emph{not} a subspace. If plugging in zeros for the variables
  breaks the given condition, you're done.
\end{itemize}

\section{C03: Sum and Direct Sum of
Subspaces}\label{c03-sum-and-direct-sum-of-subspaces}

\begin{quote}
\textbf{Exam Weight}: Tier 1-2 \textbar{} \textbf{Appeared in}: 2019,
2023, 2024 \textbar{} \textbf{Typical Marks}: 2--6
\end{quote}

When you have multiple subspaces within a vector space, you can combine
them. The most important way to combine them is via a ``Direct Sum'',
which guarantees a unique decomposition of vectors.

\subsection{Definitions and Examples}\label{definitions-and-examples}

\textbf{Sum of Subspaces (\(U + W\))}: Let \(U\) and \(W\) be subspaces
of a vector space \(V\). The sum \(U + W\) is the set of all possible
vectors \(u + w\) where \(u \in U\) and \(w \in W\). - \emph{Example}:
In \(\mathbb{R}^2\), let \(U\) be the \(x\)-axis and \(W\) be the line
\(y=x\). Their sum \(U+W\) covers the entire \(\mathbb{R}^2\) plane.

\textbf{Direct Sum (\(U \oplus W\))}: The sum \(U + W\) is called a
\textbf{direct sum} (written \(U \oplus W\)) if every element in the sum
can be written \emph{uniquely} as \(u + w\). - This occurs if and only
if the only vector shared by both subspaces is the zero vector:
\(U \cap W = \lbrace 0\rbrace\). - \emph{Example}: In \(\mathbb{R}^2\),
let \(U\) be the \(x\)-axis (vectors \((x,0)\)) and \(W\) be the
\(y\)-axis (vectors \((0,y)\)). Their intersection is only the origin
\((0,0)\). Thus,
\[
\mathbb{R}^2 = U \oplus W
\]
, because any vector \((x,y)\) has exactly one decomposition:
\((x,0) + (0,y)\).

\subsection{Must-Know Proof (PYQ 2019)}\label{must-know-proof-pyq-2019}

\textbf{Problem}: Prove that the vector space \(V\) is the direct sum of
its subspaces \(U\) and \(W\) if and only if (i) \(V = U + W\), and (ii)
\(U \cap W = \lbrace 0\rbrace\).

\textbf{Proof}:

\textbf{Part 1: Forward Direction (\(\implies\))} Assume \(V\) is the
direct sum of \(U\) and \(W\) (\(V = U \oplus W\)). By definition, every
vector \(v \in V\) can be \emph{uniquely} written as \(v = u + w\)
(where \(u \in U, w \in W\)). - Since every \(v\) can be written as
\(u+w\), we have \(V = U + W\). - To prove
\(U \cap W = \lbrace 0\rbrace\), let \(v \in U \cap W\). Since
\(v \in U\), we can write it as \(v = v + 0\) (where \(v \in U\) and
\(0 \in W\)). Since \(v \in W\), we can write it as \(v = 0 + v\) (where
\(0 \in U\) and \(v \in W\)). Because the decomposition must be unique,
these two expressions must be identical. Therefore, \(v = 0\). Thus,
\(U \cap W = \lbrace 0\rbrace\).

\textbf{Part 2: Reverse Direction (\(\impliedby\))} Assume (i)
\(V = U + W\) and (ii) \(U \cap W = \lbrace 0\rbrace\). Since
\(V = U + W\), every \(v \in V\) can be written as \(v = u + w\). We
must show this representation is unique. Suppose \(v\) can be
represented in two ways:
\[
v = u_1 + w_1 \quad \text{and} \quad v = u_2 + w_2
\]
Equating them:
\[
u_1 + w_1 = u_2 + w_2 \implies u_1 - u_2 = w_2 - w_1
\]
Let this new vector be \(x\). Because \(U\) is a subspace,
\[
x = u_1 - u_2 \in U.
\]
Because \(W\) is a subspace,
\[
x = w_2 - w_1 \in W.
\]
Therefore, \(x \in U \cap W\). But we are given
\(U \cap W = \lbrace 0\rbrace\). Thus, \(x = 0\).
\[
u_1 - u_2 = 0 \implies u_1 = u_2
\]
\[
w_2 - w_1 = 0 \implies w_1 = w_2
\]
Since the components are identical, the representation is unique.
Therefore, \(V = U \oplus W\). \(\blacksquare\)

\subsection{Exam Patterns}\label{exam-patterns-31}

\begin{itemize}
\tightlist
\item
  The proof above is a highly standard 6-mark linear algebra theorem.
  Memorize the trick for the forward direction (\(v=v+0\) vs \(v=0+v\))
  and the reverse direction (
\end{itemize}
\[
u_1-u_2 = w_2-w_1
\]
). - When asked to ``define'' direct sum for 2 marks, always include the
mathematical condition \(U \cap W = \lbrace 0\rbrace\) and provide the
\(x\)-axis / \(y\)-axis example in \(\mathbb{R}^2\) for full credit.

\section{C04: Linear Combination and
Span}\label{c04-linear-combination-and-span}

\begin{quote}
\textbf{Exam Weight}: Tier 2 \textbar{} \textbf{Appeared in}: 2020,
2021, 2023, 2024 \textbar{} \textbf{Typical Marks}: 3--5
\end{quote}

A vector is a \textbf{linear combination} of a set of vectors if it can
be created by scaling those vectors and adding them together.

\subsection{The Core Concept}\label{the-core-concept}

To express a vector \(v\) as a linear combination of vectors
\(e_1, e_2, \dots, e_n\), we set up the equation:
\[
v = c_1 e_1 + c_2 e_2 + \dots + c_n e_n
\]
This generates a system of linear equations where the unknowns are the
scalars \(c_1, c_2, \dots, c_n\). Solving this system (usually via an
augmented matrix and row reduction) gives the required coefficients.

\subsection{Worked Example 1: Coordinate Vectors (PYQ 2020,
2021)}\label{worked-example-1-coordinate-vectors-pyq-2020-2021}

\textbf{Problem}: Write the vector \(v = (1, -2, 5)\) as a linear
combination of the vectors
\[
e_1 = (1, 1, 1),
\]
\[
e_2 = (1, 2, 3)
\]
, and
\[
e_3 = (2, -1, 1).
\]
\textbf{Solution}: Set up the relation:
\[
(1, -2, 5) = c_1(1, 1, 1) + c_2(1, 2, 3) + c_3(2, -1, 1)
\]
This gives the system of equations (grouping by components):
\[
c_1 + c_2 + 2c_3 = 1
\]
\[
c_1 + 2c_2 - c_3 = -2
\]
\[
c_1 + 3c_2 + c_3 = 5
\]
Write the augmented matrix and apply row operations to reach row echelon
form:
\[
\begin{bmatrix}
1 & 1 & 2 & \lvert  & 1 \\
1 & 2 & -1 &  \rvert & -2 \\
1 & 3 & 1 & \lvert  & 5
\end{bmatrix}
\]
Apply \(R_2 \to R_2 - R_1\) and \(R_3 \to R_3 - R_1\):
\[
\begin{bmatrix}
1 & 1 & 2 & \lvert  & 1 \\
0 & 1 & -3 &  \rvert & -3 \\
0 & 2 & -1 & \lvert  & 4
\end{bmatrix}
\]
Apply \(R_3 \to R_3 - 2R_2\):
\[
\begin{bmatrix}
1 & 1 & 2 & \lvert  & 1 \\
0 & 1 & -3 &  \rvert & -3 \\
0 & 0 & 5 & \lvert  & 10
\end{bmatrix}
\]
Solve by back substitution: - From \(R_3\):
\[
5c_3 = 10 \implies c_3 = 2
\]
\begin{itemize}
\tightlist
\item
  From \(R_2\):
\end{itemize}
\[
c_2 - 3(2) = -3 \implies c_2 = 3
\]
\begin{itemize}
\tightlist
\item
  From \(R_1\):
\end{itemize}
\[
c_1 + 3 + 2(2) = 1 \implies c_1 = -6
\]
The linear combination is:
\[
v = -6e_1 + 3e_2 + 2e_3
\]

\subsection{Worked Example 2: Polynomial Vectors (PYQ
2024)}\label{worked-example-2-polynomial-vectors-pyq-2024}

\textbf{Problem}: Express the polynomial
\[
v = 3t^2 + 5t - 5
\]
as a linear combination of the polynomials
\[
P_1 = t^2 + 2t + 1,
\]
\[
P_2 = 2t^2 + 5t + 4
\]
, and
\[
P_3 = t^2 + 3t + 6.
\]
\textbf{Solution}: Set up the relation:
\[
v = c_1 P_1 + c_2 P_2 + c_3 P_3
\]
\[
3t^2 + 5t - 5 = c_1(t^2 + 2t + 1) + c_2(2t^2 + 5t + 4) + c_3(t^2 + 3t + 6)
\]
Expand and group terms by powers of \(t\):
\[
3t^2 + 5t - 5 = (c_1 + 2c_2 + c_3)t^2 + (2c_1 + 5c_2 + 3c_3)t + (c_1 + 4c_2 + 6c_3)
\]
Match the coefficients to form a system of equations: 1.
\[
c_1 + 2c_2 + c_3 = 3
\]
\begin{enumerate}
\def\labelenumi{\arabic{enumi}.}
\setcounter{enumi}{1}
\tightlist
\item
\end{enumerate}
\[
2c_1 + 5c_2 + 3c_3 = 5
\]
\begin{enumerate}
\def\labelenumi{\arabic{enumi}.}
\setcounter{enumi}{2}
\tightlist
\item
\end{enumerate}
\[
c_1 + 4c_2 + 6c_3 = -5
\]
You can solve this using substitution or matrix row reduction. Let's use
substitution from Eq 1 (
\[
c_1 = 3 - 2c_2 - c_3
\]
): Substitute into Eq 2:
\[
2(3 - 2c_2 - c_3) + 5c_2 + 3c_3 = 5 \implies 6 - 4c_2 - 2c_3 + 5c_2 + 3c_3 = 5 \implies c_2 + c_3 = -1
\]
Substitute into Eq 3:
\[
(3 - 2c_2 - c_3) + 4c_2 + 6c_3 = -5 \implies 3 + 2c_2 + 5c_3 = -5 \implies 2c_2 + 5c_3 = -8
\]
Now solve the \(2 \times 2\) system:
\[
c_2 = -1 - c_3.
\]
\[
2(-1 - c_3) + 5c_3 = -8 \implies -2 - 2c_3 + 5c_3 = -8 \implies 3c_3 = -6 \implies c_3 = -2
\]
Find \(c_2\):
\[
c_2 = -1 - (-2) = 1.
\]
Find \(c_1\):
\[
c_1 = 3 - 2(1) - (-2) = 3.
\]
The linear combination is:
\[
v = 3P_1 + P_2 - 2P_3
\]

\subsection{Exam Patterns}\label{exam-patterns-32}

\begin{itemize}
\tightlist
\item
  When given polynomials, simply treat their coefficients like a
  standard vector (e.g., \(3t^2 + 5t - 5\) becomes \((3, 5, -5)\)). The
  math is identical to coordinate vector problems.
\item
  If the system of equations is inconsistent (e.g., you get \(0 = 5\)
  during row reduction), then it is \emph{not possible} to write the
  vector as a linear combination of the given set.
\end{itemize}

\section{C05: Linear Dependence and
Independence}\label{c05-linear-dependence-and-independence}

\begin{quote}
\textbf{Exam Weight}: Tier 1 \textbar{} \textbf{Appeared in}: 2018,
2020, 2021, 2024 \textbar{} \textbf{Typical Marks}: 4--6
\end{quote}

A set of vectors is \textbf{linearly independent} if no vector in the
set can be written as a linear combination of the others. If at least
one vector can be, the set is \textbf{linearly dependent}.

\subsection{Definitions and The Determinant
Test}\label{definitions-and-the-determinant-test}

\textbf{Linearly Independent}: A set of vectors
\[
\lbrace v_1, v_2, \dots, v_n\rbrace
\]
is linearly independent if the equation
\[
c_1 v_1 + c_2 v_2 + \dots + c_n v_n = 0
\]
has \emph{only the trivial solution} (
\[
c_1 = c_2 = \dots = c_n = 0
\]
). \textbf{Linearly Dependent}: The set is dependent if the equation has
a solution where at least one coefficient is not zero.

\subsubsection{The Determinant Test}\label{the-determinant-test}

If you have exactly \(n\) vectors in \(\mathbb{R}^n\), you can form an
\(n \times n\) matrix with the vectors as its rows (or columns). - If
Determinant \(\neq 0 \implies\) Linearly \textbf{Independent}. - If
Determinant \(= 0 \implies\) Linearly \textbf{Dependent}.

*(Note: If you are given a theoretical proof question (PYQ 2024) asking
you to show why this determinant rule is true, set
\[
a\vec{A} + b\vec{B} + c\vec{C} = 0
\]
, convert it to a homogeneous system, and state that a homogeneous
system only has the trivial solution if the determinant of the
coefficient matrix is non-zero).*

\subsection{Worked Example 1: Finding the Dependence Relation (PYQ
2018)}\label{worked-example-1-finding-the-dependence-relation-pyq-2018}

\textbf{Problem}: Show that
\[
\vec{A} = 2\hat{i} + \hat{j} - 3\hat{k},
\]
\[
\vec{B} = \hat{i} - 4\hat{k},
\]
\[
\vec{C} = 4\hat{i} + 3\hat{j} - \hat{k}
\]
are linearly dependent, and determine a relation between them.

\textbf{Solution}: \textbf{Step 1: Check Determinant}
\[
D =
\begin{vmatrix}
2 & 1 & -3 \\
1 & 0 & -4 \\
4 & 3 & -1
\end{vmatrix} = 2(0 - (-12)) - 1(-1 - (-16)) - 3(3 - 0) = 24 - 15 - 9 = 0
\]
Since \(D = 0\), the vectors are linearly dependent.

\textbf{Step 2: Find the Relation} Set
\[
x\vec{A} + y\vec{B} + z\vec{C} = 0:
\]
\[
x(2, 1, -3) + y(1, 0, -4) + z(4, 3, -1) = (0,0,0)
\]
This yields the system: 1. \(2x + y + 4z = 0\) 2.
\(x + 3z = 0 \implies x = -3z\) 3. \(-3x - 4y - z = 0\)

Substitute \(x = -3z\) into Equation 1:
\[
2(-3z) + y + 4z = 0 \implies -6z + y + 4z = 0 \implies y = 2z
\]
Let \(z = 1\) (we can pick any non-zero value). This gives \(x = -3\)
and \(y = 2\). Substitute the coefficients back into the original
relation:
\[
-3\vec{A} + 2\vec{B} + 1\vec{C} = 0 \implies \vec{C} = 3\vec{A} - 2\vec{B}
\]

\subsection{Worked Example 2: Abstract Vectors (PYQ 2020,
2021)}\label{worked-example-2-abstract-vectors-pyq-2020-2021}

\textbf{Problem}: Suppose vectors \(u, v, w\) are linearly independent.
Show that \(u + v, u - v, u - 2v + w\) are also linearly independent.

\textbf{Solution}: Set a linear combination equal to the zero vector:
\[
c_1(u + v) + c_2(u - v) + c_3(u - 2v + w) = 0
\]
Group the terms by the original vectors \(u, v, w\):
\[
(c_1 + c_2 + c_3)u + (c_1 - c_2 - 2c_3)v + c_3w = 0
\]
Since we are \emph{given} that \(u, v, w\) are linearly independent, the
only way their linear combination equals 0 is if their coefficients are
exactly 0. Thus: 1.
\[
c_1 + c_2 + c_3 = 0
\]
\begin{enumerate}
\def\labelenumi{\arabic{enumi}.}
\setcounter{enumi}{1}
\tightlist
\item
\end{enumerate}
\[
c_1 - c_2 - 2c_3 = 0
\]
\begin{enumerate}
\def\labelenumi{\arabic{enumi}.}
\setcounter{enumi}{2}
\tightlist
\item
\end{enumerate}
\[
c_3 = 0
\]
Substitute
\[
c_3 = 0
\]
into the first two equations:
\[
c_1 + c_2 = 0
\]
\[
c_1 - c_2 = 0
\]
Adding them gives
\[
2c_1 = 0 \implies c_1 = 0
\]
. This means
\[
c_2 = 0.
\]
Since the only solution is
\[
c_1 = c_2 = c_3 = 0
\]
, the new vectors are linearly independent. \(\blacksquare\)

\subsection{Exam Patterns}\label{exam-patterns-33}

\begin{itemize}
\tightlist
\item
  If you are asked to test 3 vectors in \(\mathbb{R}^3\), immediately
  calculate the determinant.
\item
  If asked to ``find a relation'', you must set up the \(x, y, z\)
  system, express two variables in terms of the third, arbitrarily set
  the third to \(1\), and write the final formula.
\item
  Abstract linear independence proofs (like Example 2) are extremely
  common and very easy to score full marks on if you group coefficients
  properly.
\end{itemize}

\section{C06: Basis and Dimension}\label{c06-basis-and-dimension}

\begin{quote}
\textbf{Exam Weight}: Tier 1 \textbar{} \textbf{Appeared in}: 2018,
2020, 2021, 2023, 2024 (twice) \textbar{} \textbf{Typical Marks}: 4--5
\end{quote}

A basis is the most efficient set of vectors needed to define a space.
It contains enough vectors to reach everywhere in the space (span), but
no unnecessary vectors (linear independence).

\subsection{Definitions}\label{definitions-1}

\begin{itemize}
\tightlist
\item
  \textbf{Basis}: A set of vectors
\end{itemize}
\[
B = \lbrace v_1, v_2, \dots, v_n\rbrace
\]
is a basis of a vector space \(V\) if: 1. The set \(B\) is linearly
independent. 2. The set \(B\) spans \(V\) (any vector in \(V\) can be
written as a linear combination of vectors in \(B\)). -
\textbf{Dimension}: The dimension of a vector space is exactly equal to
the number of vectors in its basis.

\subsection{Finding a Basis and Dimension (The Row Reduction
Method)}\label{finding-a-basis-and-dimension-the-row-reduction-method}

A very common exam question gives you a set of vectors and asks you to
find a basis and the dimension of the subspace they span.

\textbf{The Algorithm}: 1. Write the given vectors as the \textbf{rows}
of a matrix. 2. Apply elementary row operations to reduce the matrix to
\textbf{Row Echelon Form} (zeros below the leading ones). 3. The
\textbf{non-zero rows} of the resulting matrix form the basis. 4. The
\textbf{number} of non-zero rows is the dimension.

\emph{(Note: If a row becomes entirely zeros, it means that original
vector was redundant and was a linear combination of the others).}

\subsection{Worked Example 1: Vectors in (PYQ
2018)}\label{worked-example-1-vectors-in-pyq-2018}

\textbf{Problem}: Let \(W\) be the subspace of \(\mathbb{R}^4\)
generated by \((1, -2, 5, -3)\), \((2, 3, 1, -4)\), and
\((3, 8, -3, -5)\). Find a basis and the dimension of \(W\).

\textbf{Solution}: Write vectors as rows:
\[
\begin{bmatrix}
1 & -2 & 5 & -3 \\
2 & 3 & 1 & -4 \\
3 & 8 & -3 & -5
\end{bmatrix}
\]
Apply row operations to eliminate the first column below row 1
(\(R_2 \to R_2 - 2R_1\) and \(R_3 \to R_3 - 3R_1\)):
\[
\begin{bmatrix}
1 & -2 & 5 & -3 \\
0 & 7 & -9 & 2 \\
0 & 14 & -18 & 4
\end{bmatrix}
\]
Eliminate the second column below row 2 (\(R_3 \to R_3 - 2R_2\)):
\[
\begin{bmatrix}
1 & -2 & 5 & -3 \\
0 & 7 & -9 & 2 \\
0 & 0 & 0 & 0
\end{bmatrix}
\]
The non-zero rows form the basis:
\[
\text{Basis} = \lbrace (1, -2, 5, -3), (0, 7, -9, 2) \rbrace
\]
Since there are 2 vectors in the basis:
\[
\text{Dimension} = 2
\]

\subsection{Worked Example 2: Polynomial Spaces (PYQ 2020,
2023)}\label{worked-example-2-polynomial-spaces-pyq-2020-2023}

\textbf{Problem}: Find a basis and dimension of the space generated by
the polynomials:
\[
v_1 = t^3 - 2t^2 + 4t + 1
\]
\[
v_2 = 2t^3 - 3t^2 + 9t - 1
\]
\[
v_3 = t^3 + 6t - 5
\]
\[
v_4 = 2t^3 - 5t^2 + 7t + 5
\]
\textbf{Solution}: Represent each polynomial as a coordinate vector in
the standard basis \(\lbrace t^3, t^2, t, 1\rbrace\) and form a matrix:
\[
\begin{bmatrix}
1 & -2 & 4 & 1 \\
2 & -3 & 9 & -1 \\
1 & 0 & 6 & -5 \\
2 & -5 & 7 & 5
\end{bmatrix}
\]
Apply \(R_2 \to R_2 - 2R_1\), \(R_3 \to R_3 - R_1\),
\(R_4 \to R_4 - 2R_1\):
\[
\begin{bmatrix}
1 & -2 & 4 & 1 \\
0 & 1 & 1 & -3 \\
0 & 2 & 2 & -6 \\
0 & -1 & -1 & 3
\end{bmatrix}
\]
Apply \(R_3 \to R_3 - 2R_2\) and \(R_4 \to R_4 + R_2\):
\[
\begin{bmatrix}
1 & -2 & 4 & 1 \\
0 & 1 & 1 & -3 \\
0 & 0 & 0 & 0 \\
0 & 0 & 0 & 0
\end{bmatrix}
\]
Convert the non-zero rows back into polynomials to form the basis:
\[
\text{Basis} = \lbrace t^3 - 2t^2 + 4t + 1, \quad t^2 + t - 3 \rbrace
\]
Since there are 2 polynomials in the basis:
\[
\text{Dimension} = 2
\]

\subsection{Exam Patterns}\label{exam-patterns-34}

\begin{itemize}
\tightlist
\item
  This exact row-reduction algorithm appeared 5 times in the past 6
  years. It is highly reliable.
\item
  If asked ``Determine whether these 4 vectors form a basis of
  \(\mathbb{R}^4\)'', perform the row reduction. If you end up with 4
  non-zero rows, yes they do. If any row zeros out (meaning dimension
  \textless{} 4), they do not form a basis. (See PYQ 2021/2024).
\end{itemize}

\section{C07: Linear Transformations}\label{c07-linear-transformations}

\begin{quote}
\textbf{Exam Weight}: Tier 3 \textbar{} \textbf{Appeared in}: 2023
\textbar{} \textbf{Typical Marks}: 4
\end{quote}

A linear transformation (or mapping) is a function between two vector
spaces that preserves the operations of vector addition and scalar
multiplication.

\subsection{The Linearity Conditions}\label{the-linearity-conditions}

Let \(V\) and \(U\) be vector spaces over the same field \(K\). A
function \(F: V \to U\) is a \textbf{linear mapping} if for all vectors
\(u, v \in V\) and all scalars \(k \in K\):

\begin{enumerate}
\def\labelenumi{\arabic{enumi}.}
\tightlist
\item
  \textbf{Preservation of Addition}: \(F(u + v) = F(u) + F(v)\)
\item
  \textbf{Preservation of Scalar Multiplication}: \(F(ku) = kF(u)\)
\end{enumerate}

\emph{Shortcut}: These two conditions can be combined into one:
\(F(au + bv) = aF(u) + bF(v)\).

\textbf{Crucial Property}: Any linear transformation MUST map the zero
vector of the domain to the zero vector of the codomain. If
\(F(0_V) \neq 0_U\), the mapping is definitely not linear.

\subsection{Worked Example: Verifying Linearity (PYQ
2023)}\label{worked-example-verifying-linearity-pyq-2023}

\textbf{Problem}: Identify if the following mappings \(F\) are linear or
not linear: (i)
\[
F : \mathbb{R}^3 \to \mathbb{R}
\]
defined by \(F(x, y, z) = 2x - 3y + 4z\) (ii)
\[
F : \mathbb{R}^2 \to \mathbb{R}^3
\]
defined by \(F(x, y) = (x+1, 2y, x+y)\)

\textbf{Solution}:

\textbf{(i) Mapping \(F(x, y, z) = 2x - 3y + 4z\)} Let
\[
u = (x_1, y_1, z_1)
\]
and
\[
v = (x_2, y_2, z_2).
\]
\emph{Check Addition}:
\[
F(u + v) = F(x_1 + x_2, y_1 + y_2, z_1 + z_2) = 2(x_1 + x_2) - 3(y_1 + y_2) + 4(z_1 + z_2)
\]
\[
F(u + v) = (2x_1 - 3y_1 + 4z_1) + (2x_2 - 3y_2 + 4z_2) = F(u) + F(v)
\]
\emph{Check Scalar Multiplication}:
\[
F(ku) = F(kx_1, ky_1, kz_1) = 2(kx_1) - 3(ky_1) + 4(kz_1) = k(2x_1 - 3y_1 + 4z_1) = kF(u)
\]
Both conditions hold. The mapping is \textbf{linear}.

\textbf{(ii) Mapping \(F(x, y) = (x+1, 2y, x+y)\)} We use the
zero-vector shortcut. A linear mapping must map the zero vector
\((0,0)\) to the zero vector \((0,0,0)\).
\[
F(0, 0) = (0 + 1, 2(0), 0 + 0) = (1, 0, 0) \neq (0, 0, 0)
\]
Since \(F(0_V) \neq 0_U\), the mapping is \textbf{not linear}.

\subsection{Exam Patterns}\label{exam-patterns-35}

\begin{itemize}
\tightlist
\item
  When given a function with a constant addition term (like \(x+1\)), it
  is automatically non-linear because it is an ``affine'' shift, not a
  linear one. Use the \(F(0) \neq 0\) trick to prove it in one line.
\item
  For functions without constants (like \(2x-3y\)), you must grind out
  the full \(F(u+v)\) and \(F(ku)\) proofs.
\end{itemize}

\section{C08: Kernel, Image, Rank, and
Nullity}\label{c08-kernel-image-rank-and-nullity}

\begin{quote}
\textbf{Exam Weight}: Tier 1-2 \textbar{} \textbf{Appeared in}: 2019
\textbar{} \textbf{Typical Marks}: 6
\end{quote}

Every linear mapping \(F: V \to U\) is associated with two crucial
subspaces: the Kernel (inside the domain \(V\)) and the Image (inside
the codomain \(U\)).

\subsection{Definitions}\label{definitions-2}

\begin{itemize}
\tightlist
\item
  \textbf{Kernel (Null Space)}: The set of all vectors in \(V\) that map
  exactly to the zero vector in \(U\).
\end{itemize}
\[
\text{Ker}(F) = \lbrace v \in V \mid F(v) = 0_U \rbrace
\]
\begin{itemize}
\tightlist
\item
  \textbf{Image (Range Space)}: The set of all vectors in \(U\) that are
  ``hit'' by the mapping (i.e., they have at least one pre-image in
  \(V\)).
\end{itemize}
\[
\text{Im}(F) = \lbrace u \in U \mid \exists v \in V \text{ such that } F(v) = u \rbrace
\]

\subsection{Must-Know Proof (PYQ
2019)}\label{must-know-proof-pyq-2019-1}

\textbf{Problem}: Let \(F : V \to U\) be a linear mapping. Show that the
Kernel of \(F\) is a subspace of \(V\) and the Image of \(F\) is a
subspace of \(U\).

\textbf{Part 1: Prove \(\text{Ker}(F)\) is a subspace of \(V\)} We must
verify the three subspace properties for \(\text{Ker}(F)\): 1.
\textbf{Zero Vector}: Since \(F\) is linear,
\[
F(0_V) = 0_U
\]
. Therefore, \(0_V \in \text{Ker}(F)\). 2. \textbf{Closure under
Addition}: Let \(v_1, v_2 \in \text{Ker}(F)\). By definition,
\[
F(v_1) = 0_U
\]
and
\[
F(v_2) = 0_U.
\]
\[
F(v_1 + v_2) = F(v_1) + F(v_2) = 0_U + 0_U = 0_U
\]
. Thus, \(v_1 + v_2 \in \text{Ker}(F)\). 3. \textbf{Closure under Scalar
Multiplication}: Let \(v \in \text{Ker}(F)\) and let \(k\) be a scalar.
\[
F(kv) = kF(v) = k(0_U) = 0_U
\]
. Thus, \(kv \in \text{Ker}(F)\). Since all three hold,
\(\text{Ker}(F)\) is a subspace. \(\blacksquare\)

\textbf{Part 2: Prove \(\text{Im}(F)\) is a subspace of \(U\)} We must
verify the three subspace properties for \(\text{Im}(F)\): 1.
\textbf{Zero Vector}: Since
\[
F(0_V) = 0_U
\]
, the zero vector \(0_U\) has a pre-image. Thus,
\(0_U \in \text{Im}(F)\). 2. \textbf{Closure under Addition}: Let
\(u_1, u_2 \in \text{Im}(F)\). Then there exist \(v_1, v_2 \in V\) such
that
\[
F(v_1) = u_1
\]
and
\[
F(v_2) = u_2.
\]
\[
u_1 + u_2 = F(v_1) + F(v_2) = F(v_1 + v_2).
\]
Since \(v_1 + v_2 \in V\), the sum \(u_1 + u_2\) is the image of a
vector in \(V\). Thus, \(u_1 + u_2 \in \text{Im}(F)\). 3.
\textbf{Closure under Scalar Multiplication}: Let \(u \in \text{Im}(F)\)
and let \(k\) be a scalar. There exists \(v \in V\) such that
\(F(v) = u\). \(ku = kF(v) = F(kv)\). Since \(kv \in V\), \(ku\) is the
image of a vector in \(V\). Thus, \(ku \in \text{Im}(F)\). Since all
three hold, \(\text{Im}(F)\) is a subspace. \(\blacksquare\)

\subsection{Rank-Nullity Theorem}\label{rank-nullity-theorem}

The dimensions of these two subspaces are given special names: -
\textbf{Nullity}: The dimension of the Kernel. - \textbf{Rank}: The
dimension of the Image.

The Rank-Nullity Theorem states that for any linear mapping
\(F: V \to U\), the dimensions perfectly balance out to equal the
dimension of the starting domain \(V\):
\[
\dim(\text{Ker}(F)) + \dim(\text{Im}(F)) = \dim(V)
\]
\[
\text{Nullity} + \text{Rank} = \dim(V)
\]

\subsection{Exam Patterns}\label{exam-patterns-36}

\begin{itemize}
\tightlist
\item
  The proof that the Kernel and Image are subspaces is a classic 6-mark
  theoretical question. Memorize the 3-step subspace check for both.
\item
  If asked to find the Kernel of a matrix mapping, set the matrix times
  a vector \(\vec{x}\) to the zero vector:
\end{itemize}
\[
A\vec{x} = 0
\]
Then solve for \(\vec{x}\) using row reduction.

\section{C09: Matrix Representation and Change of
Basis}\label{c09-matrix-representation-and-change-of-basis}

\begin{quote}
\textbf{Exam Weight}: Tier 3 \textbar{} \textbf{Appeared in}: Curriculum
Reference \textbar{} \textbf{Typical Marks}: 3--4
\end{quote}

Every linear transformation between finite-dimensional vector spaces can
be entirely represented by a matrix. Changing the basis of the vector
space changes the matrix representation in a predictable way.

\subsection{Matrix Representation of a Linear
Operator}\label{matrix-representation-of-a-linear-operator}

Let \(T: V \to V\) be a linear operator on a vector space \(V\) with
dimension \(n\). Let
\[
B = \lbrace v_1, v_2, \dots, v_n\rbrace
\]
be a basis for \(V\).

When you apply the transformation \(T\) to the first basis vector
\(v_1\), the result \(T(v_1)\) is still a vector in \(V\). Therefore, it
can be written as a linear combination of the basis vectors:
\[
T(v_1) = a_{11}v_1 + a_{21}v_2 + \dots + a_{n1}v_n
\]
The coefficients
\[
(a_{11}, a_{21}, \dots, a_{n1})
\]
form the \textbf{first column} of the matrix representation \([T]_B\).
You repeat this for every basis vector. The resulting \(n \times n\)
matrix is the \textbf{Matrix Representation} of the linear operator.

\subsection{Change of Basis Matrix}\label{change-of-basis-matrix}

Suppose you have two different bases for the same vector space \(V\): an
old basis \(B\) and a new basis \(B'\). You can write each vector of the
\emph{new} basis \(B'\) as a linear combination of the \emph{old} basis
\(B\).

Taking the column coefficients of these combinations creates the
\textbf{Change of Basis Matrix} (or Transition Matrix), denoted as
\(P\). - The matrix \(P\) converts coordinates from the new basis to the
old basis:
\[
[x]_B = P[x]_{B'}.
\]
\begin{itemize}
\tightlist
\item
  The inverse matrix \(P^{-1}\) converts coordinates from the old basis
  to the new basis:
\end{itemize}
\[
[x]_{B'} = P^{-1}[x]_B.
\]

\subsection{Similarity and Matrix
Equivalency}\label{similarity-and-matrix-equivalency}

When you change the basis of the vector space from \(B\) to \(B'\), the
matrix representation of the linear operator \(T\) also changes.

If
\[
[T]_B = A
\]
(the matrix in the old basis) and
\[
[T]_{B'} = C
\]
(the matrix in the new basis), they are related by the formula:
\[
C = P^{-1} A P
\]
Two matrices \(A\) and \(C\) that satisfy this relationship for some
invertible matrix \(P\) are called \textbf{Similar Matrices}. Similar
matrices represent the \emph{exact same linear transformation}, just
viewed from the perspective of two different coordinate systems. Because
they represent the same transformation, similar matrices share
identical: - Determinants - Traces - Eigenvalues

\subsection{Exam Patterns}\label{exam-patterns-37}

\begin{itemize}
\tightlist
\item
  While deep calculations involving Change of Basis rarely appear in
  short exams, understanding the formula
\end{itemize}
\[
C = P^{-1} A P
\]
is critical for diagonalization questions (which appear in the Matrix
Theory sections).

\section{C10: Inner Product Spaces}\label{c10-inner-product-spaces}

\begin{quote}
\textbf{Exam Weight}: Tier 2 \textbar{} \textbf{Appeared in}: Assignment
/ Final \textbar{} \textbf{Typical Marks}: 4--5
\end{quote}

An inner product space is a vector space equipped with a way to multiply
vectors together, resulting in a scalar. This generalizes the concept of
the dot product, allowing us to define angles and lengths in abstract
spaces.

\subsection{Inner Product Axioms}\label{inner-product-axioms}

Let \(V\) be a vector space over the complex numbers \(\mathbb{C}\) (or
real numbers \(\mathbb{R}\)). An inner product is a function
\(\langle x, y \rangle\) that maps two vectors to a scalar, satisfying
the following axioms for all \(x, y, z \in V\) and scalar \(\alpha\):

\begin{enumerate}
\def\labelenumi{\arabic{enumi}.}
\tightlist
\item
  \textbf{Conjugate Symmetry}:
  \(\langle x, y \rangle = \overline{\langle y, x \rangle}\) \emph{(If
  the field is real, this is simply commutativity:
  \(\langle x, y \rangle = \langle y, x \rangle\))}
\item
  \textbf{Linearity in the First Argument}:
  \(\langle x + y, z \rangle = \langle x, z \rangle + \langle y, z \rangle\)
  \(\langle \alpha x, y \rangle = \alpha \langle x, y \rangle\)
\item
  \textbf{Positive-Definiteness}: \(\langle x, x \rangle \geq 0\), and
  \(\langle x, x \rangle = 0\) if and only if \(x = 0\).
\end{enumerate}

\subsection{Must-Know Proof 1: The Induced
Norm}\label{must-know-proof-1-the-induced-norm}

\textbf{Problem}: Prove that every inner product space is a normed
space.

\textbf{Proof}: Let \(V\) be an inner product space. We define the
induced norm as
\[
\|x\| = \sqrt{\langle x, x \rangle}
\]
. We must prove that this satisfies the three axioms of a norm.

\begin{enumerate}
\def\labelenumi{\arabic{enumi}.}
\tightlist
\item
  \textbf{Positivity}: By the positive-definite axiom of inner products,
  \(\langle x, x \rangle \geq 0\). Thus,
\end{enumerate}
\[
\|x\| = \sqrt{\langle x, x \rangle} \geq 0
\]
. Additionally, \(\langle x, x \rangle = 0 \iff x = 0\), so
\[
\|x\| = 0 \iff x = 0.
\]
\begin{enumerate}
\def\labelenumi{\arabic{enumi}.}
\setcounter{enumi}{1}
\tightlist
\item
  \textbf{Absolute Homogeneity}:
\end{enumerate}
\[
\|\alpha x\| = \sqrt{\langle \alpha x, \alpha x \rangle} = \sqrt{\alpha \overline{\alpha} \langle x, x \rangle} = \sqrt{\lvert \alpha \rvert^2 \langle x, x \rangle} = \lvert \alpha \rvert \sqrt{\langle x, x \rangle} = \lvert \alpha \rvert \|x\|
\]
\begin{enumerate}
\def\labelenumi{\arabic{enumi}.}
\setcounter{enumi}{2}
\tightlist
\item
  \textbf{Triangle Inequality}:
\end{enumerate}
\[
\|x + y\|^2 = \langle x + y, x + y \rangle = \langle x, x \rangle + \langle x, y \rangle + \langle y, x \rangle + \langle y, y \rangle
\]
\[
\|x + y\|^2 = \|x\|^2 + \langle x, y \rangle + \overline{\langle x, y \rangle} + \|y\|^2 = \|x\|^2 + 2\text{Re}(\langle x, y \rangle) + \|y\|^2
\]
By the Cauchy-Schwarz inequality (
\[
\lvert \langle x, y \rangle \rvert \leq \|x\| \|y\|
\]
) and knowing
\[
\text{Re}(z) \leq \lvert z \rvert:
\]
\[
2\text{Re}(\langle x, y \rangle) \leq 2\lvert \langle x, y \rangle \rvert \leq 2\|x\|\|y\|
\]
\begin{verbatim}
Substituting this back:
\end{verbatim}
\[
\|x + y\|^2 \leq \|x\|^2 + 2\|x\|\|y\| + \|y\|^2 = (\|x\| + \|y\|)^2
\]
\begin{verbatim}
Taking the square root gives $\|x + y\| \leq \|x\| + \|y\|$.
\end{verbatim}

Since all three axioms hold, the induced norm makes the inner product
space a normed space. \(\blacksquare\)

\subsection{Must-Know Proof 2: The Parallelogram
Law}\label{must-know-proof-2-the-parallelogram-law}

\textbf{Problem}: Prove the parallelogram law in an inner product space.

\textbf{Proof}: The parallelogram law states that the sum of the squares
of the lengths of the four sides of a parallelogram equals the sum of
the squares of the lengths of the two diagonals:
\[
\|x + y\|^2 + \|x - y\|^2 = 2\|x\|^2 + 2\|y\|^2
\]
Expand the left-hand side using the inner product:
\[
\|x + y\|^2 = \langle x + y, x + y \rangle = \langle x, x \rangle + \langle x, y \rangle + \langle y, x \rangle + \langle y, y \rangle
\]
\[
= \|x\|^2 + \langle x, y \rangle + \langle y, x \rangle + \|y\|^2
\]
Expand the second term:
\[
\|x - y\|^2 = \langle x - y, x - y \rangle = \langle x, x \rangle - \langle x, y \rangle - \langle y, x \rangle + \langle y, y \rangle
\]
\[
= \|x\|^2 - \langle x, y \rangle - \langle y, x \rangle + \|y\|^2
\]
Add the two equations together:
\[
\|x + y\|^2 + \|x - y\|^2 = (\|x\|^2 + \langle x, y \rangle + \langle y, x \rangle + \|y\|^2) + (\|x\|^2 - \langle x, y \rangle - \langle y, x \rangle + \|y\|^2)
\]
The cross terms \(\langle x, y \rangle\) and \(-\langle x, y \rangle\)
cancel out, as do \(\langle y, x \rangle\) and
\(-\langle y, x \rangle\).
\[
\|x + y\|^2 + \|x - y\|^2 = 2\|x\|^2 + 2\|y\|^2
\]
\(\blacksquare\)

\subsection{Exam Patterns}\label{exam-patterns-38}

\begin{itemize}
\tightlist
\item
  If asked to prove something is an inner product space, you just need
  to grind through the 3 axioms (linearity, symmetry,
  positive-definite).
\item
  The Parallelogram Law proof is incredibly straightforward and a very
  easy 5 marks if it appears.
\end{itemize}

\section{C11: Normed Linear Spaces \& Banach
Spaces}\label{c11-normed-linear-spaces-banach-spaces}

\begin{quote}
\textbf{Exam Weight}: Tier 1-2 \textbar{} \textbf{Appeared in}:
Assignment / Final \textbar{} \textbf{Typical Marks}: 3--5
\end{quote}

While an inner product gives a vector space ``angles'', a norm gives a
vector space ``length''. When a normed space is mathematically
``complete'', it is called a Banach Space.

\subsection{Normed Linear Space
Definitions}\label{normed-linear-space-definitions}

A normed linear space is a vector space \(V\) equipped with a norm
\(\|\cdot\| : V \to \mathbb{R}\), which assigns a non-negative real
length to each vector. It must satisfy three properties for all
\(x, y \in V\) and scalars \(\alpha\): 1. \textbf{Positivity}:
\(\|x\| \geq 0\), and
\[
\|x\| = 0 \iff x = 0.
\]
\begin{enumerate}
\def\labelenumi{\arabic{enumi}.}
\setcounter{enumi}{1}
\tightlist
\item
  \textbf{Absolute Homogeneity}:
\end{enumerate}
\[
\|\alpha x\| = \lvert \alpha \rvert \|x\|.
\]
\begin{enumerate}
\def\labelenumi{\arabic{enumi}.}
\setcounter{enumi}{2}
\tightlist
\item
  \textbf{Triangle Inequality}: \(\|x + y\| \leq \|x\| + \|y\|\).
\end{enumerate}

\textbf{Example Verification}: Verify whether
\[
\|x\| = \lvert x_1 \rvert + \lvert x_2 \rvert
\]
defines a norm on \(\mathbb{R}^2\). 1. \textbf{Positivity}:
\(\lvert x_1 \rvert \geq 0\) and \(\lvert x_2 \rvert \geq 0\), so
\[
\lvert x_1 \rvert + \lvert x_2 \rvert \geq 0
\]
. If
\[
\|x\| = 0
\]
, then
\[
\lvert x_1 \rvert=0
\]
and
\[
\lvert x_2 \rvert=0
\]
, so \(x=(0,0)\). 2. \textbf{Homogeneity}:
\[
\|\alpha x\| = \lvert \alpha x_1 \rvert + \lvert \alpha x_2 \rvert = \lvert \alpha \rvert(\lvert x_1 \rvert + \lvert x_2 \rvert) = \lvert \alpha \rvert\|x\|.
\]
\begin{enumerate}
\def\labelenumi{\arabic{enumi}.}
\setcounter{enumi}{2}
\tightlist
\item
  \textbf{Triangle Inequality}:
\end{enumerate}
\[
\|x + y\| = \lvert x_1 + y_1 \rvert + \lvert x_2 + y_2 \rvert \leq (\lvert x_1 \rvert + \lvert y_1 \rvert) + (\lvert x_2 \rvert + \lvert y_2 \rvert) = \|x\| + \|y\|.
\]
Since all three hold, it is a valid norm (the \(L^1\) norm).

\subsection{Banach Space Definitions}\label{banach-space-definitions}

A \textbf{Banach space} is defined as a \textbf{complete} normed linear
space. ``Complete'' means that every Cauchy sequence of vectors in the
space successfully converges to a limit that is \emph{also contained
within the space}. There are no ``missing holes'' in the space.

\subsubsection{Example of a Banach
Space}\label{example-of-a-banach-space}

The space \(\ell^\infty\), which consists of all bounded infinite
sequences of complex numbers
\[
x = (x_1, x_2, \dots)
\]
, is equipped with the supremum norm:
\[
\|x\|_\infty = \sup_{i} \lvert x_i \rvert
\]
Because \(\mathbb{C}\) is complete, Cauchy sequences of coordinates
converge to bounded limits within \(\ell^\infty\).

\subsubsection{Example of an INCOMPLETE Normed
Space}\label{example-of-an-incomplete-normed-space}

The space \(P[a, b]\) of all polynomial functions on \([a, b]\),
equipped with the supremum norm. It is not complete because a Cauchy
sequence of polynomials (like the Taylor series expansion for \(e^t\))
will converge to a limit function (\(e^t\)) that is continuous but
\emph{not} a polynomial. Since the limit escapes the space of
polynomials, the space is not complete.

\subsection{Must-Know Proof: Finite-Dimensional
Spaces}\label{must-know-proof-finite-dimensional-spaces}

\textbf{Problem}: Prove that every finite-dimensional normed linear
space is complete.

\textbf{Proof}: Let \(V\) be an \(n\)-dimensional normed linear space.
Let \(\lbrace e_1, \dots, e_n\rbrace\) be a basis. Any vector \(x\) can
be uniquely written as
\[
x = \sum_{i=1}^n \alpha_i e_i.
\]
Define a reference norm:
\[
\|x\|_1 = \sum_{i=1}^n \lvert \alpha_i \rvert
\]
. In a finite-dimensional space, all norms are equivalent, meaning there
exist constants \(c, C > 0\) such that
\[
c\|x\|_1 \leq \|x\| \leq C\|x\|_1.
\]
Let \((x^{(m)})\) be a Cauchy sequence in \(V\) with respect to
\(\|\cdot\|\). Because the norms are equivalent, it is also Cauchy with
respect to \(\|\cdot\|_1\). Let
\[
x^{(m)} = \sum_{i=1}^n \alpha_i^{(m)} e_i
\]
. Since \((x^{(m)})\) is Cauchy in \(\|\cdot\|_1\), for a given
\(\epsilon > 0\), for large \(p, q\):
\[
\|x^{(p)} - x^{(q)}\|_1 = \sum_{i=1}^n \lvert \alpha_i^{(p)} - \alpha_i^{(q)} \rvert < \epsilon
\]
This implies that for each coordinate index \(i\), the sequence
\[
(\alpha_i^{(m)})_{m=1}^\infty
\]
is a Cauchy sequence of scalars. Since \(\mathbb{R}\) (and
\(\mathbb{C}\)) is complete, each coordinate sequence converges to a
limit \(\alpha_i\).

Let
\[
x = \sum_{i=1}^n \alpha_i e_i
\]
. Clearly, \(x \in V\). We check convergence:
\[
\|x^{(m)} - x\|_1 = \sum_{i=1}^n \lvert \alpha_i^{(m)} - \alpha_i \rvert \to 0
\]
as \(m \to \infty\). Because norms are equivalent,
\[
\|x^{(m)} - x\| \leq C\|x^{(m)} - x\|_1 \to 0
\]
, meaning the sequence converges to \(x\) in the original norm. Thus,
every Cauchy sequence converges within \(V\), making \(V\) complete.
\(\blacksquare\)

\subsection{Exam Patterns}\label{exam-patterns-39}

\begin{itemize}
\tightlist
\item
  The definitions are highly testable. Remember: Banach = Complete +
  Normed.
\item
  You may be asked for a 3-mark definition and example of a Banach
  space. Provide the \(\ell^\infty\) or \(C[a,b]\) example.
\end{itemize}

\section{C12: Hilbert Spaces}\label{c12-hilbert-spaces}

\begin{quote}
\textbf{Exam Weight}: Tier 1-2 \textbar{} \textbf{Appeared in}:
Assignment / Final \textbar{} \textbf{Typical Marks}: 4--5
\end{quote}

If a Banach space is a complete space with a norm (length), a Hilbert
space goes one step further: it is a complete space with an inner
product (length and angle).

\subsection{Definition}\label{definition-16}

A \textbf{Hilbert Space} is a complete inner product space. - It has an
inner product \(\langle x, y \rangle\). - This inner product induces a
norm
\[
\|x\| = \sqrt{\langle x, y \rangle}.
\]
\begin{itemize}
\tightlist
\item
  The space is complete with respect to the metric defined by this
  induced norm (every Cauchy sequence converges to a limit within the
  space).
\end{itemize}

\subsection{The Hierarchy of Spaces}\label{the-hierarchy-of-spaces}

The spaces you have learned form a strict nested hierarchy: 1.
\textbf{Vector Space}: Has addition and scalar multiplication. 2.
\textbf{Normed Space}: A Vector Space that also has a concept of length
(\(\|x\|\)). 3. \textbf{Banach Space}: A Normed Space that is
``complete'' (no missing limits). 4. \textbf{Inner Product Space}: A
Vector Space that has a dot-product-like operation (angles). \emph{Note:
All inner product spaces are automatically Normed Spaces via the induced
norm, but they are not necessarily complete.} 5. \textbf{Hilbert Space}:
An Inner Product Space that is complete.

\subsubsection{Must-Know Proof: Every Hilbert is a Banach (PYQ /
Assignment)}\label{must-know-proof-every-hilbert-is-a-banach-pyq-assignment}

\textbf{Problem}: Prove that every Hilbert space is a Banach space.

\textbf{Proof}: By definition, a Hilbert space is an inner product space
that is complete with respect to the distance metric induced by its
inner product. The induced metric is based on the induced norm, defined
as
\[
\|x\| = \sqrt{\langle x, x \rangle}.
\]
As proven previously (see C10), the induced norm
\[
\|x\| = \sqrt{\langle x, x \rangle}
\]
satisfies all the axioms of a norm (positivity, absolute homogeneity,
and the triangle inequality via the Cauchy-Schwarz inequality).
Therefore, every inner product space is automatically a \textbf{normed
linear space} when equipped with this induced norm.

Because a Hilbert space is, by definition, \textbf{complete} with
respect to this specific induced norm, it fulfills both conditions of a
Banach space: 1. It is a normed linear space. 2. It is complete. Thus,
every Hilbert space is a Banach space. \(\blacksquare\) \emph{(Note: The
converse is not true. Not all norms are induced by an inner product. If
a norm does not satisfy the Parallelogram Law, it cannot come from an
inner product, so a space with that norm could be Banach but not
Hilbert).}

\subsection{Worked Example: The Space}\label{worked-example-the-space}

\textbf{Problem}: Show that the space \(\ell^2\) is a Hilbert space.

\textbf{Solution Overview}: The space \(\ell^2\) consists of all
infinite sequences of complex numbers
\[
x = (x_1, x_2, \dots)
\]
such that the sum of the squares of their absolute values converges:
\[
\sum_{i=1}^\infty \lvert x_i \rvert^2 < \infty.
\]
\textbf{1. Define the Inner Product} The inner product is defined as
\[
\langle x, y \rangle = \sum_{i=1}^\infty x_i \overline{y_i}.
\]
By the Cauchy-Schwarz inequality for sequences, this series converges
absolutely, so the inner product is mathematically well-defined. It
trivially satisfies linearity, conjugate symmetry, and
positive-definiteness.

\textbf{2. Prove Completeness} Let \((x^{(n)})\) be a Cauchy sequence in
\(\ell^2\). For \(\epsilon > 0\), there is an \(N\) such that for
\(n, m > N\), we have
\[
\|x^{(n)} - x^{(m)}\|_2 < \epsilon.
\]
This implies that for each coordinate \(i\), the sequence
\[
(x_i^{(n)})_{n=1}^\infty
\]
is a Cauchy sequence in \(\mathbb{C}\). Since \(\mathbb{C}\) is
complete, each coordinate converges to a limit \(x_i \in \mathbb{C}\).
Let
\[
x = (x_1, x_2, \dots)
\]
. By taking the limit as \(m \to \infty\) in the finite partial sums of
the Cauchy condition, and extending to the infinite sum, we find
\(x^{(n)} - x \in \ell^2\). Since \(x^{(n)} \in \ell^2\) and \(\ell^2\)
is a vector space,
\[
x = x^{(n)} - (x^{(n)} - x) \in \ell^2.
\]
Finally, this shows \(\|x^{(n)} - x\|_2 \to 0\) as \(n \to \infty\),
meaning the sequence converges to \(x\) in the \(\ell^2\) norm. Because
every Cauchy sequence converges to an element within \(\ell^2\), the
space is complete, making it a Hilbert space.

\subsection{Exam Patterns}\label{exam-patterns-40}

\begin{itemize}
\tightlist
\item
  ``Define a Hilbert space and give an example'' is a highly probable
  short question. Provide the definition and the \(\ell^2\) sequence
  space example.
\item
  The 4-mark theoretical proof ``Every Hilbert is Banach'' is very easy
  if you clearly state definitions.
\end{itemize}

\end{document}
